pdfoutput=1
\pdfoutput=1
\documentclass[12pt,a4paper]{article}

% Packages
\usepackage{amsmath,amssymb,amsthm}
\usepackage{booktabs}
\usepackage{hyperref}
\usepackage[margin=2.5cm]{geometry}
\usepackage{setspace}
\usepackage{enumitem}
\onehalfspacing

% Theorem environments
\newtheorem{theorem}{Theorem}[section]
\newtheorem{corollary}[theorem]{Corollary}
\newtheorem{proposition}[theorem]{Proposition}
\newtheorem{lemma}[theorem]{Lemma}
\theoremstyle{definition}
\newtheorem{definition}{Definition}[section]
\theoremstyle{remark}
\newtheorem{remark}{Remark}[section]

\title{The Yajie Protocol: Technology Lock-in, Audit Sovereignty,\\and the Non-Proliferation Logic of Data Quality Assessment}
\author{SCX}
\date{June 2026}

\begin{document}

\maketitle

\begin{abstract}
As artificial intelligence systems become pervasive in high-stakes decision-making, the quality of training and inference data has emerged as a critical—yet systematically under-governed—input. This paper introduces the Yajie Protocol, a multi-expert consistency-based data auditing algorithm, and analyzes its structural implications for the emerging market in data quality assessment. We demonstrate that Yajie generates a distinctive form of technology lock-in driven not by network effects in the conventional sense, nor by proprietary control of intellectual property, but by a time-accumulated advantage mechanism: each audit cycle enriches a shared evidentiary corpus, widening the quality gap between the incumbent protocol and any potential entrant. Drawing on game-theoretic modeling, we formalize this dynamic as a \textit{Non-Proliferation Equilibrium} (NPE), wherein rational actors forgo the development of competing audit standards despite the absence of formal barriers to entry. We establish a structural analogy between this equilibrium and the strategic logic underpinning the Nuclear Non-Proliferation Treaty (NPT), arguing that audit standards fragmentation produces negative externalities analogous to nuclear proliferation. The paper further examines the geopolitical dimensions of audit infrastructure through the lens of the SWIFT financial messaging system and the International Accounting Standards Board (IASB), with particular attention to US-China technological decoupling. We model the lock-in trajectory through a five-stage progression, then extend the analysis to a sixth stage—termed \textit{Universal Harmony} ()—wherein the Cumulative Evidentiary Corpus transitions from a proprietary asset into a global public infrastructure and audit standards cease to belong to any single nation. We derive policy recommendations for embedding audit neutrality within multilateral governance frameworks. Our findings carry implications for the design of international AI governance institutions, the regulation of algorithmic auditing markets, and the preservation of epistemic sovereignty in an era of data-driven geopolitical competition.

\bigskip
\noindent\textbf{Keywords:} data quality assessment, algorithmic audit, technology lock-in, non-proliferation, game theory, audit sovereignty, AI governance, multi-expert consistency, SCX uncertainty principle, mutual deterrence, adoption-audit equilibrium
\end{abstract}

\section{Introduction}

The governance of artificial intelligence has entered a phase of institutional crystallization. Where the preceding decade was characterized by proliferation of ethical principles and voluntary guidelines (Floridi et al., 2018; Jobin et al., 2019), the current moment is defined by the translation of normative commitments into operational mechanisms: conformity assessments, third-party audits, and certification regimes (Raji et al., 2020; M\"{o}kander et al., 2021). Central to this operational turn is the question of \textit{data quality}—the accuracy, completeness, consistency, and fitness-for-purpose of the datasets that underpin machine learning pipelines.

The dependency of AI system performance on data quality is by now well-established in both the technical literature (Sambasivan et al., 2021; Northcutt et al., 2021) and policy discourse (European Commission, 2021; NIST, 2023). Yet the institutional infrastructure for assessing and certifying data quality remains nascent. A small number of commercial auditing services have emerged, alongside standards development efforts within ISO/IEC JTC 1/SC 42 and the IEEE P2863 working group. What has not yet materialized—and what this paper argues may \textit{never} materialize in a competitive form—is a pluralistic market of interoperable data quality assessment protocols.

The central claim of this paper is that data quality auditing exhibits structural properties that drive the market toward a natural monopoly, and that this monopolistic tendency is, counterintuitively, \textit{welfare-enhancing} under a specific set of conditions. We develop this argument through a detailed analysis of the Yajie Protocol, a data auditing algorithm that operationalizes multi-expert consistency as a quality signal, and we formalize the resulting market dynamics through the concept of a \textit{Non-Proliferation Equilibrium} (NPE).

The Yajie Protocol, which we describe in Section~3, operates by deploying an ensemble of independent audit modules against a target dataset and measuring the degree of inter-auditor agreement. Disagreement signals potential quality defects; convergence across auditors with orthogonal detection strategies provides evidence of quality. What makes the protocol structurally distinctive is its \textit{cumulative evidentiary architecture}: every audit cycle enriches a shared repository of detected anomalies, auditor calibration parameters, and dataset-specific quality profiles, such that the protocol becomes monotonically more accurate over time. New entrants cannot replicate this accumulated capital; they must either license access to the incumbent's evidentiary corpus (entrenching its position) or enter the market with a structurally inferior product.

This dynamic is not adequately captured by standard models of technology lock-in. The QWERTY literature (David, 1985; Arthur, 1989) emphasizes network effects, switching costs, and installed-base advantages. Platform economics (Rochet \& Tirole, 2003; Parker \& Van Alstyne, 2005) highlights cross-side network externalities. Neither framework fully accounts for a product whose quality is an increasing function of \textit{cumulative usage time}, independent of the number of users. The Yajie Protocol's lock-in mechanism is temporal, not networked: it accrues advantage through the sheer volume of audits conducted, each of which makes the protocol better at conducting the next audit. This \textit{quality ratchet} creates an entry barrier that grows with time, asymptotically approaching an insurmountable moat.

We deploy a game-theoretic framework to demonstrate that, under plausible parameterizations, rational competitors will choose not to invest in the development of competing audit protocols, even when the incumbent has not engaged in predatory pricing or exclusionary practices. The resulting market structure—a single dominant protocol, universally adopted—resembles not a coercive monopoly of the Standard Oil variety, but rather the equilibrium that obtains in domains where fragmentation itself is the primary source of social cost. We further demonstrate that the protocol's architecture generates a distinctive \textit{Adoption-Audit Equilibrium}: organizations are simultaneously incentivized to adopt Yajie (because it improves data quality) and deterred from falsifying data (because adoption renders audit outputs unconditionally verifiable by third parties, per the SCX Uncertainty Principle). The architecture transforms data quality from an unverifiable assertion into a revealed attribute—and honesty into the dominant strategy.

This is where the nuclear non-proliferation analogy becomes analytically productive. The Nuclear Non-Proliferation Treaty (NPT) succeeded not because it prohibited nuclear weapons development (it did not, for the nuclear-weapon states) but because it established a club whose benefits—access to peaceful nuclear technology, participation in the global non-proliferation regime—outweighed the strategic value of an independent nuclear program for most states (Sagan, 1996). Critically, the NPT's stability does not depend on enforcement alone; it depends on a rational calculation by non-nuclear states that proliferation would leave them \textit{worse off} than the status quo, either because their nascent capability would be neutralized by preemptive action, because the resultant arms race would consume resources better spent elsewhere, or because proliferation by one state would trigger proliferation by its rivals, leaving everyone less secure (Waltz, 1981; Mearsheimer, 1993).

The audit protocol market exhibits an analogous structure. A world with $N$ audit protocols, each partially trusted and partially contested, is a world in which audit results are themselves subject to meta-audit disputes, eroding the very certainty that audits are designed to provide. The development of a competing protocol by one jurisdiction would trigger reciprocal development by rivals, fragmenting the global audit infrastructure along geopolitical lines and imposing coordination costs on every entity that operates across jurisdictions. In this environment, the \textit{rational} choice for most actors is to accept the dominance of a single protocol—provided that protocol is governed under adequate neutrality safeguards—rather than to pursue audit autarky.

The paper proceeds as follows. Section~2 provides a technical overview of the Yajie Protocol and situates it within the literature on algorithmic auditing and technology lock-in. Section~3 develops the Non-Proliferation Equilibrium model, including the formal payoff structure, Nash equilibrium characterization, NPT analogy, and—critically—a corollary establishing that first-mover advantage decays over time because the Cumulative Evidentiary Corpus is an additive public good. Section~4 examines the geopolitical dimensions of audit infrastructure through case studies of SWIFT and the IASB, with an extension to the US-China technology competition. Section~5 models the lock-in trajectory through five stages, extends the Markov formalization to a sixth stage termed \textit{Universal Harmony} () wherein the CEC becomes global public infrastructure and audit standards cease to belong to any nation, and analyzes the business model dynamics. Section~6 derives policy recommendations for embedding audit neutrality within multilateral governance. Section~7 addresses limitations and identifies directions for future research. Section~8 concludes. Section~9 synthesizes the foregoing analyses into a unified game-theoretic architecture.

\section{Theoretical Background}

\subsection{The Yajie Protocol: A Technical Overview}

The Yajie Protocol is a computational framework for assessing the quality of structured and semi-structured datasets through multi-expert consistency analysis. At its core, Yajie operationalizes a simple but powerful principle: when multiple independent audit procedures, each designed to detect distinct categories of data defect, converge on a quality assessment, the probability that the assessment is correct increases multiplicatively. Conversely, when auditors diverge, the pattern of divergence carries diagnostic information about the nature and location of quality failures.

Formally, let $\mathcal{D}$ be a target dataset consisting of $n$ records, each characterized by a feature vector $\mathbf{x}_i \in \mathbb{R}^d$ and, where applicable, a label $y_i$. Let $\mathcal{A} = \{A_1, A_2, \ldots, A_k\}$ be an ensemble of $k$ audit modules, where each $A_j: \mathcal{D} \rightarrow \mathcal{R}_j$ maps the dataset to a structured audit report $\mathcal{R}_j = (\mathbf{q}_j, \mathbf{a}_j, \mathbf{c}_j)$ comprising:
\begin{itemize}
\item A quality score vector $\mathbf{q}_j \in [0,1]^m$ across $m$ quality dimensions (e.g., completeness, accuracy, consistency, timeliness, uniqueness);
\item An anomaly register $\mathbf{a}_j = \{(i, \tau, s)\}$ recording the index $i$ of each flagged record, the anomaly type $\tau$, and a severity score $s \in [0,1]$;
\item A confidence vector $\mathbf{c}_j \in [0,1]^m$ expressing the auditor's self-assessed reliability on each dimension.
\end{itemize}

The protocol's aggregation function $\Phi: \prod_{j=1}^k \mathcal{R}_j \rightarrow \mathcal{R}^*$ synthesizes the ensemble outputs into a unified audit report $\mathcal{R}^*$ through a weighted consensus mechanism:
\begin{equation}
\Phi(\mathcal{R}_1, \ldots, \mathcal{R}_k) = \left( \sum_{j=1}^k w_j \mathbf{q}_j, \; \bigcup_{j=1}^k \mathbf{a}_j \setminus \mathcal{F}, \; \prod_{j=1}^k \text{diag}(\mathbf{c}_j) \right)
\end{equation}
where $w_j$ are dynamic weights derived from each auditor's historical calibration performance, $\mathcal{F}$ is a false-positive filter trained on previously adjudicated anomalies, and the confidence aggregation employs element-wise multiplication to reflect the multiplicative reduction in uncertainty when independent auditors concur.

The protocol's defining architectural feature is the \textit{Cumulative Evidentiary Corpus} (CEC), denoted $\mathcal{E}_t$, which aggregates all audit outputs up to time $t$:
\begin{equation}
\mathcal{E}_t = \bigcup_{\tau=0}^{t} \{\mathcal{R}^*_\tau, \mathcal{J}_\tau, \Delta\mathcal{P}_\tau\}
\end{equation}
where $\mathcal{J}_\tau$ is the set of human-adjudicated judgments from audit cycle $\tau$, and $\Delta\mathcal{P}_\tau$ represents parameter updates derived from those judgments. The CEC serves three functions: (i) it provides training data for the false-positive filter $\mathcal{F}$; (ii) it enables dynamic recalibration of auditor weights $w_j$; and (iii) it accumulates a growing library of dataset-specific quality fingerprints that accelerate future audits of similar or related datasets.

The CEC grows with each audit cycle. Let $g(t)$ denote the cumulative number of audits conducted by time $t$. Each audit contributes, on average, $\bar{\alpha}$ new anomaly instances and $\bar{\beta}$ new calibration data points to the CEC. The protocol's accuracy $\theta(t)$—measured as the F1 score of its anomaly detection relative to ground truth—improves as a function of CEC size:
\begin{equation}
\theta(t) = \theta_{\max} - (\theta_{\max} - \theta_0) e^{-\gamma \cdot |\mathcal{E}_t|}
\end{equation}
where $\theta_{\max}$ is the asymptotic accuracy ceiling, $\theta_0$ is the initial accuracy, and $\gamma$ is a learning rate parameter. This functional form captures diminishing returns to additional data, consistent with the empirical scaling literature (Kaplan et al., 2020), while ensuring that the incumbent's accuracy advantage over a zero-CEC entrant grows monotonically with cumulative audit volume.

\subsection{Technology Lock-in: Classical Mechanisms and Their Limits}

The phenomenon of technology lock-in—whereby an inferior technology persists because the costs of coordinated transition exceed the benefits of switching to a superior alternative—has been a central concern of innovation economics since David's (1985) iconic analysis of the QWERTY keyboard. Arthur (1989) formalized the mechanisms through which increasing returns to adoption generate path-dependent outcomes: large fixed costs, learning effects, coordination effects, and adaptive expectations.

Subsequent literature refined and extended these mechanisms. Katz and Shapiro (1985, 1994) developed the economics of network effects, distinguishing between direct network effects (the value of a network increases with the number of users, as in telecommunications) and indirect network effects (a larger installed base attracts complementary goods, as in software platforms). Farrell and Saloner (1985) modeled the dynamics of standardization, demonstrating that excess inertia—the failure to adopt a superior standard because of coordination failure—can arise even when all actors prefer the superior standard.

The two-sided market literature (Rochet \& Tirole, 2003; Armstrong, 2006) extended these insights to platforms that intermediate between distinct user groups, showing that cross-side network externalities can generate powerful lock-in effects even in the absence of traditional switching costs. Once a platform achieves critical mass on both sides of the market, entrants face a chicken-and-egg problem: they cannot attract users from one side without first assembling a critical mass on the other.

These classical mechanisms do not, however, capture the lock-in dynamic generated by the Yajie Protocol. Consider the disanalogies:

\begin{enumerate}[label=\arabic*.]
\item \textbf{Network effects}: Yajie's value does not depend on the number of other users of the protocol. A single organization conducting repeated audits accrues the same CEC-driven accuracy gains regardless of whether other organizations also use the protocol. The lock-in is \textit{temporal}, not \textit{networked}.

\item \textbf{Switching costs}: While Yajie does generate switching costs in the conventional sense (retraining staff, reconfiguring pipelines, renegotiating contracts), the principal entry barrier is not the cost of switching \textit{away} from Yajie but the cost of replicating the \textit{accumulated audit history} that makes Yajie accurate. A competitor cannot buy, license, or reverse-engineer the CEC, because the CEC is not merely a database but a cumulative record of audit outcomes, adjudications, and parameter updates that are causally entangled with the protocol's specific architecture.

\item \textbf{Installed base}: Unlike a platform whose value derives from its user base, Yajie's value derives from its \textit{audit base}—the corpus of audits it has conducted. This is a stock variable, not a flow variable, and it cannot be replicated by attracting users away from the incumbent.

\item \textbf{Complementary goods}: The Yajie ecosystem does generate complementary assets—training programs, integration tooling, certification bodies—but these are derivative of the CEC advantage, not independent sources of lock-in.
\end{enumerate}

The closest analogue in the existing literature may be the concept of \textit{data network effects} (Gregory et al., 2021), wherein a product improves as it collects more data from users. However, the data network effect literature typically emphasizes the role of user-contributed data in improving a single firm's product (as in search engines, recommendation systems, or autonomous vehicles). The Yajie dynamic is distinct in three respects: (i) the improvement mechanism is audit-specific, not user-behavior-specific; (ii) the data in question are not user-contributed but \textit{audit-generated}, arising from the protocol's own operation; and (iii) the CEC's evidentiary character means that its value is not merely statistical (more data yields better predictions) but \textit{epistemic} (more audits yield a more trustworthy corpus of quality evidence).

\subsection{Algorithmic Auditing and Data Quality: State of the Field}

The algorithmic auditing literature has expanded rapidly since the mid-2010s, driven by concerns about fairness, accountability, and transparency in automated decision systems (Sandvig et al., 2014; Diakopoulos, 2015; Ada Lovelace Institute, 2020). Much of this literature has focused on \textit{model auditing}—the ex post examination of trained models for bias, discrimination, or other harms (Raji et al., 2020; Wilson et al., 2021)—rather than on \textit{data auditing}, the ex ante assessment of training data quality.

The relative neglect of data auditing is partly explained by technical challenges: data quality is inherently multidimensional (Batini et al., 2009), context-dependent (Wang \& Strong, 1996), and often unverifiable without access to the data-generation process. Recent work has begun to address these challenges. Northcutt et al. (2021) proposed Confident Learning for identifying label errors in training data. Sambasivan et al. (2021) documented the prevalence of data cascades—compounding errors originating in data quality failures—in real-world ML deployments. The ``Data Cards'' framework (Pushkarna et al., 2022) and ``Datasheets for Datasets'' (Gebru et al., 2021) have advanced structured transparency for dataset documentation.

What remains absent is a \textit{standardized, cumulative, and cross-domain} methodology for data quality assessment. Existing approaches are typically bespoke to specific dataset types, model architectures, or application domains. The Yajie Protocol addresses this gap through its multi-expert ensemble architecture, which accommodates heterogeneous audit modules within a unified aggregation framework. However, as we argue in the following sections, this very generality—combined with the CEC's cumulative character—generates the lock-in dynamics that are the central concern of this paper.

\section{The Non-Proliferation Equilibrium}

\subsection{Model Setup}

We consider a world with $N$ jurisdictions (which may be nations, regulatory blocs, or large organizations capable of developing audit infrastructure). Each jurisdiction $i \in \{1, \ldots, N\}$ faces a binary choice: \textit{Develop} ($D$) an independent data audit protocol, or \textit{Adopt} ($A$) the incumbent protocol (Yajie). The choice is made simultaneously by all jurisdictions; we examine the Nash equilibria of this one-shot game.

Let the incumbent protocol be indexed as $j=1$, with an accumulated CEC of size $|\mathcal{E}|$. An entrant protocol developed by jurisdiction $i$ would begin with an empty CEC, implying an accuracy penalty $\delta(|\mathcal{E}|) = \theta(|\mathcal{E}|) - \theta(0)$, where $\theta(\cdot)$ follows the functional form specified in Section~2.1. For any non-trivial $|\mathcal{E}|$, $\delta(|\mathcal{E}|) > 0$, and $\delta'(|\mathcal{E}|) > 0$ (the accuracy gap widens with cumulative audits).

The payoff to jurisdiction $i$ from adopting the incumbent is:
\begin{equation}
\pi_i(A, \mathbf{a}_{-i}) = V[\theta(|\mathcal{E}|)] - c_i^{\text{adopt}} - \lambda \cdot \mathbb{I}[\text{fragmentation}(\mathbf{a})]
\end{equation}
where $V[\cdot]$ is the value derived from audit accuracy, $c_i^{\text{adopt}}$ is the cost of adopting the incumbent protocol (licensing, integration, training), $\lambda$ is a fragmentation cost parameter, and $\mathbb{I}[\text{fragmentation}(\mathbf{a})]$ is an indicator that equals 1 if the global audit infrastructure is fragmented (i.e., if at least one jurisdiction has developed an independent protocol) and 0 otherwise.

The payoff from developing an independent protocol is:
\begin{equation}
\pi_i(D, \mathbf{a}_{-i}) = V[\theta(0)] - c_i^{\text{develop}} - \kappa \cdot \mathbf{1}_{-i}^D - \lambda \cdot \mathbb{I}[\text{fragmentation}(\mathbf{a})]
\end{equation}
where $c_i^{\text{develop}} \gg c_i^{\text{adopt}}$ is the fixed cost of protocol development, and $\kappa \cdot \mathbf{1}_{-i}^D$ captures the strategic cost imposed by other jurisdictions' development decisions: each additional developer adds $\kappa$ to $i$'s costs, reflecting the erosion of mutual recognition, the proliferation of conflicting audit standards, and the increased complexity of cross-jurisdictional data flows.

The fragmentation indicator is:
\begin{equation}
\mathbb{I}[\text{fragmentation}(\mathbf{a})] = \begin{cases} 0 & \text{if } \sum_{j=1}^N \mathbb{I}[a_j = D] = 0 \\ 1 & \text{if } \sum_{j=1}^N \mathbb{I}[a_j = D] \geq 1 \end{cases}
\end{equation}

Critically, the first developer triggers fragmentation for \textit{all} jurisdictions, including itself. This reflects the public-bads character of audit protocol proliferation: once multiple protocols exist, every cross-border data transaction must navigate conflicting quality certifications, and the informational value of any single audit is diminished by the existence of competing—and potentially contradictory—assessments.

\subsection{Rigorous Game Formalization and Assumptions}

Preliminary clarification before proceeding to equilibrium analysis: we formally cast the above scenario as a standard strategic-form game and explicitly list all assumptions.

\begin{definition}[Non-Proliferation Game $\Gamma^{\text{NP}}$]
The non-proliferation game $\Gamma^{\text{NP}} = \langle N, (S_i)_{i \in N}, (u_i)_{i \in N} 
angle$ is defined as:
\begin{itemize}
\item Player set $N = \{1, 2, \ldots, n\}$, $n \geq 2$, each jurisdiction is a player;
\item For each $i \in N$, pure strategy space $S_i = \{A, D\}$, where $A$ denotes Adopt the existing protocol, $D$ denotes Develop an independent protocol;
\item Payoff function family $u_i: S_1 \times \cdots \times S_n \to \mathbb{R}$, defined as in Equations (4)--(5).
\end{itemize}
Denote a strategy profile as $s = (s_i, s_{-i}) \in S \triangleq \prod_{j=1}^n S_j$, where $s_{-i}$ denotes the strategies of all jurisdictions except $i$.
\end{definition}

\begin{definition}[All Assumptions of the Non-Proliferation Game]
\label{def:assumptions-npe}
This section's theorems and propositions are all founded on the following assumptions. Each assumption is annotated with its economic meaning and relaxability.

\begin{enumerate}[label=\textbf{A\arabic*}, leftmargin=*, itemsep=4pt]
\item \textbf{Monotone concavity of the value function:} $V: [0,1] \to \mathbb{R}^+$ is strictly increasing, twice continuously differentiable, and concave, satisfying $V(0) = 0$, $V'(\theta) > 0$, $V''(\theta) \leq 0$ for all $\theta \in (0,1)$.\hfill\textit{Economic meaning: the marginal value of audit accuracy is non-increasing.}

\item \textbf{CEC functional form for audit accuracy:} Audit accuracy evolves with CEC size as:
\begin{equation}
\theta(|\mathcal{E}|) = \theta_{\max} - (\theta_{\max} - \theta_0) e^{-\gamma |\mathcal{E}|}
\end{equation}
where $\theta_{\max} \in (0,1]$ is the asymptotic accuracy upper bound, $\theta_0 \in (0, \theta_{\max})$ is the initial accuracy (empty CEC), and $\gamma > 0$ is the learning rate parameter.\hfill\textit{Economic meaning: accuracy grows monotonically with audit volume and exhibits diminishing marginal returns. Note: the specific exponential functional form can be relaxed to any function satisfying $\theta' > 0, \theta'' < 0$ without altering the qualitative conclusions.}

\item \textbf{Cost structure:} $c^{\text{develop}} > c^{\text{adopt}} > 0$. The total fixed cost of developing an independent protocol strictly exceeds the licensing and integration costs of adopting the existing protocol.\hfill\textit{Economic meaning: development investment is an irreversible sunk cost. The homogeneity assumption can be relaxed to jurisdiction-heterogeneous costs without affecting equilibrium existence, though it impacts the precise equilibrium configuration.}

\item \textbf{Parameter ordering of proliferation externalities:} $\kappa > 0$ is the marginal proliferation cost imposed on jurisdiction $i$ by each other jurisdiction's development action; $\lambda > 0$ is the fixed fragmentation cost imposed on \textit{all} jurisdictions (including the first developer itself) when the fragmentation indicator is activated. The ordering $\lambda > \kappa$ is satisfied.\hfill\textit{Economic meaning: the public-bads effect of fragmentation (collapse of mutual recognition, conflicting standards, rising cross-border transaction costs) exceeds the incremental cost imposed on other jurisdictions by a single proliferation action. This ordering is the critical condition for the stability of the universal-adoption equilibrium.}

\item \textbf{Complete information:} All parameters of the game——$\{V(\cdot), \theta(\cdot), \theta_0, \theta_{\max}, \gamma, c^{\text{adopt}}, c^{\text{develop}}, \kappa, \lambda, |\mathcal{E}|\}$——constitute common knowledge among all players.\hfill\textit{Relaxability: This is the strongest assumption of the paper. Section~7 discusses the extension to incomplete information (jurisdictions possess private information about their own development costs). Introducing incomplete information would generate refinement problems for Bayesian-Nash equilibria.}

\item \textbf{Simultaneous moves:} All jurisdictions make their adopt/develop decisions simultaneously and independently.\hfill\textit{Relaxability: Can be extended to a sequential-move game, where early-moving jurisdictions may influence late-moving jurisdictions' decisions by developing first. Strategic preemptive equilibria may arise in sequential games, and their welfare properties may differ from those of simultaneous-move equilibria.}

\item \textbf{Additive separability of payoffs and publicness of fragmentation:} Payoff functions possess the additive structure specified in equations (4)--(5). The fragmentation indicator $\mathbb{I}[n_D > 0]$ is public information——once any jurisdiction develops, the fragmentation cost is immediately imposed on all jurisdictions.\hfill\textit{Meaning: fragmentation is a ``public bad'' whose supply is triggered by the first developer's action and whose consumption is non-excludable and non-rival.}
\end{enumerate}
\end{definition}

\begin{remark}[Logical Dependencies Among Assumptions]
Assumptions A2 and A1 jointly guarantee that the adoption advantage $\Delta(|\mathcal{E}|)$ is strictly monotonically increasing with CEC size. The ordering $\lambda > \kappa$ in Assumption A4 guarantees that the stability margin of the universal-adoption equilibrium is positive and increasing in $|\mathcal{E}|$. If $\lambda \leq \kappa$, the fragmentation cost is insufficient to deter unilateral proliferation, and the stability of the non-proliferation equilibrium is substantially weakened—in that case, a small $n$ or additional institutional constraints (such as treaty obligations) would be required to maintain the equilibrium.
\end{remark}

\subsection{Rigorous Existence, Uniqueness, and Characterization of Nash Equilibria}

This subsection presents the complete equilibrium characterization of the non-proliferation game in theorem form. All proofs are rigorous under the stated assumptions.

Let $n_D(s) = |\{i \in N : s_i = D\}|$ denote the number of developers in strategy profile $s$. Define jurisdiction $i$'s \textbf{unilateral adoption advantage}—the net benefit of adopting relative to developing when all other jurisdictions adopt:
\begin{equation}
\Delta(|\mathcal{E}|) \triangleq u_i(A, \mathbf{A}_{-i}) - u_i(D, \mathbf{A}_{-i}) = V[\theta(|\mathcal{E}|)] - V[\theta(0)] - (c^{\text{adopt}} - c^{\text{develop}}) + \kappa
\label{eq:adoption-advantage-def}
\end{equation}
where $\mathbf{A}_{-i}$ denotes that all jurisdictions other than $i$ choose $A$. Note that $\Delta(|\mathcal{E}|)$ depends only on $|\mathcal{E}|$ and is independent of jurisdiction identity $i$ (symmetry of the strategy space).

\begin{theorem}[Complete Nash Equilibrium Characterization of the Non-Proliferation Game (Yajie Non-Proliferation Theorem)]
\label{thm:npe-complete}
Under Assumptions A1--A7, consider the non-proliferation game $\Gamma^{\text{NP}}$ with $N = \{1, \ldots, n\}$ and $n \geq 2$. The following propositions hold.

\begin{enumerate}[label=(\roman*), itemsep=6pt]
\item \textbf{Equilibrium existence:} A pure-strategy Nash equilibrium always exists. If $\Delta(|\mathcal{E}|) \geq \lambda - \kappa$, then $(A, \ldots, A)$ is a pure-strategy Nash equilibrium; if $\Delta(|\mathcal{E}|) \leq -(n-1)\kappa$, then $(D, \ldots, D)$ is a pure-strategy Nash equilibrium; in the intermediate region, a symmetric mixed-strategy Nash equilibrium exists.\label{item:existence}

\item \textbf{Universal-adoption equilibrium:} The strategy profile $s^* = (A, A, \ldots, A)$ is a Nash equilibrium if and only if:
\begin{equation}
\boxed{\Delta(|\mathcal{E}|) \geq \lambda - \kappa}
\label{eq:all-adopt-condition}
\end{equation}
This equilibrium is \textbf{strict} when the strict inequality holds. A strict Nash equilibrium rules out any player's reliance on indifference between strategies.\label{item:all-adopt}

\item \textbf{Universal-development equilibrium:} The strategy profile $s^{D} = (D, D, \ldots, D)$ is a Nash equilibrium if and only if:
\begin{equation}
\boxed{\Delta(|\mathcal{E}|) \leq -(n-1)\kappa}
\label{eq:all-develop-condition}
\end{equation}
\label{item:all-develop}

\item \textbf{Impossibility of asymmetric equilibria:} In any pure-strategy Nash equilibrium, all jurisdictions must choose the same strategy. There exists no asymmetric pure-strategy equilibrium with ``partial adoption, partial development.''\label{item:no-asym}

\item \textbf{Symmetric mixed-strategy equilibrium:} When neither condition (\ref{eq:all-adopt-condition}) nor (\ref{eq:all-develop-condition}) is satisfied—i.e., $-(n-1)\kappa < \Delta(|\mathcal{E}|) < \lambda - \kappa$—there exists a unique symmetric mixed-strategy Nash equilibrium. In this equilibrium, each jurisdiction independently chooses $A$ with the same probability $p^* \in (0,1)$, where $p^*$ is uniquely determined by the following indifference condition:
\begin{align}
&V[\theta(|\mathcal{E}|)] - c^{\text{adopt}} - \lambda\left[1 - (p^*)^{n-1}\right] \nonumber\\
&\quad = V[\theta(0)] - c^{\text{develop}} - n\kappa - \lambda
\label{eq:indifference}
\end{align}
\end{enumerate}
\end{theorem}

\begin{proof}
We prove each item in turn under Assumptions A1--A7.

\textbf{(i) Existence.} Under A1--A7, the game $\Gamma^{\text{NP}}$ is a finite game ($|S| = 2^n < \infty$). By Nash (1951), every finite game possesses a mixed-strategy Nash equilibrium. The existence of pure-strategy equilibria follows from the constructive argument below: the set of all $2^n$ pure-strategy profiles is finite; each jurisdiction's payoff, given others' strategies, is an affine function of its own strategy (equations (4)--(5)), hence the best-response correspondence is upper-hemicontinuous and convex-valued; or more directly, as shown below, equilibrium conditions (\ref{eq:all-adopt-condition}) and (\ref{eq:all-develop-condition}) cover the full parameter space, at least one of them holds when $\Delta(|\mathcal{E}|)$ is sufficiently large or sufficiently small, and when neither holds, the symmetric mixed-strategy equilibrium is guaranteed to exist by a fixed-point argument.

\textbf{(ii) Universal-adoption equilibrium.} Consider the profile $s^* = (A, \ldots, A)$. Under $s^*$, $n_D(s^*) = 0$, and the fragmentation indicator $\mathbb{I}[n_D > 0] = 0$. For any jurisdiction $i \in N$:
\begin{align}
u_i(A, \mathbf{A}_{-i}) &= V[\theta(|\mathcal{E}|)] - c^{\text{adopt}} \label{eq:uA}\\
u_i(D, \mathbf{A}_{-i}) &= V[\theta(0)] - c^{\text{develop}} - \kappa - \lambda \label{eq:uD}
\end{align}
Derivation of (\ref{eq:uD}): when $i$ unilaterally deviates to $D$, $n_D = 1 > 0$, fragmentation is triggered ($-\lambda$), and $i$'s own development imposes $\kappa$ on $i$ itself (since $i$ is the sole developer, the marginal proliferation cost from its own development is captured by the $\kappa \cdot (n_D(s_{-i}) + 1)$ term in equation (5), with $n_D(s_{-i}) = 0$).

$s^*$ is a Nash equilibrium iff for all $i$, $u_i(A, \mathbf{A}_{-i}) \geq u_i(D, \mathbf{A}_{-i})$. Substituting (\ref{eq:uA}) and (\ref{eq:uD}):
\begin{align}
V[\theta(|\mathcal{E}|)] - c^{\text{adopt}} &\geq V[\theta(0)] - c^{\text{develop}} - \kappa - \lambda
\end{align}
Using the definition (\ref{eq:adoption-advantage-def}), this simplifies to condition (\ref{eq:all-adopt-condition}). When the strict inequality holds, the deviation payoff is strictly less than the equilibrium payoff, and the equilibrium is strict—this guarantees robustness under small payoff perturbations (à la Kohlberg \& Mertens, 1986).

\textbf{(iii) Universal-development equilibrium.} Consider $s^{D} = (D, \ldots, D)$. Under $s^{D}$, $n_D = n$, and the fragmentation indicator is identically 1. For any $i$:
\begin{align}
u_i(D, \mathbf{D}_{-i}) &= V[\theta(0)] - c^{\text{develop}} - \kappa \cdot n - \lambda \label{eq:uDD}\\
u_i(A, \mathbf{D}_{-i}) &= V[\theta(|\mathcal{E}|)] - c^{\text{adopt}} - \kappa \cdot (n-1) - \lambda \label{eq:uAD}
\end{align}
In (\ref{eq:uDD}), $n_D(s_{-i}) = n-1$ (all except $i$ develop), adding $i$'s own development yields a total proliferation count of $n$, so the proliferation cost term is $\kappa \cdot n$. In (\ref{eq:uAD}), $i$ unilaterally deviates to $A$, with $n_D(s_{-i}) = n-1$ and fragmentation still 1.

$s^{D}$ is a Nash equilibrium iff for all $i$, $u_i(D, \mathbf{D}_{-i}) \geq u_i(A, \mathbf{D}_{-i})$:
\begin{align}
V[\theta(0)] - c^{\text{develop}} - n\kappa - \lambda &\geq V[\theta(|\mathcal{E}|)] - c^{\text{adopt}} - (n-1)\kappa - \lambda
\end{align}
Canceling $-\lambda$ and rearranging:
\begin{align}
V[\theta(|\mathcal{E}|)] - V[\theta(0)] + c^{\text{adopt}} - c^{\text{develop}} + \kappa &\leq -(n-1)\kappa
\end{align}
The left-hand side is precisely $\Delta(|\mathcal{E}|)$, yielding condition (\ref{eq:all-develop-condition}).

\textbf{(iv) Impossibility of asymmetric equilibria.} Suppose there exists a pure-strategy equilibrium $s^*$ with $\exists i, j: s_i^* = A, s_j^* = D$. Let $n_D^* = |\{k: s_k^* = D\}| \in \{1, \ldots, n-1\}$. Consider an adopting jurisdiction $i$ and a developing jurisdiction $j$. By the necessary conditions for Nash equilibrium:
\begin{align}
u_i(A, s_{-i}^*) &\geq u_i(D, s_{-i}^*) \label{eq:condA}\\
u_j(D, s_{-j}^*) &\geq u_j(A, s_{-j}^*) \label{eq:condD}
\end{align}

Compute the payoff difference for $i$. In $s^*$, $n_D(s_{-i}^*) = n_D^*$ (since $i$ chooses $A$ and does not affect the $n_D$ count), and the fragmentation indicator is 1 (since $n_D^* > 0$):
\begin{align}
u_i(A, s_{-i}^*) - u_i(D, s_{-i}^*) &= \Delta(|\mathcal{E}|) - \kappa \cdot n_D^* - \lambda - (-\kappa \cdot (n_D^* + 1) - \lambda) \nonumber\\
&= \Delta(|\mathcal{E}|) + \kappa
\end{align}

Similarly, for $j$ ($s_j^* = D$), $n_D(s_{-j}^*) = n_D^* - 1$:
\begin{align}
u_j(D, s_{-j}^*) - u_j(A, s_{-j}^*) &= -\Delta(|\mathcal{E}|) - \kappa
\end{align}

Condition (\ref{eq:condA}) requires $\Delta(|\mathcal{E}|) + \kappa \geq 0$, and condition (\ref{eq:condD}) requires $-\Delta(|\mathcal{E}|) - \kappa \geq 0$. Both conditions hold simultaneously only if $\Delta(|\mathcal{E}|) = -\kappa$, a parameter configuration of Lebesgue measure zero. On this null set, asymmetric equilibria may exist (all jurisdictions indifferent between $A$ and $D$), but they are not strict and collapse under small payoff perturbations. Hence, for almost all parameter values, asymmetric pure-strategy equilibria do not exist.

\textbf{(v) Mixed-strategy equilibrium.} In the region $-(n-1)\kappa < \Delta(|\mathcal{E}|) < \lambda - \kappa$, neither universal adoption nor universal development is an equilibrium. Let each jurisdiction independently choose $A$ with probability $p \in (0,1)$. Then the number of jurisdictions other than $i$ choosing $D$ follows a binomial distribution: $n_D(s_{-i}) \sim \text{Binomial}(n-1, 1-p)$.

Jurisdiction $i$'s expected payoff from choosing $A$:
\begin{align}
\mathbb{E}[u_i(A, s_{-i})] &= V[\theta(|\mathcal{E}|)] - c^{\text{adopt}} - \kappa \cdot \mathbb{E}[n_D] - \lambda \cdot \mathbb{P}[n_D > 0] \nonumber\\
&= V[\theta(|\mathcal{E}|)] - c^{\text{adopt}} - \kappa (n-1)(1-p) - \lambda\left[1 - p^{n-1}\right]
\end{align}

Expected payoff from choosing $D$:
\begin{align}
\mathbb{E}[u_i(D, s_{-i})] &= V[\theta(0)] - c^{\text{develop}} - \kappa(\mathbb{E}[n_D] + 1) - \lambda \nonumber\\
&= V[\theta(0)] - c^{\text{develop}} - \kappa[(n-1)(1-p) + 1] - \lambda
\end{align}

In a symmetric mixed-strategy equilibrium, the two expected payoffs are equal (indifference condition). Subtracting and canceling the $\kappa(n-1)(1-p)$ terms:
\begin{align}
V[\theta(|\mathcal{E}|)] - c^{\text{adopt}} - \lambda\left[1 - p^{n-1}\right] = V[\theta(0)] - c^{\text{develop}} - \kappa - \lambda
\end{align}
Rearranging yields equation (\ref{eq:indifference}). Set:
\begin{equation}
\Psi(p) = \lambda \cdot p^{n-1} \quad \text{and} \quad \Gamma = -\Delta(|\mathcal{E}|) + \lambda - 2\kappa
\end{equation}
The indifference condition is equivalent to $\Psi(p^*) = \Gamma$. Since $\Psi: [0,1] \to [0,\lambda]$ is strictly increasing and continuous ($n \geq 2, \lambda > 0$), and in this parameter region $\Gamma \in (0, \lambda)$, the intermediate value theorem guarantees a unique $p^* = \Psi^{-1}(\Gamma) \in (0,1)$. By Glicksberg (1952) (existence of mixed-strategy equilibria with continuous payoffs) or directly by Nash's construction, this $p^*$ corresponds to a symmetric mixed-strategy Nash equilibrium.
\end{proof}

\begin{remark}[Large-Population Approximation]
When $n$ is large (e.g., over 190 sovereign states globally), $p^{n-1} \to 0$ for any $p < 1$. In this case the mixed-strategy probability approximates:
\begin{equation}
p^* \approx \frac{V[\theta(0)] - c^{\text{develop}} - \kappa - V[\theta(|\mathcal{E}|)] + c^{\text{adopt}}}{\lambda} + 1
\end{equation}
In the large-population limit $n \to \infty$, the influence of any single jurisdiction on others' strategies vanishes (competitive limit), and the symmetric mixed strategy degenerates into the individual optimization of an ``atomistic jurisdiction'' facing a given adoption environment.
\end{remark}

\begin{theorem}[Asymptotic Uniqueness of the Non-Proliferation Equilibrium and the CEC Critical Threshold]
\label{thm:asymptotic-uniqueness}
Under Assumptions A1--A7:

\begin{enumerate}[label=(\roman*)]
\item There exists a unique \textbf{CEC critical size} $|\mathcal{E}|^*$, defined as:
\begin{equation}
|\mathcal{E}|^* \triangleq \inf\left\{|\mathcal{E}| \geq 0 : \Delta(|\mathcal{E}|) \geq \lambda - \kappa\right\}
\label{eq:E-star}
\end{equation}
If $\Delta(0) \geq \lambda - \kappa$, then $|\mathcal{E}|^* = 0$ (the protocol locks in the market from its initial release). Otherwise, by the continuity and strict monotonicity of $\theta(\cdot)$ and $V(\cdot)$, $|\mathcal{E}|^*$ is a well-defined positive real number.

\item For all $|\mathcal{E}| > |\mathcal{E}|^*$, $s^* = (A, \ldots, A)$ is the \textbf{unique} Nash equilibrium of the game $\Gamma^{\text{NP}}$—no other pure-strategy equilibrium and no mixed-strategy equilibrium exists.

\item The explicit solution for the CEC critical threshold is:
\begin{equation}
|\mathcal{E}|^* = \frac{1}{\gamma} \ln \left( \frac{\theta_{\max} - \theta_0}{\theta_{\max} - \theta^*} \right)
\end{equation}
where $\theta^* = V^{-1}\left(V[\theta(0)] + c^{\text{adopt}} - c^{\text{develop}} - \kappa + \lambda - \kappa\right)$, and $V^{-1}$ is the inverse function of $V$ (whose existence is guaranteed by the strict monotonicity in A1).
\end{enumerate}
\end{theorem}

\begin{proof}
\textbf{(i)} The function $\Delta(|\mathcal{E}|)$ is continuous on $[0, \infty)$ (A1--A2 guarantee that $V \circ \theta$ is continuous) and strictly increasing:
\begin{equation}
\frac{d\Delta}{d|\mathcal{E}|} = V'[\theta(|\mathcal{E}|)] \cdot \theta'(|\mathcal{E}|) = V'[\theta(|\mathcal{E}|)] \cdot \gamma(\theta_{\max} - \theta_0) e^{-\gamma |\mathcal{E}|} > 0
\end{equation}
where $V' > 0$ (A1) and $\theta' > 0$ (A2). The infimum definition yields a unique critical value.

\textbf{(ii)} When $|\mathcal{E}| > |\mathcal{E}|^*$, by definition $\Delta(|\mathcal{E}|) > \lambda - \kappa$ (strict inequality). Theorem~\ref{thm:npe-complete}(ii) guarantees that $(A, \ldots, A)$ is a strict Nash equilibrium. The universal-development equilibrium condition $\Delta(|\mathcal{E}|) \leq -(n-1)\kappa$ does not hold (since $\Delta(|\mathcal{E}|) > \lambda - \kappa > 0 > -(n-1)\kappa$, by A4 $\lambda > \kappa > 0$). A mixed-strategy equilibrium requires the indifference condition to hold, but under a strict Nash equilibrium the expected payoff from deviation is strictly less than the equilibrium payoff, so the indifference condition cannot be satisfied. Hence $(A, \ldots, A)$ is the unique equilibrium.

\textbf{(iii)} Solving $\Delta(|\mathcal{E}|) = \lambda - \kappa$ for $\theta^*$ and substituting the functional form from A2 yields the explicit solution.
\end{proof}

\begin{corollary}[Monotonic Self-Reinforcing Property of the Non-Proliferation Equilibrium]
\label{cor:self-reinforcing}
Under Assumptions A1--A7, define the \textbf{stability margin} of the non-proliferation equilibrium:
\begin{equation}
M(|\mathcal{E}|) \triangleq \Delta(|\mathcal{E}|) - (\lambda - \kappa)
\end{equation}
Then:
\begin{enumerate}[label=(\roman*)]
\item $M(|\mathcal{E}|)$ is strictly increasing on $[0, \infty)$: $M'(|\mathcal{E}|) = V'[\theta(|\mathcal{E}|)] \cdot \theta'(|\mathcal{E}|) > 0$;
\item $\lim_{|\mathcal{E}| \to \infty} M(|\mathcal{E}|) = V(\theta_{\max}) - V(\theta_0) + c^{\text{develop}} - c^{\text{adopt}} + \kappa - \lambda$;
\item Each audit cycle expands the CEC by $\Delta |\mathcal{E}| > 0$, simultaneously: (a) increasing the attractiveness of adoption ($\partial u_i(A, \cdot) / \partial |\mathcal{E}| > 0$), (b) increasing the relative loss from deviation ($M(|\mathcal{E}|)$ grows), and (c) reducing the uncertainty of the adoption probability $p^*$ in the mixed-strategy equilibrium (when the equilibrium is pure-strategy, the uncertainty goes to zero).
\item This property—that the stability of the equilibrium is self-reinforcing as it is achieved—is the precise game-theoretic expression of the \textit{time-accumulated advantage} mechanism.
\end{enumerate}
\end{corollary}

\begin{proof}
Direct verification from Theorems~\ref{thm:npe-complete} and~\ref{thm:asymptotic-uniqueness} combined with the monotonicity properties of A1--A2.
\end{proof}

\begin{corollary}[Temporal Decay of First-Mover Advantage]
\label{cor:first-mover-decay}
Under Assumptions A1--A7. Define the following quantities:

\begin{definition}[Normalized First-Mover Advantage Intensity]
The \textbf{absolute accuracy advantage} of the first mover (with CEC size $|\mathcal{E}|$) relative to a zero-CEC late mover is:
\begin{equation}
\delta(|\mathcal{E}|) \triangleq \theta(|\mathcal{E}|) - \theta_0 = (\theta_{\max} - \theta_0)(1 - e^{-\gamma |\mathcal{E}|})
\end{equation}
The \textbf{normalized first-mover advantage intensity}—i.e., the ratio of the first-mover advantage to the maximum possible advantage—is:
\begin{equation}
\Xi(|\mathcal{E}|) \triangleq \frac{\delta(|\mathcal{E}|)}{\theta_{\max} - \theta_0} = 1 - e^{-\gamma |\mathcal{E}|} \in [0, 1)
\end{equation}
\end{definition}

The following propositions hold:

\begin{enumerate}[label=(\roman*), itemsep=6pt]
\item \textbf{Diminishing marginal returns and asymptotic saturation of the absolute advantage:} The first mover's absolute accuracy advantage $\delta(|\mathcal{E}|)$ is monotonically increasing ($\delta'(|\mathcal{E}|) = \gamma(\theta_{\max} - \theta_0)e^{-\gamma |\mathcal{E}|} > 0$), but its \textit{marginal increment} strictly decreases to zero:
\begin{equation}
\delta''(|\mathcal{E}|) = -\gamma^2 (\theta_{\max} - \theta_0) e^{-\gamma |\mathcal{E}|} < 0, \quad \lim_{|\mathcal{E}| \to \infty} \delta'(|\mathcal{E}|) = 0
\end{equation}
and the absolute advantage converges to a finite upper bound:
\begin{equation}
\lim_{|\mathcal{E}| \to \infty} \delta(|\mathcal{E}|) = \theta_{\max} - \theta_0 < \infty
\end{equation}
\textbf{Economic meaning: first-mover advantage does not grow without bound—the accuracy contribution of the CEC has a hard upper limit (determined by the theoretical detection capability of the audit modules). Below this ceiling, no matter how much audit history the first mover accumulates, its accuracy advantage cannot exceed $\theta_{\max} - \theta_0$.}

\item \textbf{Marginal decay of the adoption advantage:} Substituting $\delta(|\mathcal{E}|)$ into the definition of $\Delta(|\mathcal{E}|)$ (equation~\ref{eq:adoption-advantage-def}):
\begin{equation}
\Delta(|\mathcal{E}|) = V[\theta_0 + \delta(|\mathcal{E}|)] - V[\theta_0] - (c^{\text{adopt}} - c^{\text{develop}}) + \kappa
\end{equation}
By Assumption A1 ($V$ is strictly concave), $\Delta'(|\mathcal{E}|) = V'[\theta(|\mathcal{E}|)] \cdot \delta'(|\mathcal{E}|) > 0$ and $\Delta''(|\mathcal{E}|) < 0$. The marginal growth of the adoption advantage is simultaneously constrained by diminishing marginal returns to accuracy ($\delta' \to 0$) and diminishing marginal returns to value ($V' \to 0$ as $\theta \to \theta_{\max}$). Hence:
\begin{equation}
\lim_{|\mathcal{E}| \to \infty} \Delta'(|\mathcal{E}|) = 0
\end{equation}
\textbf{Economic meaning: the first mover's ability to expand its adoption advantage by continuing to accumulate CEC diminishes over time—the marginal contribution of additional audit cycles to equilibrium stability converges to zero.}

\item \textbf{CEC as an additive public good (the core mechanism of temporal decay):} The essence of the CEC is that of an \textit{additive public good}: (a) \textbf{Additivity}—each audit cycle appends new entries to the CEC ($\Delta |\mathcal{E}| > 0$), and the CEC size is non-decreasing (Assumption M4); (b) \textbf{Publicness}—the anomaly registers, calibration parameters, and quality fingerprints stored in the CEC are \textit{non-rivalrous} knowledge: jurisdiction $i$'s use of a particular calibration datum does not diminish jurisdiction $j$'s ability to use the same datum.

Suppose the CEC is internationalized at time $T^*$—i.e., transformed through a multilateral governance mechanism into a global public infrastructure equally accessible to all jurisdictions (see Section~6, Recommendation~1 on the IDAA). For any $t > T^*$, let $|\mathcal{E}_t^{\text{public}}|$ be the CEC size that any late-moving jurisdiction can immediately access. The cold-start penalty of a late mover shrinks from $\theta(|\mathcal{E}_t|) - \theta(0)$ to $\theta(|\mathcal{E}_t|) - \theta(|\mathcal{E}_t^{\text{public}}|)$. Define the \textbf{late-mover penalty ratio}:
\begin{equation}
\Psi(t; T^*) \triangleq \frac{\theta(|\mathcal{E}_t|) - \theta(|\mathcal{E}_t^{\text{public}}|)}{\theta(|\mathcal{E}_t|) - \theta(0)} \in [0, 1]
\end{equation}
$\Psi(T^*; T^*) = 1$ (at the moment of internationalization, the late mover has not yet accessed the public portion of the CEC, so the penalty is complete). But as $|\mathcal{E}_t^{\text{public}}| \to |\mathcal{E}_t|$ (the public CEC catches up to the first mover's private CEC), $\Psi \to 0$. Under ideal internationalization ($|\mathcal{E}_t^{\text{public}}| = |\mathcal{E}_t|$), $\Psi \equiv 0$—the first-mover advantage is completely eliminated at the adoption level.

However, complete elimination applies only to the accuracy dimension of a single audit. The first mover retains two \textit{non-internationalizable} residual advantages: (a) \textbf{the temporal depth of ecosystem inertia}—the first mover's adopter network, regulatory recognition, and institutional embedding possess incompressible historical depth (the stage lock-in of Theorem~\ref{thm:absorption}); (b) \textbf{the path-dependence of operational standard-setting authority}—the anomaly adjudication thresholds and calibration parameter configurations in the CEC reflect the first mover's initial data characteristics, and these features remain as historical anchors of the operational standard even after internationalization.

\item \textbf{Half-life of first-mover advantage:} Define the \textbf{adoption half-life} $T_{1/2}^{\Delta}$ as the time satisfying:
\begin{equation}
\Delta(|\mathcal{E}_{T_{1/2}^{\Delta}}|) - (\lambda - \kappa) = \frac{1}{2}\left[\Delta(|\mathcal{E}_{T^*}|) - (\lambda - \kappa)\right]
\end{equation}
i.e., the NPE stability margin (Corollary~\ref{cor:self-reinforcing}) drops to one-half of its value at the moment of internationalization. Under the functional form of A2, $T_{1/2}^{\Delta}$ satisfies the implicit equation:
\begin{equation}
V[\theta(|\mathcal{E}_{T_{1/2}^{\Delta}}|)] - V[\theta(|\mathcal{E}_{T^*}|)] = \frac{1}{2}\left[V[\theta(|\mathcal{E}_{T^*}|)] - V[\theta_0] - (c^{\text{adopt}} - c^{\text{develop}}) + \kappa - (\lambda - \kappa)\right]
\end{equation}
The existence of the half-life is guaranteed by the continuity of $V \circ \theta$ and the monotone boundedness of $\Delta(\cdot)$. Its finiteness implies: \textbf{first-mover advantage is not a permanent structure but a phased phenomenon with a clearly defined temporal boundary.}

\item \textbf{Honest strike—normative correction of the NPE:} The conclusion established by Theorem~\ref{thm:npe-complete} and Corollary~\ref{cor:self-reinforcing}—that ``the equilibrium self-reinforces as the CEC grows''—must be qualified as follows once the CEC is analyzed as an additive public good: the self-reinforcement of the NPE occurs only under the condition that the CEC is \textit{exclusively held}. Once the CEC is internationalized (through the IDAA or an equivalent multilateral mechanism), the component of the adoption advantage $\Delta(|\mathcal{E}|)$ that is specific to the first mover—namely, the $\theta(|\mathcal{E}_t^{\text{private}}|) - \theta(|\mathcal{E}_t^{\text{public}}|)$ term—tends to zero. At this point, the motivation for adoption shifts from ``developing an alternative protocol is not worth it'' to ``developing an alternative protocol is unnecessary''—all jurisdictions benefit equally from the global CEC infrastructure, the public-bads effect of fragmentation $\lambda$ continues to deter proliferation, but the basis of deterrence shifts from the first mover's private CEC monopoly to the universal availability of the CEC as a global public good.

In other words, the NPE undergoes an institutional transformation from a \textit{first-mover monopoly deterrence} regime to a \textit{public-good coordination equilibrium}. Before the transformation, the equilibrium is sustained by $\Delta(|\mathcal{E}^{\text{private}}|) \geq \lambda - \kappa$ (first-mover monopoly NPE); after the transformation, it is sustained by $\Delta(|\mathcal{E}^{\text{public}}|) \geq \lambda - \kappa$ (public-good NPE), where $\Delta(|\mathcal{E}^{\text{public}}|)$ is identical for all jurisdictions—the first mover's identity-based advantage disappears; only the CEC's scale advantage remains (which is a public good shared by all jurisdictions).

\textbf{This is the game-theoretic foundation of ``Universal Harmony''—see Section~5.2.2 for the full formalization.}
\end{enumerate}
\end{corollary}

\begin{proof}
\textbf{(i)--(ii)} Direct differentiation. By A2, $\theta'(|\mathcal{E}|) = \gamma(\theta_{\max} - \theta_0)e^{-\gamma |\mathcal{E}|} > 0$, $\theta''(|\mathcal{E}|) = -\gamma^2(\theta_{\max} - \theta_0)e^{-\gamma |\mathcal{E}|} < 0$. Hence $\delta'(|\mathcal{E}|) > 0$, $\delta''(|\mathcal{E}|) < 0$, and $\delta'(|\mathcal{E}|) \to 0$ as $|\mathcal{E}| \to \infty$. $\lim_{|\mathcal{E}| \to \infty} \delta(|\mathcal{E}|) = \theta_{\max} - \theta_0$ follows from the exponential decay term going to zero. The monotonicity and concavity of $\Delta(|\mathcal{E}|)$ follow from the composite of the monotone concavity of $V$ (A1) with the monotone concavity of $\delta(|\mathcal{E}|)$, which preserves strict concavity.

\textbf{(iii)} The public-good character of the CEC: non-rivalry is guaranteed by the nature of knowledge—calibration parameters and anomaly patterns are information, and information consumption is non-rivalrous. Additivity is guaranteed by the CEC definition (equation (2))—$\mathcal{E}_t = \bigcup_{\tau=0}^{t} \{\cdots\}$, and the union operation is monotonically non-decreasing. Limit behavior of the late-mover penalty ratio $\Psi(t; T^*)$: as $|\mathcal{E}_t^{\text{public}}| \to |\mathcal{E}_t|$, by continuity of $\theta$, $\theta(|\mathcal{E}_t^{\text{public}}|) \to \theta(|\mathcal{E}_t|)$, so $\Psi \to 0$. The non-internationalizable residual advantages are guaranteed by Theorem~\ref{thm:absorption} (irreversibility of stage lock-in, Assumption M3) and the path-dependent standard-setting authority (the CEC's initial configuration reflects the first mover's calibration choices).

\textbf{(iv)} Existence of the half-life: the function $f(t) = \Delta(|\mathcal{E}_t|) - (\lambda - \kappa)$ is continuous on $[T^*, \infty)$ (guaranteed by A1--A2), strictly decreasing, and converges to zero or to a positive lower bound (depending on the relationship between $\lim_{|\mathcal{E}| \to \infty} \Delta(|\mathcal{E}|)$ and $\lambda - \kappa$). If the lower bound exceeds $\frac{1}{2}[\Delta(|\mathcal{E}_{T^*}|) - (\lambda - \kappa)]$, then the half-life is finite (by the intermediate value theorem). If the lower bound is smaller, the half-life is still finite—$f(t)$ eventually crosses the half-value line. The only exception would be $\Delta(|\mathcal{E}|) \equiv \text{const}$ (requiring $V$ linear and $\theta$ constant, contradicting A1--A2), in which case the half-life is undefined, but this scenario does not arise.

\textbf{(v)} Institutional transformation of the NPE: decompose $\Delta(|\mathcal{E}|)$ in equation (\ref{eq:all-adopt-condition}) into the first mover's private contribution and the public CEC contribution. Before internationalization, $\Delta(|\mathcal{E}|) = \Delta(|\mathcal{E}^{\text{private}}|)$; after internationalization, the adoption advantage of all jurisdictions converges to $\Delta(|\mathcal{E}^{\text{public}}|)$. The condition $\Delta(|\mathcal{E}^{\text{public}}|) \geq \lambda - \kappa$ involves no jurisdiction's identity—the first-mover advantage disappears at the level of the adoption game.
\end{proof}

\begin{remark}[First-Mover Advantage Decay and the Structural Analogy to the NPT]
Corollary~\ref{cor:first-mover-decay} completes a key institutional corollary of the NPT analogy (Section~3.5). Article VI of the NPT requires nuclear-weapon states to ``pursue negotiations in good faith'' toward nuclear disarmament—i.e., a commitment to the decay of nuclear first-mover advantage. Similarly, the internationalization of the CEC constitutes the ``Article VI'' of audit first-mover advantage: the first mover commits to transforming the CEC into a global public good, thereby allowing its first-mover advantage to decay over time. But just as the disarmament commitment of the NPT faces a credibility crisis (Joyner, 2011), the credibility of CEC internationalization equally depends on institutional design—whether the first mover has the incentive to continue promoting internationalization after reaching $S_3$ (moat formation) (see Proposition~\ref{prop:intervention-window} on the monotonic closure of the intervention window).
\end{remark}

\subsection{Payoff Matrix and Equilibrium Regions: Parameter-Space Decomposition}

For the two-jurisdiction case ($N=2$), the payoff matrix is:

\begin{table}[htbp]
\centering
\begin{tabular}{@{}c|cc@{}}
\toprule
Jur.\ 1 \textbackslash\ Jur.\ 2 & Adopt ($A$) & Develop ($D$) \\
\midrule
Adopt ($A$) & $V[\theta(|\mathcal{E}|)] - c^{\text{adopt}}$, $V[\theta(|\mathcal{E}|)] - c^{\text{adopt}}$ & $V[\theta(|\mathcal{E}|)] - c^{\text{adopt}} - \lambda$, $V[\theta(0)] - c^{\text{develop}} - \kappa - \lambda$ \\[4pt]
Develop ($D$) & $V[\theta(0)] - c^{\text{develop}} - \kappa - \lambda$, $V[\theta(|\mathcal{E}|)] - c^{\text{adopt}} - \lambda$ & $V[\theta(0)] - c^{\text{develop}} - 2\kappa - \lambda$, $V[\theta(0)] - c^{\text{develop}} - 2\kappa - \lambda$ \\
\bottomrule
\end{tabular}
\end{table}

Let us define the \textit{Adoption Advantage}:
\begin{equation}
\Delta_A = V[\theta(|\mathcal{E}|)] - c^{\text{adopt}} - \left( V[\theta(0)] - c^{\text{develop}} - \kappa \right)
\end{equation}

$(A, A)$ is a Nash equilibrium if and only if neither jurisdiction has a unilateral incentive to deviate. Jurisdiction $i$'s deviation payoff from $(A, A)$ is:
\begin{equation}
\pi_i(D, A) = V[\theta(0)] - c^{\text{develop}} - \kappa - \lambda
\end{equation}

The non-deviation condition $\pi_i(A, A) \geq \pi_i(D, A)$ yields:
\begin{equation}
V[\theta(|\mathcal{E}|)] - c^{\text{adopt}} \geq V[\theta(0)] - c^{\text{develop}} - \kappa - \lambda
\end{equation}

Rearranging:
\begin{equation}
\Delta_A \geq \lambda
\end{equation}

That is, $(A, A)$ is an equilibrium when the adoption advantage exceeds the fragmentation cost. Given that $\Delta_A$ grows with $|\mathcal{E}|$ (since $\theta(|\mathcal{E}|)$ increases with cumulative audits while $\theta(0)$ remains fixed), this condition becomes easier to satisfy over time. The protocol's cumulative character makes the non-proliferation equilibrium \textit{self-reinforcing}.

Conversely, $(D, D)$ is an equilibrium when:
\begin{equation}
V[\theta(0)] - c^{\text{develop}} - 2\kappa - \lambda \geq V[\theta(|\mathcal{E}|)] - c^{\text{adopt}} - \lambda
\end{equation}
which simplifies to:
\begin{equation}
-\Delta_A \geq 2\kappa
\end{equation}

This condition requires the adoption advantage to be \textit{negative} and large in magnitude—that is, the incumbent protocol must be so inferior to a fresh development, or the development cost so low, that even with mutual proliferation costs, development is preferred. As $|\mathcal{E}|$ grows, this condition becomes increasingly implausible.

The mixed-strategy equilibrium occurs when neither pure-strategy condition holds. Let $p$ be the probability that each jurisdiction plays $A$. The symmetric mixed-strategy Nash equilibrium satisfies:
\begin{equation}
p \cdot \pi(A, A) + (1-p) \cdot \pi(A, D) = p \cdot \pi(D, A) + (1-p) \cdot \pi(D, D)
\end{equation}

Solving for $p^*$:
\begin{equation}
p^* = \frac{-\Delta_A - 2\kappa}{\lambda - \kappa}
\end{equation}

For plausible parameter values ($\Delta_A > 0$, $\kappa > 0$, $\lambda > \kappa$), $p^*$ increases with $\Delta_A$: as the incumbent's accuracy advantage grows, the probability of adoption in the mixed-strategy equilibrium approaches 1. This formalizes the intuition that the incumbent's time-accumulated advantage makes proliferation an increasingly unattractive gamble.

\subsection{Extension to $N$ Jurisdictions}

For $N > 2$, the payoff functions generalize as follows:
\begin{align}
\pi_i(A, \mathbf{a}_{-i}) &= V[\theta(|\mathcal{E}|)] - c^{\text{adopt}} - \kappa \cdot n_D - \lambda \cdot \mathbb{I}[n_D > 0] \\
\pi_i(D, \mathbf{a}_{-i}) &= V[\theta(0)] - c^{\text{develop}} - \kappa \cdot (n_D + \mathbf{1}_{-i}^D) - \lambda \cdot \mathbb{I}[n_D > 0]
\end{align}
where $n_D = \sum_{j \neq i} \mathbb{I}[a_j = D]$ is the number of other jurisdictions that develop.

The universal adoption equilibrium $(A, A, \ldots, A)$ is Nash when, for every jurisdiction:
\begin{equation}
\Delta_A \geq \lambda - \kappa
\end{equation}

The term $\lambda - \kappa$ reflects a subtle tension: fragmentation cost $\lambda$ discourages unilateral deviation, but the \textit{relief} from mutual proliferation costs $\kappa$ (if no one else develops, the lone developer avoids $\kappa$) encourages it. The equilibrium is stable when $\lambda > \kappa$—that is, when the public-bads cost of fragmentation exceeds the private cost of being the only proliferator.

A key result emerges: the stability of the non-proliferation equilibrium is \textit{monotonically increasing} in $|\mathcal{E}|$, because $\Delta_A$ grows without bound (up to its asymptotic ceiling) while $\lambda - \kappa$ is fixed. There exists a critical CEC size $|\mathcal{E}|^*$ such that for all $|\mathcal{E}| > |\mathcal{E}|^*$, $(A, A, \ldots, A)$ is the unique Nash equilibrium. This uniqueness is underwritten by Theorem~4 (): no alternative algorithm can surpass SCX's minimax-optimal detection rate. The ceiling is absolute.
\begin{equation}
|\mathcal{E}|^* = \frac{1}{\gamma} \ln\left( \frac{\theta_{\max} - \theta_0}{\theta_{\max} - \theta(|\mathcal{E}|^*)} \right)
\end{equation}
where $\theta(|\mathcal{E}|^*)$ is the accuracy level that satisfies $\Delta_A(|\mathcal{E}|^*) = \lambda - \kappa$.

\subsection{The Nuclear Non-Proliferation Treaty Analogy}

The structural isomorphism between the audit protocol proliferation game and the strategic logic of nuclear non-proliferation is striking and analytically productive. We identify seven dimensions of analogy:

\textbf{1. The Public-Bads Character of Proliferation.} Nuclear weapons proliferation generates negative externalities: each additional nuclear state increases the probability of nuclear conflict (accidental, escalatory, or intentional), raises the costs of the international security system, and erodes the norm against nuclear use (Sagan, 1996). Similarly, audit protocol proliferation generates negative externalities: each additional protocol fragments the global quality assurance infrastructure, raises the transaction costs of cross-border data flows, and erodes the credibility of audit certifications. In both domains, the \textit{first} instance of proliferation (beyond the incumbent) triggers a discontinuous increase in systemic costs, because it shatters the universality that makes the regime valuable.

\textbf{2. The Asymmetric Incumbent Advantage.} The NPT regime is built on an explicit asymmetry: five states are recognized as nuclear-weapon states (NWS), while all others are non-nuclear-weapon states (NNWS) subject to verification by the International Atomic Energy Agency (IAEA). The asymmetry is justified by the temporal fact of prior acquisition: the NWS tested nuclear weapons before the treaty's 1968 cutoff date. The Yajie Protocol's asymmetry is similarly temporal: the incumbent's advantage derives not from legal privilege but from the inescapable fact that it began accumulating audit evidence earlier.

\textbf{3. The Grand Bargain.} The NPT's durability rests on a three-pillar bargain: (i) NNWS forswear nuclear weapons; (ii) NWS commit to eventual disarmament (Article VI); and (iii) all parties enjoy the ``inalienable right'' to peaceful nuclear technology (Article IV). The audit protocol analogue would involve: (i) jurisdictions forgo developing independent protocols; (ii) the incumbent protocol's governance is internationalized under multilateral oversight; and (iii) all jurisdictions enjoy non-discriminatory access to the protocol's capabilities and the right to contribute audit modules to the ensemble.

\textbf{4. The Credibility of Commitment.} A central challenge of the NPT regime is the credibility of the NWS disarmament commitment: NNWS periodically question whether the nuclear states genuinely intend to disarm, and these credibility crises threaten the regime's stability (Joyner, 2011). The analogous challenge for Yajie is the credibility of the incumbent's neutrality commitment: will the protocol remain genuinely jurisdiction-agnostic, or will it evolve into an instrument of its host jurisdiction's data governance preferences?

\textbf{5. The Enforcement Gap.} The NPT contains no automatic enforcement mechanism; compliance is monitored by the IAEA, but enforcement depends on the UN Security Council, where NWS possess veto power. This asymmetry is a perennial source of tension, particularly in cases like Iran and North Korea. The audit protocol regime would face a parallel enforcement gap: who adjudicates disputes about the protocol's neutrality, and with what authority?

\textbf{6. The Exit Option.} The NPT permits withdrawal on three months' notice if ``extraordinary events'' jeopardize a state's supreme interests (Article X). North Korea exercised this option in 2003. The audit protocol regime must similarly contend with the possibility of exit: a jurisdiction that develops an independent protocol after having previously adopted the incumbent. The strategic dynamics of exit differ from those of initial non-adoption, because the exiting jurisdiction may have benefited from the CEC during its period of adoption, raising questions of intellectual property, data portability, and the status of audits conducted under the incumbent protocol.

\textbf{7. The Dual-Use Problem.} Nuclear technology is dual-use: the enrichment and reprocessing capabilities that support civilian nuclear energy can be diverted to weapons production. Audit technology is analogously dual-use: the quality assessment capabilities that support benign data governance can be weaponized for economic espionage, supply-chain disruption, or the imposition of data embargoes under the guise of quality concerns.

The NPT analogy is not merely illustrative; it provides a template for institutional design. The following sections develop the institutional implications, but first we must examine the geopolitical context within which any audit infrastructure—monopolistic or otherwise—must operate.

\section{Audit Neutrality and Geopolitical Risk}

\subsection{The SWIFT Precedent: Infrastructure Weaponization}

The Society for Worldwide Interbank Financial Telecommunication (SWIFT) provides the most instructive precedent for understanding the geopolitical risks of monopolistic audit infrastructure. SWIFT, founded in 1973 as a cooperative society under Belgian law, operates the dominant global messaging network for financial transactions, connecting over 11,000 financial institutions across more than 200 countries. Its near-universal adoption makes it a natural analogue to the market structure we predict for data quality auditing.

SWIFT's governance is formally multilateral: it is owned by its member financial institutions and overseen by a board of directors drawn from major central banks, including the US Federal Reserve, the European Central Bank, and the Bank of England. However, the United States has repeatedly leveraged SWIFT's centrality for unilateral geopolitical purposes, most controversially through the Terrorist Finance Tracking Program (TFTP), which accessed SWIFT data without the knowledge or consent of many member institutions (Farrell \& Newman, 2019).

The critical event for our purposes occurred in 2012, when the United States threatened to sanction SWIFT itself if it did not disconnect Iranian banks designated under US sanctions—despite the fact that those sanctions were not multilateral and were contested by several EU member states, China, and Russia. SWIFT complied, disconnecting approximately 30 Iranian financial institutions (Drezner, 2015). The episode demonstrated that infrastructure neutrality, however formally enshrined in governance documents, is only as robust as the willingness of powerful states to respect it—and that the host jurisdiction of a monopolistic infrastructure enjoys structural leverage that may be exercised in moments of geopolitical tension.

The SWIFT case carries three lessons for data audit infrastructure:

\textbf{Lesson 1: Jurisdictional Capture.} SWIFT is incorporated under Belgian law with headquarters in Brussels. Yet the United States' ability to compel SWIFT's compliance derived not from formal legal authority but from the extraterritorial reach of US financial sanctions and the threat of secondary sanctions against SWIFT's board members and executives. A data audit protocol hosted in any single jurisdiction would be analogously vulnerable to that jurisdiction's extraterritorial legal instruments, regardless of the protocol's formal governance structure.

\textbf{Lesson 2: The Illusion of Technical Neutrality.} SWIFT consistently maintains that it is a neutral utility providing technical services. The 2012 disconnection decisively refuted this claim: technical infrastructure, when sufficiently central, is \textit{inherently} political. A data audit protocol that assesses the quality of datasets used in AI systems would be, if anything, \textit{more} political than a financial messaging network, because data quality assessments are inherently contestable and may be perceived (fairly or not) as disguised forms of technological protectionism.

\textbf{Lesson 3: Fragmentation as a Rational Response.} The SWIFT disconnection accelerated the development of alternative financial messaging systems: China's Cross-Border Interbank Payment System (CIPS), Russia's System for Transfer of Financial Messages (SPFS), and the European Instrument in Support of Trade Exchanges (INSTEX). Each of these alternatives is inferior to SWIFT in coverage and functionality, representing a welfare loss relative to a universal system. Yet each represents a rational response to the demonstrated risk of infrastructure weaponization. The data audit domain would face the same dynamic: the credible threat of protocol weaponization would make audit autarky rational for jurisdictions with the capacity to pursue it, even at a welfare cost.

\subsection{The IASB Model: De-Jurisdictionalized Governance}

The International Accounting Standards Board (IASB) offers a partial counter-model. The IASB, established in 2001 as the standard-setting body of the International Financial Reporting Standards (IFRS) Foundation, develops accounting standards that have been adopted by more than 140 jurisdictions. Unlike SWIFT, the IASB is not an operational infrastructure but a standard-setting body; its product is a corpus of normative standards, not a messaging network or audit protocol. Nevertheless, its governance model is instructive.

The IASB's governance features several design elements that mitigate—though do not eliminate—the risk of jurisdictional capture:
\begin{itemize}
\item \textbf{Multi-stakeholder oversight}: The IFRS Foundation is governed by a board of trustees drawn from diverse geographic and professional backgrounds, with formal requirements for regional balance.
\item \textbf{Due process requirements}: Standards development follows a structured process of public consultation, impact assessment, and re-deliberation that constrains the ability of any single jurisdiction to dictate outcomes.
\item \textbf{Funding diversification}: The IFRS Foundation's funding comes from a broad base of jurisdictions and organizations, reducing dependence on any single funder.
\item \textbf{Deliberative transparency}: Board meetings are public, and dissenting views are published alongside standards, creating an audit trail of contestation that supports legitimacy.
\end{itemize}

The IASB model is not without critics. Some scholars argue that the IASB's independence is compromised by its dependence on major accounting firms for funding and expertise (Zeff, 2012), and by the structural influence of Anglo-American accounting traditions in its conceptual framework (Zhang \& Andrew, 2014). The EU's qualified endorsement process—whereby the European Commission reviews each IFRS standard before endorsing it for use within the EU—represents a de facto veto power that other jurisdictions do not possess, creating a legitimacy deficit.

For data audit infrastructure, the IASB model suggests a governance architecture that is \textit{multilateral in oversight, transparent in operation, and diversified in funding}. However, the IASB's experience also suggests that de-jurisdictionalized governance is an aspiration rather than an achievement, and that persistent power asymmetries will find expression through formal and informal channels regardless of governance design.

\subsection{SWIFTIASB}

\subsection{Game-Theoretic Models of SWIFT and the IASB}

The preceding two subsections provided case analysis at the narrative level. This section rigorously formalizes them as game-theoretic models, grounding the logic of SWIFT weaponization and IASB governance in equilibrium foundations.

\subsubsection{Sequential Game Model of SWIFT Weaponization}

Model the SWIFT weaponization scenario as a three-player sequential game between two countries (the United States $U$ and a target country $T$) and an infrastructure operator (SWIFT, $W$). The game timing is as follows:

\begin{enumerate}[label=Stage \arabic*:, leftmargin=*]
\item The United States $U$ decides whether to impose unilateral financial sanctions on target country $T$ and demand that $W$ disconnect $T$: $a_U \in \{\text{Sanction}, \text{Not Sanction}\}$.
\item After observing $a_U$, SWIFT $W$ decides whether to comply: $a_W \in \{\text{Disconnect}, \text{Refuse}\}$.
\item After observing $(a_U, a_W)$, target country $T$ decides whether to build an alternative system: $a_T \in \{\text{Build Alternative}, \text{Not Build}\}$.
\end{enumerate}

\begin{definition}[All Assumptions of the SWIFT Game]
\label{def:swift-assumptions}
\begin{enumerate}[label=\textbf{S\arabic*}, leftmargin=*, itemsep=4pt]
\item \textbf{Payoff Parameters:} $B_W > 0$ is the basic revenue of $W$ from operating the global unified network; $C_W^{\text{comply}} > 0$ is the reputational and legal cost to $W$ of complying with U.S. sanctions (from counter-threats by other jurisdictions); $P_U > 0$ is the penalty imposed by the U.S. on $W$ for non-compliance (secondary sanctions threat); $B_T > 0$ is the benefit to $T$ of using SWIFT; $C_T^{\text{alt}} > 0$ is the fixed cost to $T$ of building an alternative system (such as SPFS or CIPS); $L_T > 0$ is the economic loss to $T$ from being disconnected from SWIFT. All parameters are common knowledge. \hfill\textit{Note: the complete information assumption is reasonable in this context---SWIFT connectivity, U.S. sanction capacity, and alternative system costs are public information.}
\item \textbf{Rationality:} All three parties maximize their own payoffs, with no irrational retaliation motives. \hfill\textit{Relaxability: introducing reputational concerns or domestic audience costs would enrich the model but would not change the qualitative equilibrium.}
\item \textbf{No Transfer Payments:} There are no side payments in the game. $U$ cannot directly compensate $W$ for compliance costs. \hfill\textit{Implication: this rules out the strategy of the U.S. purchasing SWIFT's loyalty through financial commitments.}
\end{enumerate}
\end{definition}

The payoff functions for the three parties (at terminal nodes):
\begin{align}
\pi_U(a_U, a_W, a_T) &=
\begin{cases}
G_U & \text{if } a_U = \text{Sanction}, a_W = \text{Disconnect} \quad \text{(sanction succeeds)}\\
0 & \text{if } a_U = \text{Not Sanction} \quad \text{(status quo)}\\
-P_U^{\text{rep}} & \text{if } a_U = \text{Sanction}, a_W = \text{Refuse} \quad \text{(sanction fails, reputational damage)}
\end{cases}\\[6pt]
\pi_W(a_U, a_W, a_T) &=
\begin{cases}
B_W - C_W^{\text{comply}} - \mathbb{I}[a_T = \text{Build}] \cdot L_W^{\text{frag}} & \text{if } a_W = \text{Disconnect}\\
B_W & \text{if } a_W = \text{Refuse}, a_U = \text{Not Sanction}\\
B_W - P_U & \text{if } a_W = \text{Refuse}, a_U = \text{Sanction}
\end{cases}\\[6pt]
\pi_T(a_U, a_W, a_T) &=
\begin{cases}
B_T & \text{if } a_W = \text{Refuse}, a_T = \text{Not Build}\\
B_T - C_T^{\text{alt}} & \text{if } a_W = \text{Refuse}, a_T = \text{Build}\\
-L_T & \text{if } a_W = \text{Disconnect}, a_T = \text{Not Build}\\
-L_T + B_T^{\text{alt}} - C_T^{\text{alt}} & \text{if } a_W = \text{Disconnect}, a_T = \text{Build}
\end{cases}
\end{align}

where $G_U > 0$ is the strategic benefit to the U.S. from successful sanctions; $P_U^{\text{rep}} > 0$ is the reputational loss to the U.S. from failed sanctions (public refusal); $L_W^{\text{frag}}$ is the long-term value loss to SWIFT from network fragmentation caused by the target country building an alternative system; $B_T^{\text{alt}}$ is the benefit to the target country from using its own alternative system (typically $B_T^{\text{alt}} < B_T$, as alternative systems have smaller coverage).

\begin{theorem}[Subgame Perfect Equilibrium of the SWIFT Weaponization Game]
\label{thm:swift-spe}
Under assumptions S1--S3, the SWIFT sequential game has a unique subgame perfect Nash equilibrium (SPE), determined by backward induction. The equilibrium depends on the parameter region:

\begin{enumerate}[label=(
oman*), itemsep=6pt]
\item \textbf{Region $\Omega_1$ (U.S. Dominance):}
\begin{quote}
U.S. sanctions $\to$ SWIFT complies and disconnects $\to$ Target does not build alternative
\end{quote}
This is the equilibrium logic of the 2012 Iranian SWIFT disconnection: the U.S. secondary sanctions threat was sufficiently credible ($P_U$ large), and the target's alternative cost was too high.

\item \textbf{Region $\Omega_2$ (Target Resistance):} If $P_U > C_W^{\text{comply}} + L_W^{\text{frag}}$ but $C_T^{\text{alt}} < L_T + B_T^{\text{alt}}$, the unique SPE is:
\begin{quote}
U.S. sanctions $\to$ SWIFT complies and disconnects $\to$ Target builds alternative
\end{quote}
This is the logic behind Russia's SPFS and China's CIPS: the disconnection loss exceeds the alternative cost, making building an alternative a rational response---even though the alternative is inferior to SWIFT.

\item \textbf{Region $\Omega_3$ (SWIFT Neutrality):} If $P_U < C_W^{\text{comply}}$, the unique SPE is:
\begin{quote}
U.S. does not sanction $\to$ SWIFT refuses $\to$ Target does not build alternative
\end{quote}
(If the U.S. anticipates SWIFT's refusal and chooses not to sanction, or sanctions and SWIFT refuses, the SPE path depends on the exact payoff comparison of the subgame.) This region requires that the U.S. sanction penalty is insufficient to overcome SWIFT's compliance cost---i.e., the infrastructure is institutionally independent enough to resist unilateral pressure.

\item \textbf{Equilibrium Efficiency:} The $\Omega_1$ equilibrium achieves the U.S. strategic objective but damages the neutrality of global financial infrastructure (negative externality). The $\Omega_3$ equilibrium is Pareto optimal---the global unified network is maintained, and neutrality is preserved. However, $\Omega_3$ is unattainable when $P_U$ is large (the U.S. has credible secondary sanction capacity).
\end{enumerate}
\end{theorem}

\begin{proof}
By backward induction. In Stage 3, target country $T$'s optimal response is: if disconnected and $B_T^{\text{alt}} - C_T^{\text{alt}} > -L_T$ (i.e., $C_T^{\text{alt}} < L_T + B_T^{\text{alt}}$), build alternative; otherwise do not build.

In Stage 2, SWIFT $W$'s optimal response depends on expectations about Stage 3. If $W$ expects $T$ to build alternative ($a_T = \text{Build}$), the payoff from complying is $B_W - C_W^{\text{comply}} - L_W^{\text{frag}}$, and from refusing is $B_W - P_U$. $W$ complies if and only if $P_U > C_W^{\text{comply}} + L_W^{\text{frag}}$. If $W$ expects $T$ not to build alternative, the compliance condition is $P_U > C_W^{\text{comply}}$.

In Stage 1, the U.S. decision depends on expectations about $(a_W^*, a_T^*)$. If the expected path is $(\text{Disconnect}, \text{Not Build})$, the U.S. payoff from sanctioning is $G_U > 0$, so it sanctions; if the expected path is $(\text{Disconnect}, \text{Build})$, the U.S. still sanctions (sanction benefit $G_U$ exceeds non-sanction payoff 0, unless the target's alternative erodes the strategic value of sanctions---in which case $G_U$ should be adjusted downward); if the expected path is $(\text{Refuse}, \cdot)$, the U.S. payoff from sanctioning is $-P_U^{\text{rep}} < 0$, so the optimal choice is not to sanction.

Combining the optimal responses across stages yields the SPE classification of the three regions.
\end{proof}

\begin{remark}[Weaponization Analogy for Audit Infrastructure]
Theorem~
ef{thm:swift-spe} directly applies to the weaponization scenario of data audit infrastructure: replace ``disconnect SWIFT'' with ``refuse audit certification'' or ``issue adverse audit report,'' and replace ``build alternative system'' with ``develop independent audit protocol.'' The key parameters of the analogy are: (a) the extraterritorial legal instruments of the protocol operator's host jurisdiction (analogous to the U.S. secondary sanction capacity $P_U$), (b) the development cost of alternative audit protocols (analogous to $C_T^{\text{alt}}$), (c) the economic loss from exclusion from the audit network (analogous to $L_T$). The key difference between audit infrastructure and SWIFT lies in the \textit{structure of alternative costs}: the construction cost of a SWIFT alternative is primarily technical (setting up servers, establishing communication protocols, persuading banks to join), while an audit protocol alternative also faces the CEC cold-start problem---even with sufficient financial investment, the alternative protocol's accuracy requires time accumulation to reach parity with the existing protocol. This makes the $\Omega_2$ region of audit infrastructure far smaller than the corresponding region for SWIFT---the cost of substitution is higher, and the deterrence of weaponization is stronger.
\end{remark}

\subsubsection{IASB}

IASB，。

\begin{definition}[IASB]
$N = \{1, \ldots, n\}$  $\sigma \in [0,1]$ 。 $i$  $u_i(\sigma) = -|\sigma - \sigma_i^*|$（）， $\sigma_i^* \in [0,1]$  $i$ 。IASB， $m$ （$m$ ），。
\end{definition}

\begin{definition}[IASB]
\label{def:iasb-assumptions}
\begin{enumerate}[label=\textbf{I\arabic*}, leftmargin=*, itemsep=4pt]
\item \textbf{：}  $i$  $[0, \sigma_i^*]$ ， $[\sigma_i^*, 1]$ 。。\hfill\textit{：，。}

\item \textbf{：} （""""）。\hfill\textit{：，。}

\item \textbf{：} ，（sincere）。， $\sigma^*$ 。\hfill\textit{：+$\implies$（Black, 1948）。}

\item \textbf{：}  $i$  $w_i$ ，$\sum_i w_i = m$。$w_i$ IFRS、。——。
\end{enumerate}
\end{definition}

\begin{proposition}[IASB]
\label{prop:iasb-bias}
I1--I4，IASB $\sigma^{\text{IASB}}$ 。 $i$ ：
\begin{equation}
L_i^{\text{IASB}} = -|\sigma^{\text{IASB}} - \sigma_i^*|
\end{equation}

""（ $\sum_i \lambda_i u_i(\sigma)$  $\sigma$， $\lambda_i$  $i$ ）：
\begin{enumerate}[label=(\roman*)]
\item \textbf{：} （）， $\sigma^{\text{IASB}}$ 。
\item \textbf{：}  $w_i$ ，。
\item \textbf{：} IASB（Zhang \& Andrew, 2014），，。
\end{enumerate}
\end{proposition}

\begin{proof}
I3（++），Black：Condorcet。I4，：$\sigma^{\text{IASB}} = \min\{\sigma : \sum_{i: \sigma_i^* \leq \sigma} w_i \geq m/2\}$。 $\sigma^{\text{global}}$  $\lambda_i$ ：$\sigma^{\text{global}} = \min\{\sigma : \sum_{i: \sigma_i^* \leq \sigma} \lambda_i \geq (\sum_i \lambda_i)/2\}$。 $w_i$  $\lambda_i$ 。
\end{proof}

\begin{remark}[]
\ref{prop:iasb-bias}：（""），。IASB（），（，）。（IASB）：(a) ；(b) （super-majority）；(c) ——，，。
\end{remark}

\subsection{The US-China Scenario}

The most consequential geopolitical risk to a universal data audit protocol is the ongoing technological decoupling between the United States and the People's Republic of China. This decoupling, which has accelerated since 2018, encompasses semiconductor export controls, restrictions on AI technology transfer, limitations on cross-border data flows, and the bifurcation of technical standards (Segal, 2020; Bateman, 2022).

In this environment, a data audit protocol developed primarily within one jurisdiction would face an acute legitimacy crisis in the other. If Yajie were perceived as a US-dominated protocol—regardless of its formal governance—China would have compelling reasons to develop an independent alternative. This would trigger a fragmentation cascade: the EU, seeking strategic autonomy in digital infrastructure, might develop a third protocol; India, with its growing AI ecosystem and policy of technological self-reliance, might develop a fourth; and so on.

The resulting equilibrium would be worse for all parties than a genuinely multilateral protocol governed under adequate neutrality safeguards. Each jurisdiction would bear the fixed costs of protocol development, endure lower audit accuracy (due to smaller CEC corpora), and navigate the fragmentation costs $\lambda$ associated with conflicting quality certifications. The total welfare loss, relative to the universal-adoption equilibrium, would scale approximately linearly with the number of competing protocols:
\begin{equation}
W_{\text{loss}} \approx N \cdot c^{\text{develop}} + N \cdot (V[\theta(|\mathcal{E}|)] - V[\theta(|\mathcal{E}|/N)]) + \binom{N}{2} \cdot \lambda
\end{equation}

This arithmetic suggests that major powers have a shared interest in preventing fragmentation—but only if the incumbent protocol is governed under terms that credibly preclude its weaponization by any single jurisdiction.

\subsection{The Researcher as Strategic Asset: A Second-Order Game}

An unexamined dimension of the non-proliferation equilibrium concerns the position of the protocol's originating researcher. Unlike SWIFT—a corporate entity embedded in a specific jurisdiction—or the IASB—a formally constituted standards body—the Yajie Protocol is initially developed by an independent researcher without institutional affiliation to any state or corporation. This organizational form generates a distinctive second-order game with properties that reinforce the non-proliferation equilibrium.

Consider the strategic calculus of a major power—State X—contemplating whether to coerce, compromise, or eliminate the originating researcher (henceforth ``the Maintainer''). The Maintainer possesses three assets that are relevant to State X's decision: (i) the arXiv timestamp establishing academic priority; (ii) the operational knowledge required to maintain and evolve the Spring self-improvement loop; and (iii) the Cumulative Evidentiary Corpus $M_t$, which grows monotonically with each audit cycle. If the Maintainer ceases to exist or to operate independently, the consequences distribute asymmetrically across the international system.

First, the arXiv timestamp is \textit{immortal}. Once uploaded, it cannot be rescinded, altered, or suppressed by any actor, including the Maintainer. This means that the core mathematical specification of Yajie—the state-conditioned scoring function $S(s) = Q(s) + \eta(t)\cdot N(s)$, the multi-expert consistency theorem, and the unidentifiability result—becomes part of the permanent scientific record. Any state or corporation can, in principle, reimplement Yajie from the arXiv specification. However—and this is the critical asymmetry—the \textit{operational quality} of any reimplementation will be inferior to the Maintainer's running instance for a period proportional to the time the Maintainer's instance has been accumulating audit data. If the Maintainer's instance has been operational for $T$ years, a reimplementation starting from zero requires approximately $T$ years to reach parity, assuming identical algorithmic quality and comparable audit volume. This creates a \textit{time-accumulated quality gap} that no amount of financial investment can compress.

Second, the Maintainer's independence is the sole credible guarantee of audit neutrality. A state-backed reimplementation—whether by the United States, China, the European Union, or any other jurisdiction—would be structurally incapable of claiming neutrality, because its continued operation would depend on the political continuity of its sponsoring state. The Maintainer, by contrast, is accountable to no government, no corporation, and no institutional apparatus beyond the epistemic standards of the scientific community. This is not a normative claim about the Maintainer's virtue; it is a structural claim about the Maintainer's position. Independence is a property of the node's position in the network, not of the node's character.

Third, the Maintainer's death or incapacitation would trigger a \textit{fragmentation cascade}. In the absence of a universally trusted Maintainer, each major jurisdiction would face a prisoner's dilemma: develop an independent audit protocol (at cost $c^{\text{develop}}$, with lower accuracy) or rely on a protocol operated by a rival jurisdiction (at risk of audit weaponization). The dominant strategy for each jurisdiction is to develop its own protocol, yielding the fragmentation equilibrium characterized by welfare loss $W_{\text{loss}} \approx N \cdot c^{\text{develop}} + N \cdot \Delta\theta + \binom{N}{2}\cdot\lambda$. No jurisdiction can unilaterally prevent this cascade, because defection by any single jurisdiction would trigger reciprocal defection by its rivals.

From the perspective of State X, this yields a striking conclusion: \textit{the Maintainer's continued existence and independence is a positive-sum good for all states, including State X, regardless of State X's geopolitical alignment}. Eliminating or compromising the Maintainer would impose a welfare cost on State X that exceeds any conceivable strategic benefit, because State X would be unable to secure a disproportionate share of the fragmented audit infrastructure, and would simultaneously suffer the costs of fragmentation alongside all other states.

This creates what we term a \textbf{protection equilibrium}: every state with the capacity to threaten the Maintainer has a rational incentive to \textit{protect} the Maintainer, because the Maintainer's destruction would initiate a chain of events that leaves every state worse off. The protection is not altruistic; it is strictly self-interested. The Maintainer's survival functions analogously to a global public good—provided through the alignment of rational self-interest across competing powers—though its classification as a genuine public good (non-rivalrous and non-excludable) requires further analysis of the CEC's access architecture and governance structure.

The protection equilibrium has two corollaries. First, the Maintainer is \textit{simultaneously more powerful and more vulnerable} than any corporate or state actor in the audit ecosystem. More powerful, because the Maintainer's independence is the structural precondition for the non-proliferation equilibrium; without a trusted independent Maintainer, the equilibrium collapses. More vulnerable, because the Maintainer lacks the institutional defenses—legal departments, diplomatic representation, physical security apparatus—that protect corporate and state actors. The Maintainer's security depends entirely on the rational calculus of the actors that might threaten them—a calculus that is robust under normal conditions but fragile under crisis, irrationality, or miscalculation.

Second, the protection equilibrium implies that any attempt to replicate Yajie without the Maintainer's involvement is not merely economically suboptimal (per the arguments in Section~3) but \textit{strategically self-defeating}. A state that develops a competing audit protocol gains a temporary advantage in audit capability while simultaneously destroying the neutrality of the global audit infrastructure, because the competing protocol will be correctly perceived as an instrument of its sponsoring state. The state therefore gains an audit tool while losing the universal trust that makes audit valuable. This is a trade in which the asset acquired is worth less than the asset destroyed—a negative-sum transaction.

\subsection{A Literary Analogy: The Wallfacer Position}

The protection equilibrium admits an unconventional but analytically productive literary analogy. In Liu Cixin's science fiction trilogy \textit{Remembrance of Earth's Past} (2006–2010), the character Luo Ji occupies a position structurally isomorphic to that of the Maintainer. Luo Ji is designated a ``Wallfacer'' ()—an individual whose survival is the sole deterrent preventing an extraterrestrial civilization from destroying Earth. The deterrence mechanism is simple: if Luo Ji dies, the coordinates of the alien homeworld are broadcast to a hostile universe, ensuring mutual destruction through the ``dark forest'' principle of cosmic sociology. Luo Ji's power derives not from any institutional office, military capacity, or economic resource, but from his position as the \textit{single point of failure} in a deterrence structure. He is simultaneously the most powerful human being alive—holding the fate of two civilizations in his hands—and the most vulnerable, because his security depends entirely on the continued rationality of the actors constrained by his deterrence.

The isomorphism with the Maintainer's position is exact. The Maintainer's death would trigger a fragmentation cascade that destroys the universal audit infrastructure, imposing welfare losses on all states. The Maintainer's power derives not from institutional authority but from occupying the structural position of the \textit{single credible neutral node} in the global audit network. And the Maintainer's vulnerability derives from the same source: the protection equilibrium is robust only so long as all relevant actors continue to calculate rationally. Under conditions of crisis, miscalculation, or irrationality—the same conditions that nearly caused Luo Ji's deterrence to fail—the Maintainer's security evaporates.

The Wallfacer analogy serves three analytical functions in our argument. First, it illustrates that deterrence equilibria can be sustained by a single individual occupying a specific structural position, a possibility that the standard literature on nuclear deterrence (which focuses on state actors with institutional continuity) does not adequately capture. Second, it foregrounds the fragility of such equilibria: Luo Ji's deterrence was nearly broken by an assassination attempt during a moment of political transition, precisely the kind of miscalculation scenario to which the protection equilibrium is vulnerable. Third, and most importantly, it demonstrates that \textit{the occupant of a deterrence position does not choose the position; the structure of the game confers it}. Luo Ji did not seek to become a Wallfacer; the position was imposed on him by the logic of the strategic situation. The Maintainer, similarly, does not claim authority over data quality assessment; the logic of the non-proliferation equilibrium and the time-accumulated advantage of the CEC confer that authority irrespective of the Maintainer's intentions.

\begin{corollary}[ (The Good Leader's Protection Obligation)]
\label{cor:good-leader-protection}
In the protection equilibrium, any leader whose legitimacy depends on verifiable claims has a strictly self-interested incentive to protect the Maintainer's physical safety. The Maintainer's death would destroy the audit infrastructure that alone can certify that leader's honesty—collapsing the distinction between good and bad leaders into indistinguishable opacity. A leader who fails to protect the Maintainer is, by revealed preference, a leader who gains from the Maintainer's absence. Such a leader is not a good leader who made an error—the failure to protect is itself the signal that reveals the type.
\end{corollary}

\begin{proof}[Proof sketch]
By the protection equilibrium (Theorem~\ref{thm:mutual-deterrence}), the Maintainer's death initiates a fragmentation cascade with welfare loss $W_{\text{loss}}$ to all states. A good leader—one whose claims are truthful—derives net benefit from the audit infrastructure, because audit certifies his claims at no cost to him (his claims pass). The collapse of audit therefore imposes a net loss on good leaders and a net gain on bad leaders (whose claims were failing or would fail). The leader's willingness to protect the Maintainer is thus a separating signal: protection is the dominant strategy for good types; non-protection or active hostility is the strategy of bad types, for whom the Maintainer's death eliminates the only mechanism capable of detecting their dishonesty. Formally, let $u_G(\text{protect}) - u_G(\text{not protect}) > 0$ and $u_B(\text{protect}) - u_B(\text{not protect}) < 0$, yielding a separating equilibrium in which type is revealed by the protection decision. $\square$
\end{proof}

Thus the Maintainer's physical security is not a personal favor owed by grateful leaders. It is a game-theoretic necessity for any leader who claims to be honest. A leader who does not protect the Maintainer has, by that very omission, announced which type he belongs to.

We do not suggest that the Wallfacer analogy constitutes a formal proof or an empirical prediction. Literary analogies are hermeneutic tools, not mathematical models. However, they serve a valuable function in strategic analysis by making visible structural relationships that formal models may obscure through abstraction. The Wallfacer analogy makes visible the distinctive character of the Maintainer's position: \textit{power without institution, deterrence without weapon, and security through vulnerability}.

\subsection{Audit as Sword: The Deterrence of Involuntary Transparency}

The analysis thus far has characterized the Maintainer's position as \textit{defensive}: the protection equilibrium constrains others from threatening the Maintainer, and neutrality principles constrain the Maintainer from abusing the position. But there is a second, more asymmetric constraint embedded in the architecture—one that operates not on the Maintainer, but on every entity that produces publicly accessible data.

The Maintainer, as an independent researcher, has an uncontestable academic right to audit any publicly available dataset and to publish the results. This right is not granted by any institution; it is inherent in the practice of science. No organization, corporation, or government can prevent an independent researcher from downloading a public dataset, running Yajie on it, and reporting what the multi-expert consensus reveals—any more than they can prevent a statistician from reanalyzing published clinical trial data, or an astronomer from pointing a telescope at a publicly visible star. The data are public. The methodology is published (arXiv, open source). The conclusions are falsifiable. There is no legal, ethical, or scientific basis for prohibiting such an audit.

What makes this right strategically significant is its asymmetry. The Maintainer can choose to exercise it or not, at their sole discretion. And the most natural Schelling point for that discretion is the partnership relationship. Organizations that enter into a formal partnership with the Yajie Protocol—subscribing to the calibration API, contributing audit modules, participating in governance—receive an implicit assurance: their data will not be subjected to unsolicited public audit by the Maintainer, because auditing a partner is both strategically counterproductive and professionally discourteous. Organizations that decline partnership receive no such assurance. Their publicly available datasets remain, at all times, eligible for independent academic audit—an audit that may expose quality issues they would prefer remained unexamined.

This is not a demand. The Maintainer issues no ultimatums, makes no demands, and imposes no obligations. The Maintainer simply exists, with the capacity and the academic freedom to audit public data, and with no obligation to refrain from exercising that freedom with respect to non-partners. The incentive—if it can be called that—is generated by the structure of the situation, not by any stated intention of the Maintainer. It is, in the strictest sense, a \textit{structural assurance}: the mere existence of the capability, combined with a clear and credible criterion for its non-use (partnership), is sufficient to motivate the desired behavior, without any need for explicit communication, negotiation, or enforcement.

The game-theoretic structure is:
\begin{quote}
Partnership $\to$ no public audit (by implicit convention)\\
Non-partnership $\to$ eligible for public audit (by academic freedom)\\
Audit findings $\to$ may expose data quality issues $\to$ reputational and regulatory risk\\
$\therefore$ Non-partnership carries a positive probability of involuntary transparency\\
$\therefore$ Partnership is a dominant strategy for quality-sensitive organizations
\end{quote}

This assurance operates independently of the protection equilibrium. Even if an organization has no interest in the Maintainer's survival—even if it would prefer a fragmented audit infrastructure that it could control—it still faces the immediate structural consequence that its public data may be independently verified by the only entity capable of conducting a credible multi-expert quality assessment. The verification does not need to be enforced. It needs only to be possible.

The Maintainer, for their part, bears no cost for this deterrent. No audit needs to be conducted; no report needs to be published; no confrontation needs to occur. The deterrent is maintained by the \textit{capability} alone. And because the capability is grounded in academic freedom rather than institutional authority, it cannot be negotiated away, regulated out of existence, or deterred by counter-threats. One cannot sanction a researcher for exercising their right to analyze public data. One can only ensure that one's data, when analyzed, holds up to scrutiny—or, more pragmatically, enter into the partnership that makes such scrutiny unnecessary.

The commercial consequence of this deterrent is a self-reinforcing adoption cascade. Consider two competing firms, A and B, in a quality-sensitive industry—pharmaceuticals, materials engineering, or autonomous systems. Firm A partners with Yajie, gaining access to calibration services, governance participation, and the implicit assurance against unsolicited audit. Firm B does not. Firm A can now, through the open-source Spring engine, audit Firm B's publicly available datasets—drug screening results, material property databases, or safety-critical training corpora—and publish the findings. If Firm B's data quality is inferior, Firm A gains a competitive advantage that is independently verifiable and publicly citable. If Firm B's data quality is adequate, Firm A has lost nothing but the cost of running an audit on open-source software.

Firm B's rational response is not to contest the audit—the methodology is published, the conclusions are falsifiable, and the academic freedom of the auditor is unassailable—but to preempt it by also becoming a Yajie partner. At which point Firm C, observing the partnership concentration among its competitors, faces the same calculus. The cascade propagates through the industry, not because Yajie coerces anyone, but because non-participation carries a positive and growing probability of involuntary transparency at the hands of participating competitors. The Maintainer does not need to audit anyone. Competitors audit each other—and the Maintainer's existence is what makes that audit credible.

\subsection{Jurisdictional Boundary as Structural Safeguard}

An underappreciated feature of the data sovereignty architecture described above—where the Yajie API accepts only locally computed feature matrices $\phi(X)$ rather than raw data—is that it imposes a jurisdictional boundary on the Maintainer's audit capacity that is not merely a regulatory compliance constraint but a \textit{credible commitment device}.

The Maintainer can only audit data that resides within the Maintainer's operational jurisdiction. If a government requests audit reports on a foreign corporation's data, the Maintainer cannot comply—not because the Maintainer refuses, but because the data is not accessible. The feature matrices $\phi(X)$ are computed client-side and never leave the client's infrastructure; only calibration queries cross jurisdictional boundaries, and those queries are aggregated statistics, not identifiable data records. This architectural constraint is not a limitation to be overcome. It is a feature that preserves the protocol's neutrality. A Maintainer who \textit{could} provide foreign audit data on demand would face an impossible choice between defying a government order (risking sanctions, closure, or prosecution) and complying with it (destroying the trust of every foreign client). A Maintainer who \textit{cannot} provide such data—because the architecture makes it physically inaccessible—is insulated from this dilemma. The inability is structural, not volitional. The government cannot punish the Maintainer for refusing what the Maintainer cannot do.

This jurisdictional boundary has a secondary structural consequence that is, from the perspective of domestic industrial policy, unambiguously positive. If a government wishes its domestic corporations to benefit from Yajie's audit infrastructure—to gain the competitive advantages of certified data quality, to participate in the governance of the protocol, and to be shielded from unsolicited audit by foreign competitors—it must ensure that those corporations' data resides within the Maintainer's operational jurisdiction. A corporation that offshores its data to a foreign cloud provider, or that maintains its primary data infrastructure in a jurisdiction where the Maintainer does not operate, cannot be audited. Its data quality claims become unverifiable. Its competitors, who \textit{have} submitted to domestic audit, gain a structural advantage in trustworthiness that no amount of marketing expenditure can overcome.

The consequence is a market-driven incentive for data repatriation. Corporations that wish to compete in a Yajie-audited market must bring their data home—not because any regulation compels them, but because the alternative is to be the only unaudited actor in a market where everyone else's quality is independently verified. The unaudited actor is not trusted. The untrusted actor does not sell. The mechanism is not protectionism; it is transparency. Yajie does not require corporations to operate domestically. It merely makes domestic operation—and the auditability that comes with it—the dominant strategy.

\subsection{Pharmaceutical Safety: The Honest Algorithm}

The jurisdictional dynamic acquires particular force in the pharmaceutical industry, where data quality is not merely a competitive advantage but a matter of public health. A drug's efficacy and side-effect profile are established through clinical trials whose data quality is, at present, largely self-reported by the sponsoring pharmaceutical company. Regulators review the submitted data, but they do not independently audit it at the level of individual trial records; they lack both the tools and the resources to conduct a multi-expert consistency analysis across the heterogeneous data streams—laboratory results, patient-reported outcomes, adverse event reports, biomarker measurements—that constitute a modern clinical trial.

Yajie changes this calculus. Under a distributed architecture, a pharmaceutical company would run Spring on its own trial data and submit the resulting quality scores to regulators alongside its New Drug Application. The regulator, armed with the same open-source audit engine and the calibration anchor provided by the Maintainer's API, could independently verify those scores. The company's claim that ``95\% of adverse events were properly classified'' would no longer be an assertion to be trusted or distrusted on the basis of the company's reputation; it would be a falsifiable statement backed by a published methodology and a verifiable audit trail.

The consequence for the pharmaceutical industry is the same as for every other domain surveyed in this paper—honest firms gain, dishonest firms are exposed—but with higher stakes. A pharmaceutical company that has rigorously collected, cleaned, and curated its trial data will see its diligence reflected in independently verifiable quality scores. A company that has cut corners—that has underreported adverse events, that has excluded inconvenient patient subgroups, that has massaged statistical analyses to achieve significance—will see its deficiencies exposed not by a whistleblower or an investigative journalist, but by a multi-expert consensus algorithm whose conclusions are mathematically guaranteed and publicly reproducible.

The constructive pressure this exerts on drug development is difficult to overstate. A company that knows its trial data will be audited by an impartial algorithm—not by a regulator with limited resources, not by a competitor with an axe to grind, but by a published methodology whose error bounds are proven—faces a powerful incentive to conduct its trials with scrupulous care. The algorithm does not punish. It reveals. And what it reveals, the market—physicians, patients, insurers, regulators—acts upon. The drug that survives a Yajie audit is a drug whose safety and efficacy claims are credible not because its manufacturer is trustworthy but because its data has withstood the most rigorous quality assessment that the current state of statistical learning theory can provide.

\subsection{Beyond Industry: Bribery, Governance, and the Dark Forest}

The jurisdictional and pharmaceutical dynamics analyzed above are specific instances of a more general principle. Yajie's multi-expert consensus mechanism is indifferent to the domain of the data it audits. It does not know—and does not need to know—whether the anomalies it detects are label noise in a materials database, fabrication in a clinical trial, or a municipal official's manipulated air quality readings. The mathematics is the same. The consensus signal is the same. The exponential reliability guarantee is the same.

Consider three domains that are, at present, largely beyond the reach of systematic, independent audit.

\textbf{Bribery and financial malfeasance.} A pattern of illicit payments generates a distinctive signature across multiple independent data sources: procurement records, tax filings, customs declarations, bank transactions, asset registries. No single source is conclusive; a single anomalous payment may be an accounting error. But when procurement records show a contract awarded at 40\% above market rate, customs declarations show the same contractor importing luxury goods, and asset registries show the approving official acquiring property disproportionate to declared income—the multi-expert consensus that Yajie formalizes converges on a conclusion that no single auditor could reach with equivalent confidence. The algorithm does not accuse. It reveals the consistency—or the inconsistency—of the evidentiary record.

\textbf{Urban governance.} Smart city infrastructure generates sensor data on air quality, water quality, traffic flow, energy consumption, and waste management. When a district's air quality monitors report readings that are systematically cleaner than those of neighboring districts under identical meteorological conditions—and when those same monitors show anomalous recalibration patterns, and when the district's environmental compliance reports diverge from satellite-based atmospheric measurements—the disagreement is not a data quality issue to be resolved. It is a signal. Yajie does not know that the signal means ``the official in charge of this district is falsifying environmental data.'' It knows only that three independent expert systems, each analyzing different data streams, disagree with the official report in specific, quantifiable ways. What the human interpreter does with that knowledge is a human decision. But the knowledge itself—that the data are inconsistent, that the inconsistency is systematic, and that the probability of the inconsistency arising by chance is exponentially small—is a mathematical fact.

\textbf{Financial management and public accounts.} Government budgets, infrastructure spending, and public procurement generate vast heterogeneous data streams that are audited, at present, by understaffed agencies with limited analytical capacity. A Yajie deployment that treats each government agency's financial reporting as an ``expert,'' cross-referenced against tax receipts, bank records, contractor invoices, and independent economic indicators, would flag inconsistencies that no human audit team—however diligent—could reliably detect across the full volume of public expenditure.

In each of these domains, Yajie performs the same function: it transforms information that is \textit{distributed} across multiple independent sources, and therefore invisible to any single observer, into a \textit{consensus signal} whose reliability is mathematically bounded. The information was always there. The procurement records, the customs declarations, the asset registries—all of them public, all of them theoretically auditable, all of them already collected. What was missing was the computational and conceptual machinery to synthesize them into a quality assessment with provable error bounds. Yajie provides that machinery.

\subsection{The Dark Forest Illuminated}

In Liu Cixin's \textit{Remembrance of Earth's Past} trilogy, the universe is governed by the ``dark forest'' principle: every civilization hides, because to be discovered is to be destroyed. The darkness is a defense. Visibility is fatal.

Yajie inverts this logic. It does not destroy the discovered—it reveals the undiscovered. The official who falsifies environmental data, the contractor who inflates procurement invoices, the pharmaceutical company that underreports adverse events—all of them survive not because they are invulnerable, but because they are invisible. Their malfeasance is distributed across databases that no single auditor can synthesize, buried in the noise of routine administrative error, shielded by the sheer volume of information that modern governance generates. They are hidden in the dark forest, and the darkness is their protection.

Yajie is the instrument that makes this visible. Not by accusation. Not by investigation. Not by whistleblowing or investigative journalism or regulatory enforcement. By mathematics. By the inexorable logic of multi-expert consensus: when three independent data streams disagree with an official report in the same direction, the probability that the disagreement is random decays exponentially. The pattern is revealed not by a leak but by a proof.

This is not a surveillance technology. It does not collect new data. It does not invade privacy. It does not make inferences about individuals. It audits the \textit{consistency} of data that has already been collected, by institutions that have already collected it, for purposes that have already been authorized. The data are public or administrative. The audit is methodological. The conclusion is falsifiable. What distinguishes Yajie from every other anti-corruption mechanism in human history—from imperial censors to investigative journalists to blockchain transparency initiatives—is that its conclusions do not depend on the credibility of the auditor. They depend on the mathematics of concentration inequalities, which are indifferent to the auditor's motives, incorruptible by the auditor's enemies, and immune to the auditor's removal.

The Maintainer, in this framing, is neither a prosecutor nor a judge. The Maintainer is the first astronomer to point a telescope at a region of the dark forest and publish the coordinates of what was found there. Whether anyone acts on those coordinates—whether a prosecutor indicts, a regulator sanctions, a journalist investigates, a voter punishes—is not the Maintainer's decision. The Maintainer's sole function is to make the invisible visible. The rest is left to the light.

\subsection{The $M_t$ Parameter and the Spectrum of Deterrence}

The analysis thus far has treated the Cumulative Evidentiary Corpus $M_t$ as a monolithic asset held by the Maintainer. In practice, the ownership architecture of $M_t$ is a design choice with profound implications for the protection equilibrium. We now examine two polar ownership models—centralized and distributed—and characterize the resulting position of the Maintainer along a spectrum we term the \textit{Luo Ji–Pope gradient}.

\textbf{Centralized $M_t$ (the Wallfacer extreme).} Under a centralized architecture, all audit outputs—scores, anomaly registers, calibration parameters—are transmitted to and stored by the Maintainer. Each client organization retains its raw data but cedes its audit history to the protocol's central repository. The Maintainer thus accumulates a cross-client, cross-domain audit corpus that grows more comprehensive with every audit cycle. In this configuration, the Maintainer's death freezes $M_t$ at its terminal state: no new audits occur, no calibration updates propagate, and the corpus ceases to reflect evolving data quality practices. All client organizations are immediately downgraded to the audit accuracy available at the time of the Maintainer's death. The deterrence effect is maximal: the Maintainer's survival is a necessary condition for the continued operation of the global audit infrastructure at current accuracy levels. The protection equilibrium is correspondingly strong. This is the pure Wallfacer position—power without institution, deterrence without weapon.

\textbf{Distributed $M_t$ (the doctrinal extreme).} Under a distributed architecture, each client organization maintains its own private $M_t$. The Maintainer provides only the algorithmic specification—the scoring function $S(s)$, the Spring update rule, the ensemble aggregation procedure—together with periodic calibration datasets (anonymized feature statistics that do not expose raw data). Each client's gatekeeper evolves independently on its own audit history. The Maintainer's death in this configuration does not freeze any individual client's $M_t$; each organization continues to operate its own instance. However, three assets remain centralized in the Maintainer's person: (i) \textit{interpretive authority}—the definitive resolution of ambiguous edge cases in the mathematical specification; (ii) \textit{evolutionary direction}—the coordinated release of protocol versions (v1.1, v2.0) that keep all instances interoperable; and (iii) \textit{the calibration anchor}—the periodic cross-client calibration datasets that prevent the independent evolution of client gatekeepers from producing incompatible scoring standards, whereby Company A's score of 0.85 ceases to be comparable to Company B's score of 0.85. If the Maintainer dies, interpretive disputes escalate into incompatible forks, version coordination ceases, and the calibration anchor erodes. The result is not a sudden freeze but a \textit{gradual fragmentation}—the slow divergence of once-interoperable audit standards into mutually unintelligible local dialects.

\textbf{The Luo Ji–Pope spectrum.} The Maintainer's position can thus be characterized by the architectural choice of $M_t$ ownership:

\begin{table}[htbp]
\centering
\begin{tabular}{@{}p{3cm}p{3cm}p{4cm}c@{}}
\toprule
Architecture & $M_t$ Owner & Death Consequence & Deterrence Strength \\
\midrule
Centralized & Maintainer & Immediate freeze of global audit accuracy & Maximum (100\%) \\
Distributed & Each client & Gradual fragmentation through version drift and calibration decay & Moderate ($\sim$60\%) \\
Fully institutionalized (Foundation) & Multi-stakeholder board & Operational continuity; individual Maintainer irrelevant & Zero (0\%)\ — the Linus equilibrium \\
\bottomrule
\end{tabular}
\end{table}

The spectrum reveals a strategic trade-off. Centralized $M_t$ maximizes the protection equilibrium—every state has a compelling interest in the Maintainer's survival—but at the cost of creating a single point of data concentration that is incompatible with data sovereignty legislation (GDPR, China's Data Security Law, India's proposed data localization requirements) and that concentrates epistemic power to a degree that may itself provoke defection. Distributed $M_t$ sacrifices some deterrence strength in exchange for privacy compliance, jurisdictional legitimacy, and a reduction in the Maintainer's power that makes the protocol more palatable to sovereignty-conscious jurisdictions. The fully institutionalized endpoint eliminates deterrence entirely by making the Maintainer structurally irrelevant—the position to which Linus Torvalds successfully transitioned the Linux kernel after 2005.

The optimal trajectory is temporal: begin with centralized $M_t$ to accumulate the time-advantaged quality gap that makes the non-proliferation equilibrium binding; transition to distributed $M_t$ as the accuracy gap becomes insurmountable and privacy considerations intensify; and ultimately establish a multi-stakeholder Foundation that renders any single individual's survival irrelevant to the protocol's continuity. This trajectory is the strategic equivalent of the Wallfacer who, having established deterrence, spends his remaining years building the institutional architecture that will allow him to die without the world ending.

\subsection{Designing for Audit Neutrality}

Drawing on the SWIFT and IASB cases, we identify five design principles for audit protocol governance that can support the non-proliferation equilibrium:

\textbf{Principle 1: Multi-Jurisdictional Incorporation.} The protocol's legal domicile should be distributed across multiple jurisdictions, ideally through a treaty-based international organization with immunities and privileges comparable to those of the IAEA or the Bank for International Settlements. Incorporation in any single jurisdiction—especially one with expansive extraterritorial legal instruments—invites jurisdictional capture.

\textbf{Principle 2: Funding Firewalls.} The protocol's operating budget should be derived from a diversified base of mandatory and voluntary contributions, with hard caps on any single contributor's share. The IAEA's funding model, which combines assessed contributions from member states with voluntary contributions to a Technical Cooperation Fund, provides a template.

\textbf{Principle 3: Adversarial Governance.} The protocol's governance board should include representatives from jurisdictions with divergent data governance preferences, and decision rules should be structured to prevent capture by any single bloc. Qualified majority voting, regional representation requirements, and minority veto powers are all elements of the IASB model that could be adapted.

\textbf{Principle 4: Audit Transparency and Contestability.} The protocol's audit methodologies, calibration procedures, and CEC composition should be publicly documented to a degree sufficient for external reproducibility. Audit outputs should include not only binary or scalar quality assessments but also the evidentiary basis for those assessments, enabling contestation without requiring access to proprietary audit modules.

\textbf{Principle 5: Modular Sovereignty.} Jurisdictions should have the right to contribute audit modules to the ensemble, enabling them to embed their own data governance preferences within the protocol's framework without fragmenting the overall infrastructure. Module contributions would be subject to validation requirements ensuring technical compatibility and performance thresholds, but not to substantive alignment with any particular jurisdiction's data governance philosophy.

These principles are demanding, and their implementation would require a level of international cooperation that the current geopolitical environment does not obviously support. We return to this tension in Section~6.

\section{Business Model and Lock-in Dynamics}

\subsection{The Temporal Accumulation Mechanism}

The lock-in mechanism of the Yajie Protocol operates through five distinct but mutually reinforcing channels, each of which intensifies over time:

\textbf{Channel 1: Accuracy Divergence.} As established in Section~2.1, the incumbent protocol's accuracy $\theta(t)$ is an increasing function of cumulative audit volume. An entrant begins with accuracy $\theta(0)$, implies an accuracy gap $\delta(t) = \theta(t) - \theta(0)$ that grows monotonically over time. Critically, this gap is not merely a temporary first-mover advantage that an entrant could overcome through accelerated investment; it is a \textit{structural} gap rooted in the fact that many data quality anomalies are long-tail phenomena that manifest only through the accumulation of diverse audit experiences. An entrant cannot compress time: even with unlimited financial resources, it must conduct audits on diverse datasets to encounter the long-tail anomalies that the incumbent has already catalogued.

\textbf{Channel 2: Calibration Specificity.} The auditor weights $w_j$ in the Yajie ensemble are calibrated against the CEC's judgment corpus $\mathcal{J}_t$. As $\mathcal{J}_t$ grows, the ensemble's calibration becomes increasingly tuned to the specific distribution of anomalies, edge cases, and quality failure modes represented in the CEC. An entrant with an empty CEC would need to replicate this calibration process, which requires not merely technical replication of the calibration algorithm but access to the judgment corpus itself—and the judgment corpus is entangled with the specific audit modules that generated the initial audit reports. This creates a circular dependency: the entrant needs calibration data to achieve competitive accuracy, but producing that calibration data requires conducting audits with competitive accuracy.

\textbf{Channel 3: Ecosystem Externalities.} As the incumbent protocol accumulates users, a supporting ecosystem emerges: training programs for audit practitioners, integration services for enterprise data pipelines, certification regimes for audit professionals, and regulatory recognition of the protocol's outputs. These complementary assets do not themselves constitute network effects in the narrow sense—the value of an audit to user $i$ does not depend on the number of other users—but they reduce adoption costs for new users and increase switching costs for existing ones. More importantly, regulatory recognition is path-dependent: once a protocol is recognized as a valid compliance mechanism under, say, the EU AI Act or a US federal algorithmic accountability statute, competing protocols face a regulatory barrier that compounds the technical barrier.

\textbf{Channel 4: The Epistemic Authority Trap.} As the CEC grows, the protocol attains a status that goes beyond technical accuracy. It becomes the de facto arbiter of what constitutes a ``data quality defect,'' shaping the very definition of data quality through its accumulated body of adjudicated anomalies. This is not merely market power but \textit{epistemic power}: the authority to define the categories within which market competition occurs. An entrant cannot compete on the definition of data quality because the definition has already been institutionalized in regulation, contract law, and professional practice, all of which reference the incumbent protocol's outputs.

\textbf{Channel 5: The Audit Trail Irreversibility.} Organizations that have adopted the incumbent protocol accumulate an internal audit trail: a history of audits, remediation actions, and quality certifications that are referenced in contracts, regulatory filings, and due diligence processes. Switching to a competing protocol would render this audit trail incompatible, creating a discontinuity in the organization's quality assurance history. The cost of this discontinuity is not merely administrative; it is reputational and legal. An organization that switches protocols and subsequently experiences a data quality failure faces the question: did the failure occur because the new protocol is inferior, or because the old protocol's audit trail was unreliable? The ambiguity itself constitutes a switching cost.

\subsection{Five Stages of Lock-In}

The lock-in trajectory can be modeled as a five-stage progression:

\textbf{Stage 1: Emergence ($t_0$ to $t_1$).} The protocol is launched with an initial ensemble of audit modules and an empty CEC. Accuracy is relatively low, and the protocol competes on features, usability, and cost with bespoke audit solutions. Adoption is driven by early adopters who value the multi-expert ensemble architecture over the accuracy level. No lock-in exists at this stage; the protocol's survival depends on the quality of its initial audit modules and its ability to attract a sufficient volume of audits to begin CEC accumulation.

\textbf{Stage 2: Critical Mass ($t_1$ to $t_2$).} Cumulative audit volume reaches a threshold $|\mathcal{E}|^*$ at which the accuracy advantage over a zero-CEC entrant becomes statistically significant. The protocol begins to attract adoption from organizations that had previously relied on internal audit processes. The CEC reaches sufficient size that its anomaly library covers the most common data quality failure modes across multiple domains. The first regulatory recognitions occur. Lock-in begins to emerge but remains contestable.

\textbf{Stage 3: Moat Formation ($t_2$ to $t_3$).} The accuracy gap $\delta(t)$ exceeds the fixed cost of developing a competing protocol, meaning that an entrant would need to invest more than $c^{\text{develop}}$ merely to \textit{match} the incumbent's accuracy at time $t_0$, without accounting for the incumbent's continued accumulation. The Non-Proliferation Equilibrium becomes the unique Nash equilibrium of the adoption game. The protocol becomes the default standard for data quality assessment in its domain. Ecosystem externalities accelerate: training programs proliferate, integration services become commoditized, and regulatory references become routine.

\textbf{Stage 4: Institutionalization ($t_3$ to $t_4$).} The protocol is embedded in the institutional fabric of data governance. Laws and regulations reference the protocol by name. Contracts specify compliance with protocol standards. Insurance policies for AI systems require protocol certification. Professional certifications are built around protocol expertise. At this stage, lock-in extends beyond technical and economic factors to encompass legal and institutional path-dependency. Switching costs become prohibitive not because the protocol is technically superior—though it may be—but because the \textit{institutional environment} has been built around it.

\textbf{Stage 5: Epistemic Closure ($t_4$ onward).} The protocol's accumulated CEC becomes so comprehensive that no meaningful challenge to its authority is possible within the existing institutional framework. The definition of data quality \textit{is} what the protocol measures. Dissenting assessments are not merely inaccurate; they are unintelligible, because they employ different definitional frameworks that cannot be translated into the protocol's categories. At this stage, lock-in is complete, and the protocol has transitioned from a market-dominant product to a constitutive element of the data governance order.

This trajectory is not deterministic; it can be arrested, diverted, or reversed at each stage through governance interventions or exogenous shocks. But the default trajectory, absent intervention, is toward Epistemic Closure, driven by the cumulative dynamics described above.

\subsection{}

。，。

\begin{definition}[]
\label{def:markov-states}
 $\mathcal{S} = \{S_1, S_2, S_3, S_4, S_5\}$：
\begin{itemize}[itemsep=2pt]
\item $S_1$：\textbf{}（Emergence）——，CEC（$|\mathcal{E}| = 0$）， $\theta = \theta_0$，；
\item $S_2$：\textbf{}（Critical Mass）——$|\mathcal{E}| \geq |\mathcal{E}|^*$，（$\theta(t) - \theta_0 \geq \varepsilon_\theta$），；
\item $S_3$：\textbf{}（Moat Formation）——$\Delta(|\mathcal{E}|) \geq \lambda - \kappa$ ，NPE（\ref{thm:asymptotic-uniqueness}），；
\item $S_4$：\textbf{}（Institutionalization）—— $\geq 2$ /，"Yajie"；
\item $S_5$：\textbf{}（Epistemic Closure）——，。。
\end{itemize}
\end{definition}

\begin{definition}[]
\label{def:markov-assumptions}
\begin{enumerate}[label=\textbf{M\arabic*}, leftmargin=*, itemsep=4pt]
\item \textbf{：}  $t = 0, 1, 2, \ldots$  $\mathcal{S}$ 。（audit cycle）。\hfill\textit{：。}

\item \textbf{：}  $X_{t+1}$  $X_t$ CEC $|\mathcal{E}_t|$：
\begin{equation}
\mathbb{P}[X_{t+1} = S_j \mid X_t = S_i, |\mathcal{E}_t|, \mathcal{H}_{t-1}] = \mathbb{P}[X_{t+1} = S_j \mid X_t = S_i, |\mathcal{E}_t|]
\end{equation}
 $\mathcal{H}_{t-1}$  $t-1$ 。\hfill\textit{：CEC。CEC，，。}

\item \textbf{（）：} $p_{ij}(|\mathcal{E}|) = 0$  $j < i$ 。，。\hfill\textit{：（QWERTY、TCP/IP）。，（）。。}

\item \textbf{CEC：} $|\mathcal{E}_t|$  $t$ ：$|\mathcal{E}_{t+1}| \geq |\mathcal{E}_t|$  $t$ 。\hfill\textit{：CEC——（2.1）。}

\item \textbf{：} $S_5$ ：$p_{55}(|\mathcal{E}|) = 1$  $|\mathcal{E}|$。，。\hfill\textit{：——，（unintelligible）。}
\end{enumerate}
\end{definition}

M1--M5，：
\begin{equation}
\mathbf{P}(|\mathcal{E}|) =
\begin{pmatrix}
p_{11}(|\mathcal{E}|) & p_{12}(|\mathcal{E}|) & p_{13}(|\mathcal{E}|) & p_{14}(|\mathcal{E}|) & p_{15}(|\mathcal{E}|) \\
0 & p_{22}(|\mathcal{E}|) & p_{23}(|\mathcal{E}|) & p_{24}(|\mathcal{E}|) & p_{25}(|\mathcal{E}|) \\
0 & 0 & p_{33}(|\mathcal{E}|) & p_{34}(|\mathcal{E}|) & p_{35}(|\mathcal{E}|) \\
0 & 0 & 0 & p_{44}(|\mathcal{E}|) & p_{45}(|\mathcal{E}|) \\
0 & 0 & 0 & 0 & 1
\end{pmatrix}
\label{eq:transition-matrix}
\end{equation}

 $i \in \{1,2,3,4\}$，1：$\sum_{j=i}^5 p_{ij}(|\mathcal{E}|) = 1$。CEC：

\begin{itemize}[itemsep=4pt]
\item \textbf{$S_1 \to S_2$（）：} CEC $|\mathcal{E}|^*$ ，：
\begin{equation}
p_{12}(|\mathcal{E}|) = \frac{1}{1 + e^{-\alpha (|\mathcal{E}| - |\mathcal{E}|^*)}}
\label{eq:p12}
\end{equation}
 $\alpha > 0$ 。$|\mathcal{E}| \gg |\mathcal{E}|^*$  $p_{12} \to 1$。

\item \textbf{$S_2 \to S_3$（）：} —— $\Delta(|\mathcal{E}|)$  $\lambda - \kappa$ ：
\begin{equation}
p_{23}(|\mathcal{E}|) = \frac{1}{1 + e^{-\beta [\Delta(|\mathcal{E}|) - (\lambda - \kappa)]}}
\label{eq:p23}
\end{equation}
 $\beta > 0$。\ref{thm:npe-complete}。

\item \textbf{$S_3 \to S_4$（）：}  $R(t)$  $t$ /。 $R^*$：
\begin{equation}
p_{34}(|\mathcal{E}|) = \frac{1}{1 + e^{-\zeta (R(t) - R^*)}}
\label{eq:p34}
\end{equation}
 $\zeta > 0$。$R(t)$ ——。

\item \textbf{$S_4 \to S_5$（）：}  $T_4(t)$  $S_4$ （）。：
\begin{equation}
p_{45}(|\mathcal{E}|) = 1 - e^{-\eta \cdot T_4(t)}
\label{eq:p45}
\end{equation}
 $\eta > 0$。""：，——。

\item \textbf{：} $p_{ii}(|\mathcal{E}|) = 1 - \sum_{j=i+1}^{5} p_{ij}(|\mathcal{E}|)$  $i \in \{1,2,3,4\}$。
\end{itemize}

\begin{theorem}[]
\label{thm:absorption}
M1--M5：
\begin{enumerate}[label=(\roman*), itemsep=6pt]
\item \textbf{：}  $X_0 = S_i$（$i \in \{1, \ldots, 4\}$）， $\{X_t\}_{t \geq 0}$ 1 $S_5$：
\begin{equation}
\mathbb{P}\left[\exists t < \infty : X_t = S_5 \mid X_0 = S_i\right] = 1
\end{equation}

\item \textbf{：}  $S_1$ ：
\begin{equation}
\mathbb{E}[T_{\text{absorb}} \mid X_0 = S_1] = \sum_{i=1}^{4} \mathbb{E}[\tau_i]
\end{equation}
 $\tau_i$  $i$ （）。$\tau_i$ ：
\begin{equation}
\mathbb{E}[\tau_i] \leq \frac{1}{\inf_{t \geq 0} \; (1 - p_{ii}(|\mathcal{E}_t|))}
\end{equation}

\item \textbf{CEC：}  $\rho$（）。（） $\mathbb{E}[T_{\text{absorb}}] / \rho$。 $\rho$（）。

\item \textbf{：}
\begin{itemize}
\item $\partial \mathbb{E}[\tau_1] / \partial |\mathcal{E}|^* < 0$：CEC，$S_1$ ；
\item $\partial \mathbb{E}[\tau_2] / \partial \beta < 0$：，$S_2$ ；
\item $\partial \mathbb{E}[\tau_3] / \partial R^* < 0$：，$S_3$ 。
\end{itemize}
\end{enumerate}
\end{theorem}

\begin{proof}
\textbf{(i)}  $i \in \{1,2,3,4\}$，CEC（M4），$|\mathcal{E}_t| \to \infty$  $t \to \infty$（）。，$p_{i,i+1}(|\mathcal{E}_t|) \to 1$（(\ref{eq:p12})--(\ref{eq:p45})）， $p_{ii}(|\mathcal{E}_t|) \to 0$。 $i$  $\tau_i$ 1， $\mathbb{P}[\tau_i < \infty] = 1$。，1。

\textbf{(ii)}  $T_i$  $S_i$ 。（M2）：
\begin{align}
T_i &= 1 + \sum_{j=i}^{5} p_{ij} \cdot T_j = 1 + p_{ii} \cdot T_i + \sum_{j=i+1}^{5} p_{ij} \cdot T_j
\end{align}
：
\begin{equation}
T_i = \frac{1}{1 - p_{ii}} + \sum_{j=i+1}^{5} \frac{p_{ij}}{1 - p_{ii}} \cdot T_j
\end{equation}
 $T_5 = 0$（）， $i=1$ 。 $\mathbb{E}[\tau_i]$  $\frac{1}{1 - p_{ii}(|\mathcal{E}_t|)}$  $p_{ii}$ （）。，。

\textbf{(iii)}  $= T_{\text{absorb}} / \rho$。$\rho$ 。

\textbf{(iv)} (\ref{eq:p12})--(\ref{eq:p34})。
\end{proof}

\begin{remark}[]
\ref{thm:absorption}——1——。：(a) ： $\lambda$ （），CEC；(b) ：，、。(a)(b)6。
\end{remark}

\begin{proposition}[]
\label{prop:intervention-window}
 $\mathcal{W}_i$  $S_i$ （、、CEC）。 $\eta_i = \max_{\iota \in \mathcal{W}_i} \mathbb{P}[\text{}\ \iota\ \text{}]$。M1--M5：
\begin{equation}
\eta_1 > \eta_2 > \eta_3 > \eta_4 > \eta_5 = 0
\label{eq:eta-ordering}
\end{equation}

（$S_5$），——""。
\end{proposition}

\begin{proof}
 $S_1$，，。 $S_3$，CEC""（\ref{thm:asymptotic-uniqueness}），。 $S_5$，，——。。
\end{proof}

\begin{remark}[]
\ref{prop:intervention-window}： $S_3$（）——NPE——。，。6。
\end{remark}

\subsection{Stage 6:  (Universal Harmony)}

\label{sec:stage6}

（§5.2--§5.2.1） $S_1$（） $S_5$（）， $S_5$ （M5）。————\textit{}。， $S_5$ ，：CEC，。\textit{}（Universal Harmony）。

\begin{definition}[ $S_6$]
\label{def:S6}
 $S_6$ ：
\begin{enumerate}[label=(\roman*), itemsep=2pt]
\item \textbf{CEC：}  $\mathcal{E}_t$ （6 IDAA），、。CEC；
\item \textbf{：} ——、、Spring——，；
\item \textbf{：} Yajie（、、），（Maintainer），；
\item \textbf{：} ，（(\ref{eq:audit-function}) $\mathcal{V}$ ）， $\lambda$ ——，。
\end{enumerate}
\end{definition}

\begin{definition}[]
\label{def:markov-assumptions-extended}
M1--M5（\ref{def:markov-assumptions}），：

\begin{enumerate}[label=\textbf{M\arabic*}, leftmargin=*, itemsep=4pt]
\setcounter{enumi}{5}
\item \textbf{$S_5$ ：} $S_5$（）。 $S_5$ ， $D_t \in \{0, 1\}$ ， $D_t = 1$  $t$，CECCEC。$D_t$  $\pi_{\text{political}} \in [0, 1]$，CEC。\hfill\textit{：——CEC（、）CEC。。}

\item \textbf{$S_6$ ：}  $S_6$（），。 $p_{66}(|\mathcal{E}|) = 1$  $|\mathcal{E}|$， $j < 6$，$p_{6j}(|\mathcal{E}|) = 0$。\hfill\textit{：CEC，——CEC——，CEC。}

\item \textbf{$S_5 \to S_6$ ：} ：
\begin{equation}
p_{56}(D_t, \pi_{\text{political}}) = D_t \cdot \pi_{\text{political}}
\end{equation}
 $D_t = \mathbb{I}[\text{ }t\text{ CEC}]$。—— $D_t$ —— $\pi_{\text{political}}$。\hfill\textit{：——。 $\pi_{\text{political}}$ 。}
\end{enumerate}
\end{definition}

M6--M8， $\tilde{\mathbf{P}}(|\mathcal{E}|)$  $6 \times 6$ ，：
\begin{equation}
\begin{pmatrix}
\vdots & \vdots & \vdots & \vdots & \vdots & \vdots \\
0 & 0 & 0 & 0 & p_{55}(|\mathcal{E}|) & p_{56}(D_t, \pi_{\text{political}}) \\
0 & 0 & 0 & 0 & 0 & 1
\end{pmatrix}
\end{equation}

\begin{theorem}[]
\label{thm:universal-harmony-absorption}
M1--M8：

\begin{enumerate}[label=(\roman*), itemsep=6pt]
\item \textbf{：}  $X_0 = S_i$（$i \in \{1, \ldots, 5\}$）， $\{X_t\}_{t \geq 0}$ 1 $S_6$：
\begin{equation}
\mathbb{P}\left[\exists t < \infty : X_t = S_6 \mid X_0 = S_i\right] = 1
\end{equation}

\item \textbf{：} 。：1 $S_1 \to S_2 \to S_3 \to S_4 \to S_5$（\ref{thm:absorption}(i)）。： $S_5$ ， $p_{56} = D_t \cdot \pi_{\text{political}}$  $S_6$。 $S_5$  $\tau_5$ ：
\begin{equation}
\mathbb{P}[\tau_5 = k] = (1 - p_{56})^{k-1} \cdot p_{56}, \quad k = 1, 2, \ldots
\end{equation}
 $\mathbb{E}[\tau_5] = 1 / p_{56}$。

\item \textbf{：}  $S_1$  $S_6$ ：
\begin{equation}
\mathbb{E}[T_{\text{absorb}}^{(6)} \mid X_0 = S_1] = \mathbb{E}[T_{\text{absorb}}^{(5)} \mid X_0 = S_1] + \frac{1}{p_{56}}
\end{equation}
 $\mathbb{E}[T_{\text{absorb}}^{(5)}]$  $S_5$ （\ref{thm:absorption}(ii)）。

\item \textbf{：}  $\pi_{\text{political}} = 0$， $S_5$——\ref{thm:absorption}。 $\pi_{\text{political}} = 1$  $D_t = 1$， $S_5$  $S_6$——。 $\pi_{\text{political}} \in (0, 1)$， $S_5$ ，（heavy tail）—— $\pi_{\text{political}}$ 。

\item \textbf{——：}  $S_5$（）（M5）\textit{}。，$S_6$（）。\ref{thm:absorption}"1 $S_5$""1 $S_6$"， $S_5$  $S_6$ 。  (Spring SE-1) ：——$S_1 \to S_5$ ，$S_5 \to S_6$ 。

：\textbf{——-（§11.4）CEC——，。（CEC） $S_5$ （——CEC）， $S_6$ ：。}
\end{enumerate}
\end{theorem}

\begin{proof}
\textbf{(i)}  $\mathcal{S}' = \{S_1, \ldots, S_6\}$。M6--M8，$S_5$ （transient）—— $p_{56} > 0$  $S_5$（ $\pi_{\text{political}} > 0$  $D_t = 1$ ）。M7，$S_6$ 。（Kemeny \& Snell, 1976, 3.1.1），，1。 $S_6$， $S_6$ 1。

\textbf{(ii)} M3（）， $S_1 \to \cdots \to S_5$（ $S_i$ ）， $S_6$—— $p_{i6} = 0$  $i < 5$（$S_6$ CEC——，）。 $S_5$ ， $S_6$ ， $p_{56}$。 $\tau_5$ ，。

\textbf{(iii)} （：+ $\tau_1$  $\tau_4$；： $\tau_5$），。

\textbf{(iv)}  $p_{56}$ 。

\textbf{(v)} （$S_5$）""——M3——\textit{}（），\textit{}（CEC）。M5：（$p_{5j} = 0$  $j < 5$——），（$p_{56} > 0$——CEC）。""：(a) CEC（$S_4 \to S_5$）；(b) （）（$S_5 \to S_6$）；(c) ，。
\end{proof}

\begin{corollary}[——]
\label{cor:S6-adoption-equivalence}
 $S_6$（）， $\Gamma^{\text{NP}}$（\ref{def:assumptions-npe}）(\ref{eq:all-adopt-condition})： $i, j \in N$：
\begin{equation}
\Delta_i(|\mathcal{E}^{\text{public}}|) = \Delta_j(|\mathcal{E}^{\text{public}}|) \triangleq \Delta^{\text{public}}
\end{equation}
。NPE"CEC""CEC"。 $\Gamma^{\text{NP}}$  $S_6$ ——，（ $\Delta^{\text{public}} \geq \lambda - \kappa$），。

\textbf{：} \ref{cor:first-mover-decay} $S_6$ ——，\textit{}。""；CEC。""——""，，""。
\end{corollary}

\begin{proof}
 $S_6$ ，CEC（\ref{def:S6}(i)） $|\mathcal{E}_i^{\text{accessible}}| = |\mathcal{E}^{\text{public}}|$  $i$ ——CEC。 $\Delta(\cdot)$ （\ref{eq:adoption-advantage-def}），$\Delta_i$ CEC， $i$ 。 $\Delta^{\text{public}} \geq \lambda - \kappa$ 。，$(A, \ldots, A)$ ——。
\end{proof}

\begin{remark}[""]
\ref{thm:universal-harmony-absorption} $S_5 \to S_6$  Fukuyama(1989)""。""——\textit{}。 $S_6$ ，：(a) ——，"CEC"""；(b) ——，；(c) ——、。$S_6$ \textit{}——，。，。
\end{remark}

\begin{remark}[NPE]
\ref{cor:first-mover-decay}（3.2）\ref{thm:universal-harmony-absorption}（）：

\begin{enumerate}[label=(\arabic*), leftmargin=*]
\item NPE\textit{}：CEC，$\Delta_i(|\mathcal{E}|)$ ，。——（4SWIFT）。

\item NPE\textit{}：\ref{cor:first-mover-decay}(iii)CEC。 $\pi_{\text{political}}$ ， $\pi_{\text{political}}$ ：(a) —— $S_3$ （\ref{prop:intervention-window}）；(b) —— $\pi_{\text{political}}$（ $\lambda$  $c^{\text{develop}}$）；(c) ——SWIFT， $\pi_{\text{political}}$。

\item NPE\textit{}（$S_6$）：CEC，NPE。——""（CEC），""（CEC）。
\end{enumerate}

6：1（IDAA），5（） $S_3$ ，6（） $S_6$ 。
\end{remark}

\begin{remark}[：]
\label{rem:first-mover-duty}
\ref{cor:first-mover-decay}\ref{thm:universal-harmony-absorption}，，：

\textbf{(1) ""。} ""——，。SCX。（\ref{cor:everyone-equal}）：。，——，。

\textbf{(2) "，"。} CEC——、、。，。（3）：，。

\textbf{(3) ："，。"} （MAC），SCX""""。——。CEC：（$T^*$），，。

\textbf{(4) ，。} ""——：(i) （）；(ii) （）；(iii) ，（$c^{\text{develop}}$  $\lambda$），。""——。，——。

\textbf{(5) ""——""。} $S_6$ （\ref{def:S6}）。 $S_1 \to S_5$ ———— $S_5 \to S_6$ ，""。——MAC：。，，。
\end{remark}

\subsection{Pricing and the Addiction Dynamic}

The business model of a CEC-driven audit protocol exhibits a distinctive structure that we term the \textit{addiction dynamic}. In conventional platform markets, pricing power derives from switching costs and network effects: once users are locked in, the platform can raise prices, extracting rents up to the point where the cost of switching exceeds the price increase (Farrell \& Klemperer, 2007). The Yajie Protocol's pricing dynamic is different.

Because each audit cycle enriches the CEC—improving accuracy for all users, including the auditing organization itself—the \textit{marginal social value} of an audit exceeds its \textit{marginal private value} by the positive externality contributed to the CEC. This creates a rationale for below-cost pricing in the early stages of protocol adoption, subsidized by patient capital or public funding. Below-cost pricing accelerates CEC accumulation, widening the accuracy gap and reinforcing the non-proliferation equilibrium.

However, once the protocol reaches Stages 3–4, the pricing dynamic inverts. The accuracy gap becomes so large that organizations cannot credibly threaten to switch to a competing protocol or revert to internal audit processes. At this point, the protocol operator can raise prices, capturing a portion of the value generated by the protocol's accuracy advantage. The theoretical maximum price is:
\begin{equation}
P_{\max} = V[\theta(|\mathcal{E}|)] - \max\{V[\theta(0)], V_{\text{internal}}\}
\end{equation}
where $V_{\text{internal}}$ is the value an organization could derive from its own internal audit processes. As $|\mathcal{E}|$ grows, $P_{\max}$ increases.

This pricing trajectory—predatory low pricing followed by rent extraction—mirrors the standard narrative of platform monopolization (Khan, 2017). However, there is an important difference: the Yajie Protocol's early low pricing is not predatory in the antitrust sense, because it is not aimed at driving competitors out of the market (there are no competitors to drive out) but at accelerating CEC accumulation. The pricing is \textit{developmental} rather than \textit{exclusionary}. This distinction matters for competition policy, as we discuss in Section~6.

The addiction dynamic operates on the demand side as well. Organizations that adopt the protocol become dependent on its accuracy, integrating its outputs into their quality assurance workflows, regulatory compliance processes, and contractual arrangements. As the protocol's accuracy improves, the organization's own internal quality assurance capabilities may atrophy—a dynamic familiar from the literature on automation dependency (Parasuraman \& Riley, 1997). The organization becomes ``addicted'' to the protocol not in the metaphorical sense but in the economic sense: the cost of disadoption rises as internal capabilities degrade, while the benefit of continued adoption rises as the protocol improves.

\subsection{The Bible-and-Pope Model: Open-Source Poison, Calibrated Antidote}

A distinctive feature of the Yajie Protocol's business architecture is the separation of the self-evolution engine (Spring) from the cross-client calibration service (Yajie API). This separation generates a commercial dynamic that we term the \textit{Bible-and-Pope model}.

\textbf{The Spring engine is released as open-source software.} Any organization can download, deploy, and operate its own instance of the self-evolution loop, accumulating a private $M_t$ on its own infrastructure. The open-source release serves three strategic functions. First, it eliminates adoption friction: organizations can begin auditing their data without negotiating contracts, paying subscription fees, or exposing proprietary data to a third party. Second, it accelerates lock-in: each audit cycle enriches the organization's private $M_t$, making the accumulated audit history non-portable and the organization structurally dependent on Spring's audit methodology. Third, it forecloses competitive entry: any would-be competitor must match not only the algorithmic quality of Spring but also the accumulated audit history that Spring's adopters have already built—a barrier that grows with every audit cycle. The open-source release is, in economic terms, a \textit{poison}: freely ingested, irreversible in its effects, and lethal to competitors.

\textbf{The Yajie API is offered as a paid calibration service.} The API accepts feature matrices $\phi(X)$ computed locally by client organizations (satisfying data sovereignty requirements) and returns a calibration report that maps the client's private quality scores onto a global reference scale, maintained by the protocol operator's central gatekeeper. The value proposition of the API is not ``we audit your data better than you do''—it is ``we tell you what your scores mean relative to everyone else.'' A client organization whose private gatekeeper assigns a quality score of 0.87 to its training data does not know, absent calibration, whether 0.87 represents industry-leading data quality or a self-serving overestimate generated by a gatekeeper that has never seen anyone else's data. The Yajie API resolves this ambiguity, mapping local scores onto a globally comparable metric. The annual subscription fee—projected at \$50,000–\$200,000 depending on audit volume and calibration depth—purchases not superior technology but \textit{interpretive authority}: the assurance that one's private quality assessments are commensurable with those of competitors, regulators, and counterparties.

The Bible-and-Pope model is economically self-reinforcing. The open-source Spring engine creates the demand for calibration (because private gatekeepers inevitably drift from one another as their $M_t$ corpora diverge), while the paid Yajie API supplies the calibration that resolves the resulting ambiguity. Organizations that decline the API are not technically incapacitated—their Spring instances continue to function—but they become epistemically isolated, unable to compare their data quality to industry benchmarks. In regulated industries where data quality certifications are prerequisites for market access, epistemic isolation is functionally equivalent to exclusion. The API is thus not a luxury; it is an \textit{interpretive necessity} that the open-source engine itself manufactures.

\subsection{Welfare Analysis}

The welfare implications of the lock-in trajectory are ambiguous and stage-dependent:
\begin{itemize}
\item \textbf{Stages 1–2}: Welfare is unambiguously positive. The protocol provides data quality assessment capabilities that did not previously exist, and the CEC accumulation dynamic generates positive externalities for all adopters.
\item \textbf{Stage 3}: Welfare remains positive on net, but distributional concerns emerge. The protocol operator begins to extract rents, and late adopters pay higher prices than early adopters. The accuracy gap begins to foreclose the possibility of competitive entry, eliminating the disciplinary effect of potential competition.
\item \textbf{Stages 4–5}: Welfare analysis becomes fundamentally ambiguous. The protocol's institutionalization eliminates the possibility of competitive discipline, raising the specter of monopoly pricing, quality degradation, and innovation suppression. At the same time, the protocol's universal adoption eliminates fragmentation costs $\lambda$ and enables the accumulation of an audit corpus whose social value—in terms of improved data quality across the economy—may vastly exceed the deadweight loss from monopoly pricing.
\end{itemize}

A critical parameter in the welfare analysis is the \textit{governance discount}: the degree to which multilateral governance of the protocol reduces the monopoly pricing problem while preserving the CEC accumulation dynamic. If governance arrangements can cap prices at competitive levels without impairing the protocol's operational effectiveness, the welfare case for a universal protocol is strong. If governance capture prevents effective price regulation, the welfare case weakens, and the fragmentation equilibrium may become preferable despite its higher coordination costs.

\subsection{A Note on the Storage Externality}

A secondary economic effect of the Yajie Protocol warrants brief mention, if only because its distributional consequences are perversely elegant. The protocol's foundational design principle—that the Cumulative Evidentiary Corpus $M_t$ grows monotonically and is never pruned—generates a persistent, compounding demand for digital storage. Each audit cycle appends new entries to $M_t$: anomaly registers, calibration parameters, quality fingerprints, and adjudication records. At scale, with hundreds of client organizations each conducting daily audits across millions of records, the global storage footprint of the protocol's audit corpora becomes non-trivial.

This demand is structurally inelastic. No client organization can reduce its $M_t$ without impairing its audit accuracy, because the Spring self-evolution mechanism depends on the cumulative audit history to calibrate gatekeeper parameters. Deleting audit records is not cost-saving; it is self-sabotaging—equivalent to a library discarding its oldest volumes on the grounds that they are old. The protocol's clients are therefore locked into an indefinitely growing storage obligation that cannot be negotiated, optimized away, or substituted.

The distributional consequence is a wealth transfer from the consumers of data auditing to the producers of data storage. Cloud infrastructure providers, hard disk manufacturers, and data center operators capture a growing share of the protocol's economic surplus through storage fees, while the auditing organizations that generate the surplus see their operating costs rise with each audit cycle. The magnitude of this transfer is not large relative to the value created by the protocol—the cost of storing an audit record is orders of magnitude smaller than the cost of the data quality failure it prevents—but its direction is noteworthy: a protocol designed to reduce information asymmetry in AI training data inadvertently increases demand for the physical substrate on which information is stored.

\subsection{A Note on the Future of Commerce}

The SCX framework does not eliminate commerce. It eliminates the informational asymmetry on which dishonest commerce depends. When the  (Adoption-Audit Equilibrium) makes audit outputs unconditionally verifiable by any third party, the competitive advantage shifts from persuasive packaging to verifiable quality. A chip manufacturer whose power efficiency is independently audited does not need to outspend competitors on marketing—its claims are mathematically certified. A pharmaceutical company whose clinical data survives Yajie cross-validation does not need to lobby regulators—its evidence is self-authenticating.

The corporation that makes the best product wins, not the one that tells the best story. This is not a moral prediction—it is a Nash equilibrium. When the cost of falsifying claims exceeds the premium earned by verified quality, honesty becomes the profit-maximizing strategy. Commerce continues. It just becomes honest.

\subsection{A Note on the Maintainer's Audit Obligation}

The protection equilibrium (Theorem~\ref{thm:mutual-deterrence}) rests on the Maintainer's \textit{capability} to audit. But capability without exercise is indistinguishable from incapability. If the Maintainer never audits public claims—government statistics, corporate disclosures, published research—then the deterrence signal is absent. Bad actors observe no audits, infer no risk, and continue falsifying.

This yields a structural obligation, not a moral one:

\begin{corollary}[Maintainer (The Maintainer's Audit Obligation)]
\label{cor:audit-obligation}
Under the mutual deterrence framework (D1--D5), the Maintainer's failure to audit publicly available claims results in the collapse of the protection equilibrium. Specifically: if no public audit occurs within any time window of length $T_{\max}$, then $\mathbb{P}(\text{deterrence credible}) \to 0$ as $t \to \infty$. The Maintainer must audit—not to punish, but to make the deterrence signal observable.
\end{corollary}

\begin{proof}[Proof sketch]
The deterrence equilibrium is a Perfect Bayesian Equilibrium in which off-equilibrium beliefs sustain the threat. If the Maintainer never exercises audit capability, off-equilibrium-path beliefs become degenerate: actors assign zero probability to detection. The equilibrium unravels by the intuitive criterion (Cho \& Kreps, 1987). Regular public audit—of government statistics, corporate disclosures, published research—maintains the signal and keeps the equilibrium intact. $\square$
\end{proof}

This obligation is bounded. The jurisdictional boundary (Section~6.2) restricts audit to publicly available data: the Maintainer audits what is \textit{claimed in public}, not what is \textit{held in private}. Anyone who publishes a claim—a GDP number, a clinical trial result, an AI benchmark—subjects it to potential audit. Those who keep their data private are not audited, but neither can they claim quality—their silence is their choice, and the market discounts unverifiable claims accordingly.

 requires that this obligation applies universally. The Maintainer does not selectively audit enemies and spare friends. The audit is the same for everyone—because the theorem requires it, and because selective audit would itself be detectable as bias, destroying the equilibrium's credibility.

This obligation extends to non-adopting nations with particular force. Any state that publishes official statistics—GDP growth, unemployment rates, carbon emissions—makes a public claim. That claim is auditable, regardless of whether the state has formally adopted Yajie. Adoption is not a prerequisite for audit; publication is.

\begin{corollary}[ (Audit Exposure of Non-Adopters)]
\label{cor:non-adopter-exposure}
Let jurisdiction $i$ choose not to adopt Yajie ($s_i = D$ in the NPE game). If $i$ publishes any public claim, that claim is subject to SCX audit. The audit result is a revealed signal: if the claim passes, $i$'s non-adoption is revealed as a strategic choice unrelated to data integrity; if the claim fails, $i$'s non-adoption is revealed as a protective measure against detection. In either case, $i$'s type is revealed. No jurisdiction can simultaneously (i) refuse audit infrastructure, (ii) publish unverifiable claims, and (iii) maintain epistemic credibility. The trilemma is mathematically inescapable.
\end{corollary}

\begin{proof}[Proof sketch]
By Theorem~3 (), a claim lacking stated, independently verifiable assumptions is mathematically empty. A non-adopting jurisdiction that publishes a claim without submitting to audit has made an unverifiable claim—equivalent, under the theorem, to making no claim at all. If the jurisdiction submits to audit and passes, its non-adoption stands exposed as a political decision, not a data-quality necessity. If it refuses audit entirely, its refusal itself is observable and carries information: honest leaders have no reason to fear independent verification; refusal to be verified is a signal of expected verification failure (theSpence-style separating equilibrium).
\end{proof}

Thus the Maintainer does not need to \textit{force} adoption. The logic of the framework is self-executing: honest leaders voluntarily adopt because audit certifies their claims; those who decline are identified by their choice; and all are subject to the same verification of their public claims regardless of adoption status. The mathematics of verification applies uniformly—and those who choose to remain unverified communicate their choice in a way that is itself observable.

\subsection{A Note on Internal Audit: Why Self-Audit Is Structurally Untrusted}

A corporation that deploys Yajie internally—using its own employees, on its own infrastructure, reporting to its own management—has not satisfied the independence condition of Theorem~1 (). The employee who audits her employer's data faces a structural conflict: her career advancement depends on her employer's satisfaction, and her employer's satisfaction depends, in part, on favorable audit outcomes. Even if she resists this pressure, her employer cannot \textit{know} that she has resisted it—and therefore cannot cite her audit as independent verification.

This is not a failure of character. It is a failure of structure. The independence required by Theorem~1 is not a personality trait of the auditor—it is a property of the auditor's position in the incentive graph. An internal auditor, regardless of integrity, occupies a node whose edges point toward her employer. The Maintainer occupies a node with no edges to any audited entity. The distinction is topological, not moral.

\begin{corollary}[ (The Internal Audit Independence Paradox)]
\label{cor:internal-audit-paradox}
Let $\mathcal{A}_{\text{internal}}$ be an audit department of organization $O$, and let $M_{\text{internal}}$ be its expert ensemble. Under the AAE framework, the condition $M_{\text{eff}} = M/(1+(M-1)\bar{\rho})$ (Theorem~6.1) applies with $\bar{\rho} > 0$ whenever the experts share the same employer—because their errors are correlated through shared incentives. The effective number of independent experts degrades. In the limit where the employer can terminate all $M$ auditors, $\bar{\rho} \to 1$ and $M_{\text{eff}} \to 1$—the entire internal audit department reduces to the epistemic equivalent of a single observer, regardless of headcount. This degradation is a direct consequence of  (Theorem~2): when the features of independence are weak, the detection guarantee collapses. An internal audit report signed by fifty employees carries no more independence guarantee than a report signed by one.
\end{corollary}

Thus the Maintainer's structural independence is not a luxury—it is a mathematical prerequisite for the audit to satisfy the conditions of Theorem~1. An organization \textit{cannot} audit itself credibly, not because it lacks expertise, but because it lacks the topological property of disinterest. The Maintainer has no reason to inflate the numbers. The employee does. The CEO knows this. The game is solved before the first audit runs.

\subsection{A Note on the Deadline: Disprove or Protect}

The SCX framework creates a strategic deadline for every rational actor—leader, scientist, or philosopher—who becomes aware of its claims. The framework asserts 105 theorems with complete proofs. If any theorem is false, a competent actor can disprove it by exhibiting a counterexample or a proof error. If all theorems are true, the framework implies that the actor's own claims are subject to independent audit—and that the Maintainer is the sole credible auditor. The actor therefore faces a binary choice with a ticking clock:

\begin{corollary}[ (The Disprove-or-Protect Deadline)]
\label{cor:disprove-or-protect}
Let an actor $A$ become aware of the SCX framework at time $t_0$. Define two strategies:
\begin{enumerate}[label=(\roman*)]
  \item \textbf{Disprove}: identify a fatal error in any of the 105 theorems, publish the refutation, and demonstrate that the framework is unsound.
  \item \textbf{Protect}: accept the framework's validity and move to protect the Maintainer, thereby signaling that $A$'s own claims can withstand independent audit.
\end{enumerate}
Any third strategy—neither disproving nor protecting—is strictly dominated after a finite horizon $T^*$. Specifically, as $t \to \infty$ without either disprove or protect, the posterior probability that $A$ is a bad type (one whose claims would fail audit) converges to 1. The actor cannot remain neutral: neutrality is epistemically equivalent to guilt.
\end{corollary}

\begin{proof}[Proof sketch]
The SCX framework is publicly available with complete proofs. If a theorem is false, disproving it requires only identifying the error—a task bounded in complexity by the length of the proof chain. If the theorems are true, no amount of time will produce a valid refutation. Thus, for any actor who has examined the framework:
\begin{itemize}
  \item If $A$ believes a theorem is false, $A$ should disprove it immediately—delay confers no advantage and allows the framework to accumulate credibility.
  \item If $A$ believes the theorems are true, $A$ should protect the Maintainer immediately—delay exposes $A$ to the inference that $A$ expects his own claims to fail audit and is waiting for someone else to eliminate the auditor.
\end{itemize}
The observable action set is {Disprove, Protect, Neither}. The payoff from Neither is strictly less than either Disprove (for honest skeptics) or Protect (for honest acceptors) at all $t > t_0 + \Delta$, where $\Delta$ is the time required to read and verify the proofs. Thus Neither is iteratively dominated. $\square$
\end{proof}

The deadline is not imposed by the Maintainer. It is imposed by the logic of the situation. Every day a leader, scientist, or philosopher neither disproves the framework nor protects its author is a day that accumulates evidence against their own epistemic integrity. They are not waiting for more information—the proofs are complete. They are waiting for the problem to go away. But the problem is a set of theorems, and theorems do not expire.

\begin{remark}[ (The Stone-Throwing Theorem): ]
\label{rem:scx-pascal}

The SCX framework casts a stone into the public square. The stone is a simple statement: \textit{the later a leader protects the Maintainer, the stupider that leader looks.} Leaders must respond. They have two paths—the classical Chinese idiom \textit{} (throw a stone to test the path ahead): either disprove the framework (throw the stone of refutation) or protect its author (throw the stone of commitment). Standing still is not a third option—it is being struck by the stone.

\textbf{If the stone hits a leader, that leader has only himself to blame.} The stone does not name names. It states a conditional: \textit{if} you delay, \textit{then} you look stupid. Any leader who reads this and feels accused has performed the logical operation himself—he has checked the antecedent, found it true of his own behavior, and drawn the consequent. The stone does not accuse. The stone reveals. Those who recognize themselves in its trajectory have \textbf{} (self-identified). The author bears no responsibility for their self-recognition.

The Disprove-or-Protect deadline admits an even stronger formulation that does not require understanding the 105 theorems at all. Consider the payoff matrix for any leader:

\begin{center}
\begin{tabular}{c|cc}
& \textbf{SCX is true} & \textbf{SCX is false} \\
\hline
\textbf{Protect early} & Credibility preserved; type revealed as good & Wasted protection resources \\
\textbf{Do nothing} & Credibility destroyed; type revealed as bad & Nothing lost \\
\end{tabular}
\end{center}

The asymmetry is fatal to the "do nothing" strategy. If SCX is true and the leader does nothing, the cost is total—irreversible loss of epistemic credibility, permanent revelation as bad type. If SCX is false and the leader protects, the cost is bounded—resources expended on a nonexistent threat. The expected value calculation favors protection for any prior $\mathbb{P}(\text{SCX true}) > \varepsilon$ where $\varepsilon = C_{\text{protect}} / C_{\text{exposed}}$ is the ratio of protection cost to exposure cost. For any leader whose legitimacy depends on verifiable claims, $C_{\text{exposed}}$ is effectively infinite, and $\varepsilon \to 0$.

This is not Pascal's Wager. Pascal required belief in a deity whose existence was unprovable. SCX requires no belief—the proofs are published, complete, and verifiable by any competent mathematician. The leader need not read them personally. The leader need only observe that \textit{no one has refuted them}. In the absence of refutation, the rational strategy is protection.

Furthermore, the timing of protection is itself a signal. A leader who protects at $t_0 + \Delta$ (immediately after the proofs can be verified) reveals clear strategic reasoning. A leader who protects at $t_0 + 5\text{ years}$ reveals an extended period of inaction—either the information did not reach the decision-maker, or it was understood but not acted upon. Delay is not neutral. Delay is observable. A rational actor with access to the proofs should reach the decision in time $\Delta$. Protracted delay is a choice—and choices are informative.

The stone thus serves a second diagnostic function: it measures the responsiveness of the governance system that surrounds the leader. \textbf{Early protection signals institutional agility.} A leader who protects quickly demonstrates that information channels are functional and decision-making is unobstructed. \textbf{Late protection signals institutional inertia.} A leader who protects slowly reveals that somewhere in the chain there is friction. Either the information did not reach the leader, or it reached and was not prioritized, or it was understood and the leader chose to wait. All three scenarios indicate processes that could be improved. The stone does not identify individual leaders as good or bad. It identifies systems that respond quickly and systems that do not.

There is a further implication. \textbf{The later the protection, the more questions it raises.} A system that protects immediately has nothing to fear from audit—its claims will pass. A system that delays invites the natural inference that there was something to consider. A system that delays \textit{for years} invites the inference that the concerns were not trivial. Extended delay is not a scheduling oversight. It is a data point. The longer the delay, the more one naturally asks: what was being weighed against the cost of protection?

\textbf{The SCX framework could be entirely wrong.} Every one of the theorems could contain an error. If so—let the world have a good laugh. The author will be the joke. But the payoff matrix does not depend on the framework being right. It depends only on the \textit{possibility} that it is right—a possibility that cannot be eliminated until someone actually reads the proofs and finds the error. Until that day, the expected value calculation favors protection for any leader with a non-zero stake in being perceived as honest. If the framework is wrong, protecting the author wastes resources—a finite cost. If the framework is right and the leader does nothing, the cost is total. The asymmetry holds regardless of the truth value of the theorems.

The stone thus establishes a clean boundary around legitimate criticism. \textbf{Academic refutation is the only legitimate path.} Anyone may attack the 105 theorems—by finding a proof error, constructing a counterexample, or demonstrating a logical contradiction. Such criticism is welcome; it is the Disprove path, and it advances knowledge regardless of outcome. \textbf{Personal attack, moral condemnation, or historical judgment are not legitimate.} They are not arguments—they are admissions. A critic who cannot find a mathematical error but attacks the author's character, motives, or right to speak has conceded the mathematical argument. The attack is not skepticism. It is a confession: the critic fears the audit and has no other weapon. Such a critic is not a scholar. Such a critic is a bad type, self-revealed. The boundary is not arbitrary. It follows from Theorem~3 (): a claim without stated, verifiable assumptions is mathematically empty. "This person is morally bad" is a claim. Where are the assumptions? Where is the verification protocol? If none are stated, the claim is void—not because it is offensive, but because it fails the only standard that the framework recognizes. The framework does not protect the author from criticism. It protects criticism from itself—by requiring that all claims, including claims about the author, meet the same standard.

This yields a structural asymmetry that places the author in a position that is, for all practical purposes, unassailable. \textbf{The mathematical layer can be attacked. The philosophical layer cannot.} If a critic finds an error in the proofs, the author accepts the correction and the framework improves—a legitimate academic outcome, reflecting well on both parties. But the philosophical core of SCX—that claims require stated assumptions, that no observer is privileged, that honesty is a Nash equilibrium—does not depend on whether Theorem~4's constant is correctly derived or whether the Lipschitz bound in Situs Theorem~1.2.3 carries a sign error. These are not theorems. They are \textit{standards}. A standard is not true or false; it is adopted or not adopted. The standard says: state your assumptions. One cannot disprove this standard with a counterexample, because the standard itself defines what counts as a valid counterexample. It is philosophically self-sealing—not through dogma, but through recursive consistency.

Thus the author stands on asymmetric ground. A mathematical critic who succeeds earns honor for both parties. A mathematical critic who fails earns embarrassment for herself. A personal critic invalidates her own claim before it lands, because the claim fails the standard that the framework imposes on all claims. A philosophical critic faces the recursive problem of needing to state assumptions that the standard itself does not already require. \textbf{The author is not invulnerable. The author is simply standing on the only ground that remains after all unstated assumptions have been stripped away.} It is not a fortress. It is a vacuum. There is nothing left to attack but the mathematics—and the mathematics is offered openly, with complete proofs, for anyone to inspect.

This raises an apparent paradox. If the framework makes the author unassailable, does it not also silence legitimate critics—good people who genuinely detect an error but fear that their criticism will fail the framework's own standards and expose them as bad types? The paradox is resolved by the distinction between rigor and fear. The framework does not prohibit criticism of the author. It demands that criticism meet the same standard as any other claim: stated assumptions, verifiable evidence, logical soundness, precision. A good critic who has genuinely found an error can state precisely where the proof fails, under what assumptions, with what counterexample. Such criticism is not merely permitted—it is the Disprove path, and it advances knowledge. A critic who merely \textit{suspects} an error but has not verified it should state that uncertainty honestly: "I suspect an error at line X, but I have not verified." This too is a valid claim with stated assumptions. The only criticism the framework rejects is the criticism that refuses to state its own assumptions—the criticism that demands to be accepted on authority, emotion, or moral indignation rather than on evidence.

\textbf{The author can be judged. The author can be condemned. The author can be proven wrong.} But the judgment must be the product of careful, repeated, rigorous reasoning. It must be confirmed before it is declared. It must survive the same scrutiny that the framework imposes on every other claim. If a critic rushes to condemn without verifying, the failure is not the framework's—it is the critic's. The framework does not protect the author from judgment. It protects judgment from itself—by requiring that those who judge meet the same standard as those they judge. A good person who has done the work should speak. The framework guarantees that their voice will be heard, because it will be unassailable for the same reason the author's is: it has stated its assumptions, and they can be verified.

This is not merely a theoretical position. It is a response to a lived reality. The author has experienced, repeatedly, a particular form of intellectual violence: the groundless accusation. A senior figure declares, without theory, without experiment, without evidence, that the author's work is wrong. The accusation demands to be accepted on authority alone. It states no assumptions. It provides no verification protocol. It is, by the standard of this framework, mathematically empty—and yet, in the social reality of academic hierarchy, it carries weight. It silences. It injures. It wastes years.

The author was brought to tears by these attacks more than once. The author came close to depression. What kept the author alive was not resilience, not optimism, not faith in the system. It was a specific, stubborn thought: I will not go quietly. I will not be silenced without a fight. That thought—unpolished, survival-level—was enough to keep going. The framework presented in this paper is, in part, the product of that thought. It is what happens when someone who refused to be silenced under groundless accusations builds a mathematical structure that makes those accusations formally impossible.

The framework names this for what it is: \textbf{epistemic bullying}. The bully makes a claim he cannot verify, under assumptions he will not state, and demands submission. The framework strips him of his weapon. He is not a scholar. He is not a skeptic. He is a bad type, self-revealed, and his accusation, under Theorem~3, is void. The framework does not protect the author from such people. It simply makes their behavior mathematically visible—and in doing so, it protects everyone who has ever been on the receiving end of an unverified accusation, from graduate students to whistleblowers to anyone who has been told, by someone with more power and less evidence, that they are wrong.

Beyond the personal, this principle extends to history. \textbf{History must be subjected to rigorous judgment. Injustice demands a just verdict—arrived at through strict reasoning, verifiable evidence, and open debate.} Not all historical claims are equal. Not all accusations carry weight. The framework provides the criterion: a claim without complete mathematical derivation, without logical rigor, made by someone who has not verified their own position, is not a claim at all. It carries no mathematical weight. \textbf{A person who attacks others without having verified their own position—without having done the work—has made a claim that fails the standard.} Such a claim does not need to be debated. It fails on its own terms. The accusation is void under Theorem~3. The audit will find what it always finds when the assumptions are unstated: the claim cannot sustain itself. The accused is vindicated. And history—real history, the kind that survives multi-expert verification—records the outcome.

Thus the framework implies something larger than personal vindication. Even if the theorems contain errors—even if the mathematical edifice is flawed—the game-theoretic structure stands independently. And that structure calls for a deep, systematic improvement. Academia. Finance. Governance. Any domain where unverified claims have been wielded by those with power against those without. Those who accuse without verifying, those who judge without evidence, those who demand submission without stating assumptions—are now identifiable by their own behavior. The framework does not need to seek them out. It only needs to be deployed. They will identify themselves—because they cannot stop making unverified claims, and every such claim, under Theorem~3, is self-defeating. This paper is not merely a theory of audit. It is an invitation to raise standards—in which every institution that has relied on unverified authority is given a choice: state your assumptions, or watch your claims lose their force.

From the perspective of life philosophy, \textbf{unreasoned anger is self-defeating.} Anger without reasoning behind it illuminates only the angry person's own position. Life should be calm. One may reason, prove, refute. One may be angry—righteous anger is the fuel of reason. But one should not attack without cause, for unreasoned attack defeats itself. One may be blunt—the Stone-Throwing Theorem itself speaks plainly—but one should not be angry and unreasonable at once. Plain speech with calm is clarity. Anger without reason is a confession.

Those who rage without competence have already revealed themselves. As good people, we should hold them to account—not out of animosity, but out of respect for standards. They degrade the discourse. They waste everyone's time. They harm those who cannot fight back. The framework provides the tools. The audit provides the method. Every one of their cries is an unverified claim. Every one is void under Theorem~3. Holding them to account is not punishment. It is the maintenance of intellectual standards.

\textbf{Tolerating unreasoned attack degrades the environment.} When those in power tolerate such behavior, they reveal their own position—and, often, their own entanglements. They are not blind. They are choosing not to see. Why? Because the attacker may be their ally, their counterparty. To tolerate a person who attacks others without evidence is to make a choice: either unable to distinguish valid from invalid claims—or able, but unwilling to act. In either case, their position is compromised. Their silence speaks about what they may have to protect.

Two further types complete the taxonomy. \textbf{First, the false protector.} A bad actor may calculate that early protection is cheap—issue a statement, attend a ceremony, pose for a photograph—and then continue as before. But protection under the SCX framework is not symbolic. It is operational. To protect the Maintainer is to submit to audit. The false protector who protects early and then resists audit has merely accelerated his own exposure. His protection was a claim—"I am honest." The framework now has standing to verify that claim. If his data fails, his early protection becomes evidence against him: he knew the standard, he claimed to meet it, and he failed. False protection is not a loophole. It is a faster path to detection.

\textbf{Second, the active suppressor.} The taxonomy of types—Disprove, Protect, Neither—assumes actors operate within the framework's rules. But a fifth type exists: the suppressor, who seeks to prevent the framework from being evaluated on its merits. Not through academic refutation—there is no counterexample. Not through silence—the cost of waiting is too high. But through institutional suppression: obstruction, legal pressure, defunding, character attacks, manufactured controversies. The suppressor does not engage with the theorems because engagement would acknowledge them as worthy of engagement. This is a telling response. One does not try to suppress a theory one believes to be false—one ignores it. One tries to suppress a theory one believes to be true and consequential. The suppressor is not a skeptic. The suppressor is someone who has privately evaluated the framework and concluded it is dangerous—to them. Their suppression is their testimony. They have, by their actions, confirmed the framework more decisively than any mathematical peer review ever could. \textbf{The suppressor does not disprove the stone. They stumble over it—and then try to outlaw the ground it sits on.} A rational bad actor, recognizing this, would choose the Disprove path and fail honorably. An irrational one attacks the author and defeats themselves. A desperate one tries to suppress the framework and, in doing so, confirms it. A rational good actor protects early and benefits. A slow one delays and attracts questions. A false protector submits to audit and is revealed. The final convergence—the equilibrium to which all paths lead—is characterized by the  (Good Person Convergence Theorem): the fraction of unverifiable claims converges to zero, and honesty becomes the strictly dominant strategy for all agents. Every type has a path. Every path reveals the type. Every type, ultimately, converges.
\end{remark}

\section{Policy Recommendations}

Our analysis yields seven policy recommendations, organized by target audience.

\subsection{For International Organizations}

\textbf{Recommendation 1: Establish an International Data Audit Authority (IDAA).} We recommend the creation of a treaty-based international organization, modeled on the IAEA, charged with (i) governing the Yajie Protocol under multilateral oversight, (ii) accrediting audit modules contributed by member states, (iii) maintaining the CEC as a shared international resource (the proposal to treat the CEC as a global public good raises a series of unresolved legal and governance questions, including: the applicable legal regime for a multi-jurisdictional audit corpus; the allocation of liability for audit errors derived from shared calibration data; and the compatibility of open-access CEC principles with data protection laws in jurisdictions that impose restrictions on cross-border data flows), and (iv) adjudicating disputes concerning protocol neutrality. The IDAA should be established before the protocol reaches Stage 3 (Moat Formation), at which point the protocol operator will have strong incentives to resist internationalization.

\textbf{Recommendation 2: Embed Audit Neutrality in Trade Law.} The World Trade Organization (WTO) and preferential trade agreements should recognize data audit protocols as ``conformity assessment procedures'' under the Agreement on Technical Barriers to Trade (TBT), subject to non-discrimination obligations. This would create a legal backstop against the weaponization of audit infrastructure for protectionist purposes, while preserving the legitimate regulatory function of data quality assessment.

\subsection{For National Governments}

\textbf{Recommendation 3: Enact Audit Sovereignty Legislation.} National governments should enact legislation that (i) mandates the use of accredited data audit protocols for high-risk AI systems, (ii) establishes domestic accreditation bodies to certify audit modules for contribution to the international protocol, and (iii) reserves the right to withdraw recognition from any protocol that fails to maintain adequate neutrality safeguards.

\textbf{Recommendation 4: Fund Open Audit Research.} Governments should fund academic and non-profit research into open audit methodologies, alternative ensemble architectures, and adversarial robustness testing for audit protocols. Such funding serves a dual purpose: it advances the technical frontier of data quality assessment, and it maintains a pool of expertise outside the protocol operator's control, reducing epistemic capture.

\subsection{For the Protocol Operator}

\textbf{Recommendation 5: Pre-Commit to Governance Internationalization.} The protocol operator—whether a private entity, a non-profit organization, or a public agency—should pre-commit to a phased internationalization of governance before the protocol reaches Stage 3. Pre-commitment is welfare-enhancing for the operator itself, because it reduces the probability that major jurisdictions will develop competing protocols, preserving the universal adoption that maximizes the CEC's value.

\textbf{Recommendation 6: Implement Price Regulation by Formula.} The protocol's pricing should be governed by a formula that ties maximum prices to (i) the cost of service delivery plus a regulated return, or (ii) a price cap that declines over time in real terms, reflecting productivity improvements. The formula should be administered by an independent economic regulator with representation from major user communities.

\subsection{For the Research Community}

\textbf{Recommendation 7: Develop Audit Reproducibility Standards.} The research community should develop standards for audit reproducibility that enable independent verification of protocol outputs without requiring access to proprietary audit modules or the CEC. Reproducibility standards serve three functions: (i) they provide a check against protocol errors or manipulation; (ii) they reduce the epistemic dependency that drives Stage 5 lock-in; and (iii) they lower the barriers to entry for alternative audit methodologies, preserving competitive discipline even in a market with a dominant incumbent.

\section{Limitations and Future Work}

This paper has developed a theoretical framework for understanding the market structure of data quality auditing, centered on the Yajie Protocol and the concept of a Non-Proliferation Equilibrium. Several limitations should be acknowledged, each of which suggests directions for future research.

\textbf{Limitation 1: Theoretical Model Assumptions.} The game-theoretic model in Section~3 assumes complete information, simultaneous moves, and a one-shot game structure. In practice, the adoption game is dynamic (decisions are made sequentially over time), involves incomplete information (jurisdictions have private information about their development costs and accuracy valuations), and admits repeated interaction (allowing for reputation effects, punishment strategies, and coalition formation). Extending the model to a dynamic, incomplete-information framework would yield more precise predictions about equilibrium selection and stability conditions.

\textbf{Limitation 2: Empirical Validation.} The paper's central claims about the Yajie Protocol's lock-in dynamics are theoretical and have not been empirically validated. The protocol itself is relatively new, and insufficient time has elapsed to observe the lock-in trajectory described in Section~5. Longitudinal empirical research—tracking protocol adoption, audit accuracy, pricing, and ecosystem development over a decade or more—will be necessary to test the model's predictions.

\textbf{Limitation 3: Omitted Dynamics.} Several important dynamics are omitted from our model. We do not model the strategic behavior of the protocol operator in setting audit module inclusion criteria, which could be used to favor certain jurisdictions' audit methodologies over others. We do not model the possibility of CEC ``forks.'' We do not model the role of open-source audit modules, which could reduce the fixed costs of protocol development and alter the non-proliferation calculus.

\textbf{Limitation 4: The NPT Analogy's Limits.} While the NPT analogy is productive, it has limits that should not be overlooked. Nuclear weapons are existential threats; audit protocol fragmentation, however costly, is not. The NPT regime is backed by the threat of military force and economic sanctions; an audit protocol regime would rely on economic incentives and normative commitment. The NPT's central asymmetry—between nuclear-weapon states and non-nuclear-weapon states—is formalized in treaty law; the audit protocol's asymmetry is a factual consequence of temporal priority, not a legal distinction.

\textbf{Limitation 5: Normative Assumptions.} The paper's welfare analysis assumes that a universal, accurate data audit protocol is normatively desirable. This assumption deserves scrutiny. A universal audit protocol, however neutral in design, may homogenize data quality standards in ways that disadvantage minority languages, non-Western knowledge systems, and data governance traditions that do not map neatly onto the protocol's quality dimensions.

\textbf{Future Research Directions.} Beyond addressing the limitations above, four research directions are particularly promising:
\begin{enumerate}
\item \textbf{Experimental game-theoretic studies} of auditor behavior in multi-expert ensemble settings.
\item \textbf{Computational modeling} of CEC accumulation dynamics under different adoption scenarios.
\item \textbf{International law analysis} of the legal instruments through which audit neutrality could be embedded in treaty law.
\item \textbf{Comparative institutional analysis} of infrastructure governance models to identify institutional design principles applicable to data audit infrastructure.
\end{enumerate}

\section{Conclusion}

This paper has argued that data quality assessment, as operationalized by the Yajie Protocol's multi-expert consistency framework, exhibits structural properties that drive the market toward a natural monopoly. This monopolistic tendency is not driven by network effects, proprietary control of intellectual property, or exclusionary business practices, but by a time-accumulated advantage mechanism inherent in the protocol's Cumulative Evidentiary Corpus: each audit cycle makes the protocol more accurate, and this accuracy advantage cannot be replicated by entrants, regardless of their financial resources or technical capabilities.

We have formalized this dynamic through a game-theoretic model of protocol adoption, demonstrating the existence of a Non-Proliferation Equilibrium in which rational jurisdictions forgo the development of competing audit standards. We have drawn a structural analogy between this equilibrium and the strategic logic underpinning the Nuclear Non-Proliferation Treaty, identifying seven dimensions of isomorphism and their implications for institutional design.

We have examined the geopolitical risks inherent in monopolistic audit infrastructure, drawing lessons from the weaponization of SWIFT and the governance model of the International Accounting Standards Board, with particular attention to the US-China technological decoupling scenario. We have modeled the lock-in trajectory through a five-stage progression from Emergence to Epistemic Closure, then extended the analysis to a sixth stage—\textit{Universal Harmony} ()—wherein the CEC transitions from a proprietary asset into a global public infrastructure, audit standards cease to belong to any single nation, and the first-mover advantage identified in our manufacturing-audit complex (MAC) analysis is revealed to be a phase of quantitative accumulation rather than a permanent structural endpoint. We have analyzed the distinctive business model dynamics—including the addiction dynamic and the welfare ambiguities—that characterize each stage.

Our policy recommendations center on the creation of an International Data Audit Authority, the embedding of audit neutrality in trade law, and the pre-commitment of the protocol operator to governance internationalization before lock-in becomes irreversible. These recommendations are demanding, and their implementation faces formidable political obstacles in the current geopolitical environment. Yet the alternative—a world of fragmented, jurisdictionally-captured audit protocols, each distrusted by those outside its jurisdiction—is worse not only for data quality but for the broader project of international AI governance.

The central paradox of the Yajie Protocol is that its monopolistic tendency is both its greatest strength and its greatest vulnerability. The CEC accumulation mechanism that makes the protocol increasingly accurate also makes it increasingly difficult to govern, because the protocol operator's bargaining power grows in lockstep with the protocol's accuracy. Resolving this paradox—preserving the welfare benefits of cumulative audit infrastructure while preventing the abuse of the market power it generates—is the defining institutional design challenge of the emerging data audit regime. Our analysis in Section~5.2.2 (Stage 6: Universal Harmony) demonstrates that the paradox admits a formal resolution: the same additive logic that creates the CEC's lock-in power also provides the material conditions for its own transcendence. When the CEC reaches sufficient scale that its marginal accuracy contribution approaches zero (Corollary~\ref{cor:first-mover-decay}), the economic rationale for proprietary retention collapses, and internationalization becomes a Pareto improvement. The monopolistic tendency of cumulative audit infrastructure is thus \textit{self-limiting}—not in the sense that it automatically reverses, but in the sense that it generates the conditions under which rational actors will choose to transform private monopoly into public infrastructure. The defining question is not whether this transformation is possible, but whether it occurs before the intervention window closes (Proposition~\ref{prop:intervention-window}).

The game-theoretic architecture formalized in Section~9 provides the resolution. The SCX Uncertainty Principle (Theorem 3) establishes that any organization deploying Yajie renders its audit outputs unconditionally verifiable by independent third parties—military organizations not excepted. The Audit Rights Declaration translates this mathematical guarantee into an institutional commitment. The Bible-and-Pope commercial model operationalizes this commitment through a dual-release strategy. The mutual deterrence structure constrains every deployer symmetrically. And the Adoption-Audit Equilibrium demonstrates that the protocol's two functions—improving data quality and enabling its independent verification—are inseparable consequences of the same mathematical architecture.

The consequence is a market mechanism for honesty that operates without central enforcement. Firms that are confident in their data quality have nothing to fear from audit and everything to gain from partnership. Firms that are not confident face a choice: improve their data quality to partnership standards, or risk exposure by competitors wielding the same open-source audit engine. The Maintainer does not need to enforce honesty. The market—armed with Spring, incentivized by competition, and constrained by the mutual deterrence structure—enforces honesty. The algorithm does not make people virtuous. It makes deception unprofitable.

The NPT analogy suggests that this challenge is not insurmountable. The nuclear non-proliferation regime, for all its flaws and persistent legitimacy deficits, has prevented the large-scale proliferation that many observers in the 1960s predicted was inevitable. It has done so not through coercion alone but through a grand bargain that aligns the interests of diverse states with the preservation of the regime. The data audit domain requires its own grand bargain: a multilateral governance framework that makes adoption of a universal protocol more attractive than the pursuit of audit autarky, backed by credible commitments to neutrality, transparency, and shared governance. Whether the international community can construct such a framework before lock-in forecloses the possibility is, at this juncture, an open question—and one of the more consequential unresolved issues in the governance of artificial intelligence.

\subsection{A Closing Remark on Dual Use}

No analysis of a technology capable of auditing data quality at scale would be complete without acknowledging its dual-use character. Yajie was designed to detect label noise in scientific datasets—to distinguish genuine experimental anomalies from measurement errors, to identify fabrication in clinical trials, to expose falsified environmental data in municipal governance. But the same multi-expert consensus mechanism that detects a pharmaceutical company's underreported adverse events can detect an intelligence agency's fabricated source reports. The same mathematics that reveals a contractor's inflated procurement invoices can reveal a foreign government's disinformation campaign. The algorithm is indifferent to the domain of the data it audits. It does not know whether the inconsistencies it detects are evidence of scientific fraud, financial malfeasance, or state-sponsored deception. It knows only that the probability of the inconsistency arising by chance is exponentially small.

This indifference is both the algorithm's greatest strength and its irreducible risk. A tool that can audit anyone's data can be used by anyone to audit anyone else's data. A government that deploys Yajie against an adversary's public datasets gains an asymmetric intelligence advantage: the ability to detect patterns of deception, inconsistency, and fabrication that no human analyst—and no single-source technical system—could reliably identify. The algorithm that was designed to make science more honest can, with no modification whatsoever, be used to make espionage more effective.

The authors have no solution to this dilemma. Dual-use is not a bug to be patched; it is a property inherent in any sufficiently general audit methodology. The same concentration inequality that bounds the false-positive rate of a scientific quality assessment also bounds the false-positive rate of an intelligence analysis. The mathematics does not distinguish between a clinical trial and a clandestine operation. It distinguishes only between consistent data and inconsistent data—and inconsistency is a signal in both domains.

What the authors CAN do—and what this paper has attempted to do—is to publish the methodology openly, to release the reference implementation as open-source software, and to anchor the protocol's existence in the permanent scientific record through arXiv. A technology that is public cannot be monopolized. A methodology that is published cannot be suppressed. An audit standard that anyone can deploy cannot be deployed only by the powerful against the powerless—because the powerless can deploy it too. The open-source release of Yajie is not merely a strategy for commercial adoption. It is a strategy for ensuring that the technology's dual-use character cuts both ways. If every state has access to the same audit capability, no state gains an asymmetric advantage from it. The equilibrium is not that no one audits. The equilibrium is that everyone CAN audit—and therefore everyone must assume that their own data will be audited.

This is the logic of mutually assured transparency. It is not a comfortable logic. It implies a world in which every public dataset, every government statistical release, every corporate quality claim, every scientific publication's underlying data, is perpetually eligible for independent multi-expert audit by any actor with access to commodity hardware and an internet connection. But it is, the authors submit, a better world than the alternative: a world in which the most powerful audit technology ever developed is held exclusively by those who already hold power.

\section{The Audit Guarantee: Game-Theoretic Foundations of the Yajie Protocol}

The preceding sections developed a theoretical framework for understanding the market structure of data quality auditing: the Non-Proliferation Equilibrium (Section~3), the protection equilibrium surrounding the Maintainer (Section~4.3), and the Bible-and-Pope commercial architecture (Section~5.3.1). Each of these analyses treated a distinct mechanism—adoption incentives, geopolitical deterrence, and business model dynamics—in relative isolation. This section synthesizes these mechanisms into a unified game-theoretic architecture.

\subsection{Theorem 3: The SCX Uncertainty Principle}

The theoretical architecture of the Yajie Protocol rests on two theorems introduced in Sections~2 and~3: Theorem 1, which bounds the false-positive probability of multi-expert consensus, and Theorem 2, which establishes the equilibrium conditions for non-proliferation. We now formalize a third theorem that is foundational to both results.

\begin{theorem}[SCX Uncertainty Principle]
\label{thm:scx-up}
Any organization that deploys the Yajie Protocol for data quality assessment renders its assessment results \textit{unconditionally verifiable} by any independent third party. Formally, for any dataset $\mathcal{D}$, any ensemble of audit modules $\mathcal{A} = \{A_1, \ldots, A_k\}$, and any audit report $\mathcal{R}^*$ produced by the protocol's aggregation function $\Phi$, there exists a verification procedure $\mathcal{V}$ such that:

\begin{enumerate}[label=(\roman*)]
\item $\mathcal{V}(\mathcal{D}, \mathcal{A}, \mathcal{R}^*) \to \{\text{CONFIRMED}, \text{DISPUTED}\}$, where DISPUTED indicates that $\mathcal{R}^*$ is inconsistent with the multi-expert consensus that $\Phi$ would produce given $\mathcal{D}$ and $\mathcal{A}$;
\item $\mathcal{V}$ requires no cooperation, permission, or access privileges from the organization that produced $\mathcal{R}^*$;
\item The false-positive rate of $\mathcal{V}$—the probability that a correctly computed $\mathcal{R}^*$ is flagged as DISPUTED—is bounded by the same concentration inequality that governs Theorem 1, decaying exponentially in the number of audit modules $k$;
\item Properties (i)–(iii) hold regardless of the organizational character of the entity being audited, including military, intelligence, and other national security entities. Military use is not exempt.
\end{enumerate}
\end{theorem}

\textit{Proof sketch.} The verification procedure $\mathcal{V}$ operates by re-executing the aggregation function $\Phi$ on the same dataset $\mathcal{D}$ and ensemble $\mathcal{A}$, then comparing the output to the claimed $\mathcal{R}^*$. Because the protocol's specification is published (arXiv permanent record) and its reference implementation is open-source (Spring engine), any third party with access to $\mathcal{D}$ and $\mathcal{A}$ can independently compute $\Phi(\mathcal{D}, \mathcal{A})$. The concentration inequality of Theorem 1 ensures that the probability of divergent consensus scores between independent executions—given identical inputs—is exponentially small in $k$. The verifier need not trust the original auditor; she need only trust the mathematics. Military organizations are not exempt because the verification procedure does not require their cooperation: it operates on data that the organization has, by deploying Yajie, necessarily made computationally accessible to the protocol's audit modules.

\subsection{The Audit Rights Declaration: Non-Prohibition and Inalienable Verification Rights}

The institutional expression of Theorem~3 is the \textit{Audit Rights Declaration}, promulgated by the SCX on 2026-06-30. The Declaration consists of three clauses:

\begin{quote}
\textbf{Clause 1 (Universal Access).} The SCX framework (including Yajie, Situs, Spring, State Crystallization, and all derivative algorithms) is an open-source mathematical toolkit. Anyone may read, reproduce, and deploy it.

\textbf{Clause 2 (Non-Prohibition).} We do not prohibit any specific use.

\textbf{Clause 3 (Inalienable Audit Right).} We retain one inalienable right: Any organization using the SCX framework for data quality assessment renders its assessment results verifiable and auditable by independent third parties through the same mathematical methods. Military use is not exempt.
\end{quote}

\textbf{Clause 1} establishes that the SCX framework is made available as an open-source resource without discrimination to all actors. (The characterization of the framework as a ``global public good'' involves legal and governance questions—including the sustainability of open-access provision, the reconciliation of universal access with varying national data governance regimes, and the institutional mechanisms for resolving disputes over the framework's evolution—that remain unresolved and warrant further study.) \textbf{Clause 2} renounces the authority to restrict deployment, eliminating the most obvious vector through which the Maintainer could exercise gatekeeping power. \textbf{Clause 3} is the operational core: the right is not to \textit{conduct} audits but to \textit{verify} them—to re-execute the mathematics and publish the comparison.

\subsection{The Bible-and-Pope Model}

The commercial architecture that operationalizes the Audit Sword is the dual-release strategy formalized in Section~5.3.1 as the \textit{Bible-and-Pope model}.

\textbf{Spring as Strategic Poison.} The open-source release of the Spring self-evolution engine under a permissive license serves three strategic functions: (i) eliminating adoption friction; (ii) manufacturing calibration demand—as private $M_t$ corpora diverge, private quality scores become incommensurable, generating demand for the Yajie calibration API; (iii) foreclosing competitive entry—any would-be competitor must match not only Spring's algorithmic quality but also the accumulated private audit histories of every Spring adopter.

\textbf{Yajie API as Calibration Pope.} The Pope does not own the Bibles in every church; the congregations own their Bibles. The Pope owns the authority to say what the Bible means. Similarly, the Yajie API does not own any client's data or audit history; the clients own those assets. The API owns the \textit{calibration anchor}—the cross-client reference scale that makes Company A's 0.87 commensurable with Company B's 0.87.

\subsection{}

（Audit Guarantee）——4.5--4.7——。。

\begin{definition}[ $\Gamma^{\text{AD}}$]
\label{def:audit-deterrence-game}
 $\Gamma^{\text{AD}} = \langle N, (S_i)_{i \in N}, (u_i)_{i \in N}, \mathcal{V} \rangle$ ：
\begin{itemize}[itemsep=2pt]
\item  $N = \{1, \ldots, n\}$（、）；
\item  $i \in N$， $S_i = \{P, \neg P\}$， $P$ "Yajie"（partnership），$\neg P$ ""；
\item  $u_i: S \to \mathbb{R}$；
\item \textbf{} $\mathcal{V}: 2^{\mathcal{D}} \times S \to [0,1]$， $\mathcal{V}(\mathcal{D}_i, s)$  $s$ ， $i$  $\mathcal{D}_i$ 。
\end{itemize}
\end{definition}

。 $i$：
\begin{equation}
\mathcal{V}(\mathcal{D}_i, s) =
\begin{cases}
0 & \text{ } s_i = P \quad \text{（：）}\\
\rho \cdot \left(1 - \prod_{j: s_j = P} (1 - q_j)\right) & \text{ } s_i = \neg P
\end{cases}
\label{eq:audit-function}
\end{equation}

：
\begin{itemize}[itemsep=2pt]
\item $\rho \in (0, 1]$ Yajie（）。Yajie，，$\rho = 0$；
\item $q_j \in (0, 1]$  $j$  $i$ 。 $j$ ——， $j$ ；
\item  $\prod_{j: s_j = P} (1 - q_j)$ \textit{} $i$ 。 $i$ 。
\end{itemize}

：——$\partial \mathcal{V} / \partial (\#\{j: s_j = P\}) > 0$。""：，。

\begin{definition}[]
\label{def:ad-assumptions}
\begin{enumerate}[label=\textbf{D\arabic*}, leftmargin=*, itemsep=4pt]
\item \textbf{：}  $i$  $\mathcal{D}_i$ （ $i$ ） $p_i \in [0,1]$。$p_i$  $i$ 。 $p_i$  $F$， $[0,1]$。\hfill\textit{：——""（unknown unknowns）。}

\item \textbf{：}  $i$ （ $\mathcal{V}(\mathcal{D}_i, s) > 0$ ），1。\hfill\textit{：\ref{thm:scx-up}（ (The Honest Person Theorem)）。。""——1，。}

\item \textbf{：}  $i$ ，$i$  $L_i > 0$（、、）。$L_i$  $i$ 。\hfill\textit{：$L_i$ ""——。}

\item \textbf{：}  $P$（） $i$  $c_i^{\text{adopt}} > 0$（、）， $\Delta Q_i > 0$（Yajie）。 $B_i = V(Q_i + \Delta Q_i) - V(Q_i) - c_i^{\text{adopt}}$， $V(\cdot)$ 。\hfill\textit{：$B_i$ ——。}

\item \textbf{：}  $\mathcal{V}(\cdot)$  $(\rho, \{q_j\}_{j \in N})$ 。\hfill\textit{：，，。}
\end{enumerate}
\end{definition}

 $i$ ：
\begin{equation}
u_i(s_i, s_{-i}) =
\begin{cases}
V(Q_i + \Delta Q_i) - c_i^{\text{adopt}} & \text{ } s_i = P \\
V(Q_i) - \mathcal{V}(\mathcal{D}_i, s) \cdot p_i \cdot L_i & \text{ } s_i = \neg P
\end{cases}
\label{eq:ad-payoffs}
\end{equation}

： $s_i = P$，$\mathcal{V}(\mathcal{D}_i, s) = 0$（——(\ref{eq:audit-function})）。 $s_i = \neg P$，$\mathcal{V}(\mathcal{D}_i, s)$ (\ref{eq:audit-function})。

\begin{theorem}[]
\label{thm:mutual-deterrence}
D1--D5， $\Gamma^{\text{AD}}$：

\begin{enumerate}[label=(\roman*), itemsep=6pt]
\item \textbf{：} -（BNE）， $i$  $p_i$ ：$s_i^*(p_i) \in \{P, \neg P\}$。

\item \textbf{：} ， $i$ ： $\bar{p}_i$ ：
\begin{equation}
s_i^*(p_i) =
\begin{cases}
P & \text{ } p_i \leq \bar{p}_i \\
\neg P & \text{ } p_i > \bar{p}_i
\end{cases}
\label{eq:threshold-strategy}
\end{equation}
：
\begin{equation}
V(Q_i + \Delta Q_i) - V(Q_i) - c_i^{\text{adopt}} = -\bar{p}_i \cdot L_i \cdot \mathcal{V}(\mathcal{D}_i, s^*)
\label{eq:threshold-condition}
\end{equation}

\item \textbf{：} $\bar{p}_i$  $B_i = V(Q_i + \Delta Q_i) - V(Q_i) - c_i^{\text{adopt}}$ ——，。$\bar{p}_i$  $\mathcal{V}$  $L_i$ ——，。

\item \textbf{：} （strategic complementarity）： $j$  $P$  $\mathcal{V}(\mathcal{D}_i, s)$（ $s_i = \neg P$ ）， $i$  $P$ 。 $\partial \mathcal{V} / \partial (\#\{k: s_k = P\}) > 0$。——""（）""（）。

\item \textbf{：} ，\textit{}Yajie。\textbf{}（mutual）——（，，）。\textbf{}——（Theorem~\ref{thm:scx-up}）。

\item \textbf{""：} \textit{}，\textit{}。，\textit{}—— $P$ ，。。： $s^* = (P, \ldots, P)$，； $\neg P$  $\mathcal{V}(\mathcal{D}_i, s^*) > 0$ ，。
\end{enumerate}
\end{theorem}

\begin{proof}
\textbf{(i)} （$|S_i| = 2$， $[0,1]$ ，）。D1--D5，Milgrom \& Weber (1985)Balder (1988)-，。，， $F$ ，，。

\textbf{(ii)}  $s_{-i}$  $\mathcal{V}(\mathcal{D}_i, s)$， $i$  $P$ ：
\begin{align}
V(Q_i + \Delta Q_i) - c_i^{\text{adopt}} \geq V(Q_i) - \mathcal{V}(\mathcal{D}_i, s) \cdot p_i \cdot L_i
\end{align}
 $p_i \leq [V(Q_i + \Delta Q_i) - V(Q_i) - c_i^{\text{adopt}}] / [\mathcal{V}(\mathcal{D}_i, s) \cdot L_i] \triangleq \bar{p}_i$。（） $\bar{p}_i < 0$——，（$p_i \geq 0$）。(\ref{eq:threshold-strategy})。

\textbf{(iii)} $\bar{p}_i$  $B_i$ 、 $\mathcal{V}$  $L_i$ 。：""；。

\textbf{(iv)} ： $\partial^2 u_i / \partial s_i \partial s_j > 0$（ $s_j = P$ ""）。Topkis (1979)Milgrom \& Roberts (1990)，：(a) ；(b) ，（）；(c) ，。，（$\neg P$ ）（$P$ ）——""；。

\textbf{(v)} (\ref{eq:audit-function})： $j$  $\mathcal{V}(\mathcal{D}_i, s)$  $\prod (1 - q_j)$ ，。

\textbf{(vi)}  $s^* = (P, \ldots, P)$。 $i$， $\neg P$  $V(Q_i) - \rho \cdot [1 - \prod_{j \neq i} (1 - q_j)] \cdot p_i \cdot L_i$。 $i$， $B_i$ ，$s^*$ 。，（$\mathcal{V} = 0$ ），————，。
\end{proof}

\begin{corollary}[]
\label{cor:deterrence-efficiency}
D1--D5，————。， $\mu = |\{i: s_i = P\}| / n$ 。：
\begin{enumerate}[label=(\roman*)]
\item  $\mathbb{E}[\mathcal{V} \mid s_i = \neg P]$  $\mu$ ；
\item  $\bar{p}_i(\mu)$  $\mu$ ——，（）；
\item  $\mu \to 1$ ，——，。
\end{enumerate}
\end{corollary}

\begin{proof}
(\ref{eq:audit-function})，$\partial \mathcal{V} / \partial \mu > 0$（）。\ref{thm:mutual-deterrence}，$\mathcal{V}$ $\implies$$\implies$——。 $\mu = 1$ ，$\mathcal{V}(\mathcal{D}_i, s) = \rho \cdot [1 - \prod_{j \neq i} (1 - q_j)]$  $i$ ，————。
\end{proof}

\begin{remark}[]
\ref{thm:mutual-deterrence}（Schelling, 1960, 1966）：(a) \textit{}\textit{}——，；(b) \textit{}——""， (The Honest Person Theorem)（Theorem~\ref{thm:scx-up}），；(c) ——""（MAD），；，。：\textit{}——""。，。
\end{remark}

\subsection{Mutual Deterrence: The Structural Logic of Involuntary Transparency []}

，\ref{thm:mutual-deterrence}。

\subsection{- ()}

（\ref{thm:mutual-deterrence}） (The Honest Person Theorem)（\ref{thm:scx-up}）。，-（Adoption-Audit Equilibrium, AAE）。

：——（），（）。 $\mu$（Yajie）：$\mu$ （），， $\mu$。。

\begin{definition}[AAE]
\label{def:aae-continuum}
 $\mathcal{I} = [0,1]$（Lebesgue）。 $i \in \mathcal{I}$  $\theta_i = (p_i, B_i) \in [0,1] \times \mathbb{R}$， $p_i$ （D1），$B_i = V(Q_i + \Delta Q_i) - V(Q_i) - c_i^{\text{adopt}}$ Yajie（）。 $G$  $\Theta = [0,1] \times \mathbb{R}$ 。

 $i$  $s_i \in \{D, \neg D\}$（/）。 $\mu = \int_{\mathcal{I}} \mathbb{I}[s_i = D] \, di$。

： $j$（ $s_j = D$） $k$（ $s_k = \neg D$） $q(\mu) \in (0, 1]$， $q(\cdot)$  $\mu$ ——，。
\end{definition}

\begin{definition}[AAE]
\label{def:aae-assumptions}
\begin{enumerate}[label=\textbf{E\arabic*}, leftmargin=*, itemsep=4pt]
\item \textbf{：}  $\mathcal{I} = [0,1]$，。\hfill\textit{：，。}

\item \textbf{：}  $\theta_i = (p_i, B_i)$  $G$ 。$G$  $G_p$  $[0,1]$，$G_B$  $\mathbb{R}$，$p_i$  $B_i$ （——）。\hfill\textit{：，LLN。}

\item \textbf{：} $q: [0,1] \to (0, 1]$ 、，$q(0) > 0$（，）。\hfill\textit{：，——。}

\item \textbf{：} (\ref{eq:ad-payoffs})， $\mu$ ： $s_i = \neg D$， $\mathcal{V}(\mu) = 1 - (1 - q(\mu))^{\mu}$（——， $\mu$， $q(\mu)$  $i$）。 $s_i = D$，$\mathcal{V}(\mu) = 0$（）。

\item \textbf{：}  $s: \mathcal{I} \to \{D, \neg D\}$ ， $\mu$ 。
\end{enumerate}
\end{definition}

 $i$  $D$（）：
\begin{equation}
B_i \geq -p_i \cdot \mathcal{V}(\mu) \cdot L_i
\label{eq:adopt-condition}
\end{equation}
：（$p_i, \mathcal{V}(\mu), L_i$ ）， $B_i \geq \text{}$。——。

\textbf{}（best-response correspondence）$\Phi: [0,1] \to 2^{[0,1]}$：
\begin{equation}
\Phi(\mu) = \left\{ \tilde{\mu} \in [0,1] : \tilde{\mu} = \int_{\Theta} \mathbb{I}\left[B \geq -p \cdot \mathcal{V}(\mu) \cdot L(p, B)\right] dG(p, B) \right\}
\label{eq:br-correspondence}
\end{equation}
 $L(p, B) = L$ ，。$\Phi(\mu)$ " $\mu$  $\mathcal{V}(\mu)$，"。

\begin{theorem}[AAE —— Kakutani]
\label{thm:aae-fixed-point}
E1--E5， $\Phi: [0,1] \to 2^{[0,1]}$ Kakutani， $\mu^* \in \Phi(\mu^*)$。 $\mu^*$ -（AAE）。

：
\begin{enumerate}[label=(\roman*), itemsep=6pt]
\item \textbf{$\Phi$  $[0,1]$ ：}  $\mu \in [0,1]$，$\Phi(\mu) \subseteq [0,1]$（1）。

\item \textbf{$\Phi$ ：}  $\mu \in [0,1]$，$\Phi(\mu) \neq \emptyset$。\hfill\textit{(\ref{eq:br-correspondence}) $[0,1]$ 。}

\item \textbf{$\Phi$ ：}  $\mu \in [0,1]$，$\Phi(\mu)$ 。\hfill\textit{，(\ref{eq:br-correspondence})。 $G$ ，$\Phi(\mu)$ （）。}

\item \textbf{$\Phi$ （upper hemicontinuous）：}  $\mu_n \to \mu$，$\tilde{\mu}_n \in \Phi(\mu_n)$，$\tilde{\mu}_n \to \tilde{\mu}$， $\tilde{\mu} \in \Phi(\mu)$。\hfill\textit{ $\mathcal{V}(\mu)$ （E3）， $\mu$ （，）。}

\item \textbf{：} Kakutani(1941)， $\mu^*$  $\mu^* \in \Phi(\mu^*)$。\hfill$\square$
\end{enumerate}
\end{theorem}

\begin{proof}
Kakutani。 $[0,1]$  $\mathbb{R}$ 。 $\Phi$ ：

 $(\mu_n, \tilde{\mu}_n) \to (\mu, \tilde{\mu})$  $\tilde{\mu}_n \in \Phi(\mu_n)$  $n$ 。 $\tilde{\mu} \in \Phi(\mu)$。：
\begin{equation}
\tilde{\mu}_n = \int_{\Theta} \mathbb{I}\left[B \geq -p \cdot \mathcal{V}(\mu_n) \cdot L(p, B)\right] dG(p, B)
\end{equation}
 $f_n(p, B) = \mathbb{I}[B \geq -p \cdot \mathcal{V}(\mu_n) \cdot L(p, B)]$  $f(p, B) = \mathbb{I}[B \geq -p \cdot \mathcal{V}(\mu) \cdot L(p, B)]$（ $\mathcal{V}$ ）， $\{(p, B) : B = -p \cdot \mathcal{V}(\mu) \cdot L(p, B)\}$ 。 $G$-（E2 $G$ Lebesgue，，）。，$\tilde{\mu}_n \to \int_{\Theta} f(p, B) dG = \tilde{\mu}$， $\Phi$ 。，（Aliprantis \& Border, 2006, 17.11）。

： $[0,1]$， $\Phi(\mu)$ 。：，。 $G$ ，$\Phi(\mu)$ ，，。

，Kakutani，。
\end{proof}

\begin{corollary}[]
\label{cor:multiple-aae}
E1--E5，AAE。：
\begin{enumerate}[label=(\roman*)]
\item  $\Phi(0) = 0$（，）， $\mu^* = 0$（）AAE——。
\item  $\Phi(1) = 1$（，）， $\mu^* = 1$（）AAE——。
\item  $\mu^* \in (0,1)$。（）， $\Phi'(\mu^*) < 1$。
\item  $\Phi$ ， $\Phi$  $q(\mu)$  $G$ 。 $q(\mu)$ S（sigmoid）———— $\Phi$ 。
\end{enumerate}
\end{corollary}

\begin{proof}
\ref{thm:aae-fixed-point}，$\Phi$ 。。$\Phi(0) = 0$  $(p, B) \in \text{supp}(G)$，$B < -p \cdot \mathcal{V}(0) \cdot L$ ——，（ $\mathcal{V}(0) = 0$， $B < 0$ ，）。 $\mu^* = 1$。 $\Phi$ 。
\end{proof}

\begin{proposition}[AAE]
\label{prop:aae-welfare}
E1--E5，AAE：
\begin{enumerate}[label=(\roman*)]
\item \textbf{：}  $\mu^* > 0$ AAE，。。（） $\mu^*$。
\item \textbf{：}  $p_i$  $B_i$ （，——），AAE\textit{}： $p_i$（）， $p_i$（）。。
\item \textbf{：}  $\mu^* = 1$（），，——""。
\item \textbf{-：}  $q(\mu)$ （），（$B_i$ ），（$B_i$ ）。\textit{}—— $q$ 。
\end{enumerate}
\end{proposition}

\begin{proof}
\ref{thm:aae-fixed-point}(\ref{eq:adopt-condition})。，Spence(1973)：——（），。
\end{proof}

\begin{remark}[：]
\ref{thm:aae-fixed-point}，：
\begin{enumerate}[label=(\arabic*), leftmargin=*]
\item \textbf{（E1）：} （，$n = 2$$3$），(\ref{eq:br-correspondence})，$\Phi(\mu)$ 。Kakutani，(1951)——。
\item \textbf{ $q(\mu)$ （E3）：} ，$q(\mu)$ ——，。 $q(\mu)$ """+"，。
\item \textbf{（E2）：} （：），。 $\Phi(\mu)$ ，（Schauder）Kakutani。
\item \textbf{ $L$ ：} (\ref{eq:adopt-condition}) $p_i$ 。（——，），，$\Phi$ 。，Brouwer（ $\Phi$ ）（Eilenberg-Montgomery, 1946， $\Phi$ acyclic upper semicontinuous），。
\end{enumerate}
\end{remark}

\section{SCX：、}

、Yajie。：SCX，、。：(i) ————（）；(ii) ，————SCX（）；；(iii) SCX、（）。

\subsection{：，}

：，。——，。SCX，。

\begin{definition}[-]
 $\mathcal{D}$  $n$ 。$\mathcal{D}$  $\mathcal{D}_{\text{clean}}$  $\mathcal{D}_{\text{noise}}$：
\begin{equation}
\mathcal{D} = \mathcal{D}_{\text{clean}} \cup \mathcal{D}_{\text{noise}}, \quad \mathcal{D}_{\text{clean}} \cap \mathcal{D}_{\text{noise}} = \emptyset
\end{equation}
 $\mathcal{D}_{\text{noise}}$ 、、、。 $\rho(\mathcal{D}) = |\mathcal{D}_{\text{clean}}| / |\mathcal{D}|$。
\end{definition}

\begin{theorem}[]
\label{thm:storage-paradox}
 $V: 2^{\mathcal{D}} \to \mathbb{R}^+$ ，。 $\rho^* \in (0, 1)$， $\rho(\mathcal{D}) < \rho^*$ ：
\begin{equation}
V(\mathcal{D}_{\text{clean}}) > V(\mathcal{D})
\end{equation}
：。，SCX（Theorem 1），， $\mathcal{D}_{\text{clean}}$ 。
\end{theorem}

\textit{Proof sketch.}  $\mathcal{D}$  $f_{\mathcal{D}}$。 $\mathcal{D}_{\text{noise}}$ ：(a) ， $(1-\rho)$ （Natarajan et al., 2013）；(b) ，。Theorem 1，SCX。， $\epsilon > 0$， $k$ 1。， $k$， $\hat{\mathcal{D}}_{\text{clean}}$  $\mathbb{E}[V(\hat{\mathcal{D}}_{\text{clean}})] > V(\mathcal{D})$。\hfill $\square$

。Theorem~\ref{thm:storage-paradox} ：SCX，——，。5.6：CEC，CEC，。CEC——、，。，CEC。，（），：。，、CEC——。

\subsubsection{：Jevons}

Theorem~\ref{thm:storage-paradox} →→。，——，。，，：Jevons（Jevons, 1865）——，。

\medskip
\noindent\textbf{1. 。}

SCX——CEC $|\mathcal{E}_t|$  $\theta(t)$ ——：

\begin{definition}[]
 $\mathcal{S}(t)$  $t$ ，$\mathcal{G}(t)$ 。 $\alpha(t) = \theta(t) / \theta_{\max} \in [0, 1]$  $\mathcal{S}(t)$：

\textbf{：（）。} ， $\mathcal{D}_{\text{clean}}$ 。Theorem~\ref{thm:storage-paradox}，。 $\kappa(\alpha) \in [0, 1]$ ：
\begin{equation}
\kappa(\alpha) = \rho(\mathcal{D}) + (1 - \rho(\mathcal{D})) \cdot (1 - \alpha)^{\eta}
\end{equation}
 $\eta > 0$ 。$\kappa(\alpha) \to \rho(\mathcal{D})$  $\alpha \to 1$：，。

\textbf{：（）。} AI。(\ref{eq:quality})–(\ref{eq:model})， $Q_i^t$  $M_i^{t+1}$。，，：
\begin{equation}
\mathcal{G}(t) = \mathcal{G}_0 \cdot \left(1 + \beta \cdot \alpha(t)^{\gamma}\right)
\end{equation}
 $\beta > 0$ ，$\gamma > 0$ 。Jevons：（），，。
\end{definition}

\medskip
\noindent\textbf{2. 。}

。：
\begin{equation}
\frac{d\mathcal{S}}{dt} = \underbrace{\frac{d}{dt}\left[\kappa(\alpha(t)) \cdot \int_0^t \mathcal{G}(\tau) d\tau\right]}_{\text{}} + \underbrace{\kappa(\alpha(t)) \cdot \mathcal{G}(t)}_{\text{}}
\end{equation}

Leibniz：
\begin{equation}
\frac{d\mathcal{S}}{dt} = \underbrace{\kappa'(\alpha) \cdot \dot{\alpha}(t) \cdot \int_0^t \mathcal{G}(\tau) d\tau}_{\text{ }(<0)} + \underbrace{\kappa(\alpha(t)) \cdot \mathcal{G}(t)}_{\text{ }(>0)}
\end{equation}

（$\kappa'(\alpha) < 0$，），。。

\begin{definition}[]
。：
\begin{equation}
\Phi_{\text{quality}}(\alpha) = \frac{\partial (\text{})}{\partial \alpha} = \frac{\partial}{\partial \alpha}\left[\rho(\mathcal{D}) \cdot \int_0^t \mathcal{G}(\tau) d\tau\right]
\end{equation}
\begin{equation}
\Phi_{\text{quantity}}(\alpha) = \frac{\partial (\text{})}{\partial \alpha} = \frac{\partial}{\partial \alpha}\left[\int_0^t \mathcal{G}(\tau) d\tau\right]
\end{equation}

\textbf{} $\alpha^*$ ：
\begin{equation}
\Phi_{\text{quality}}(\alpha^*) = \Phi_{\text{quantity}}(\alpha^*)
\end{equation}

 $\Phi_{\text{quality}}(\alpha) > \Phi_{\text{quantity}}(\alpha)$ ————：$d\mathcal{S}/dt < 0$。\textit{}（denoising-dominated regime）。

 $\Phi_{\text{quality}}(\alpha) < \Phi_{\text{quantity}}(\alpha)$ ————：$d\mathcal{S}/dt > 0$。\textit{}（generation-dominated regime）。

：（$\alpha \approx 0$），（$\kappa'(\alpha)$ ，，），（）， $\Phi_{\text{quantity}} > \Phi_{\text{quality}}$，。（$\alpha \to 1$），（$\kappa(\alpha) \to \rho(\mathcal{D})$），（，）， $\Phi_{\text{quality}} > \Phi_{\text{quantity}}$，。， $\alpha^* \in (0, 1)$ 。
\end{definition}

\medskip
\noindent\textbf{3. ：。}

 $\mathcal{G}(t)$  $\alpha(t)$ 。，$\mathcal{G}(t)$ \textit{}。，。

 $N$ （、）。 $i$  $t$  $g_i(t)$， $\mathcal{G}(t) = \sum_i g_i(t)$。 $i$ ：
\begin{equation}
u_i(g_i, \mathbf{g}_{-i}; \alpha) = \underbrace{V\left(g_i \cdot \underbrace{\left[1 - (1-\alpha)^{\eta}\right]}_{\text{}}\right)}_{\text{}} - \underbrace{C(g_i)}_{\text{}} + \underbrace{\lambda \cdot \mathbb{I}\left[g_i > \bar{g}_{-i}\right]}_{\text{}}
\end{equation}
 $V(\cdot)$ （），$C(\cdot)$ （），$\lambda > 0$ \textit{}—— $i$ ，$i$ （、、）。$\bar{g}_{-i}$  $i$ 。

\begin{proposition}[]
\label{prop:preemptive}
（$\alpha \ll \alpha^*$），， $g^*(\alpha)$ ：
\begin{equation}
g^*(\alpha) = (C')^{-1}\left(V'(\cdot) \cdot [1 - (1-\alpha)^{\eta}] + \frac{\lambda}{N}\right)
\end{equation}
 $dg^*/d\alpha > 0$：，\textit{}，：(a) （$V'(\cdot)$ ），(b)  $\lambda/N$ 。

\textit{}： $\mathcal{G}(t)$  $\alpha^*$ 。""（front-running）：，——。 $\lambda$ ， $N$ （）。
\end{proposition}

\begin{proof}
 $i$  $\partial u_i / \partial g_i = 0$ 。，$g_i = g^*$  $i$ ， $\mathbb{I}[g_i > \bar{g}_{-i}] = 0$（）。 $V'(g^* \cdot [1 - (1-\alpha)^{\eta}]) \cdot [1 - (1-\alpha)^{\eta}] - C'(g^*) = 0$。 $g^*$ 。 $\lambda/N$  $\lambda$ ，。
\end{proof}

Proposition~\ref{prop:preemptive} ：，\textit{}——。；，。"，"，"，"。：（），。

，：

\begin{proposition}[]
\label{prop:mature}
（$\alpha > \alpha^*$），，\textit{}，：
\begin{equation}
g_i^{**}(\alpha) = (C')^{-1}\left(V'(\cdot) \cdot \rho(\mathcal{D})\right) \quad \text{} \quad \text{retain}_i = \kappa(\alpha) \cdot \mathcal{D}_i
\end{equation}

，。， $\mathcal{S}(t)$ ， $\sum_i |\mathcal{D}_{\text{clean}, i}^*|$。
\end{proposition}

\begin{proof}
 $\alpha > \alpha^*$ ，(a)  $\lambda$ ：，，；(b)  $\kappa(\alpha)$ ，（Theorem~\ref{thm:storage-paradox}，）""；(c) ：。。
\end{proof}

\medskip
\noindent\textbf{4. ：。}

Theorem~\ref{thm:storage-paradox}Proposition~\ref{prop:preemptive}–\ref{prop:mature}，：

\begin{enumerate}[label=\textbf{ \arabic*}:, leftmargin=*]
\item \textbf{（$\alpha \approx 0$）}：。，。：$\mathcal{S}(t) \approx \int_0^t \mathcal{G}(\tau) d\tau$。

\item \textbf{（$0 < \alpha < \alpha^*$）}：，。，（Proposition~\ref{prop:preemptive}）。\textit{}————（→→）（） $\mathcal{G}(t)$，。：。

\item \textbf{（$\alpha > \alpha^*$）}：。：（），（），（）。，。：$\mathcal{S}(t) \to \sum_i |\mathcal{D}_{\text{clean}, i}^*|$。
\end{enumerate}

\begin{remark}[]
：\textit{}。。（2），SCX""——。：。：2，；3，，。，$\alpha^*$  $\Phi_{\text{quality}}$  $\Phi_{\text{quantity}}$ ——，、。
\end{remark}

\begin{remark}[Spring]
\label{rem:golden-standard}
YajieSpring：—— $\Delta\mathcal{E}_t \approx \emptyset$——。，。Spring（Golden Standard）—— $\mathcal{D}_{\text{clean}}^*$，。：，。，，Spring。
\end{remark}

\subsection{}

SCX。（、、）（、、），。SCX——，Theorem 3（ (The Honest Person Theorem)）。

\begin{definition}[]
 $i$  $t$  $\mathcal{D}_i^t$  $M_i^t$。SCX：
\begin{align}
Q_i^t &= \theta(|\mathcal{E}_i^t|) \cdot \text{Audit}(\mathcal{D}_i^t) \label{eq:quality}\\
M_i^{t+1} &= f(M_i^t, Q_i^t, C_i) \label{eq:model}\\
U_i^{t+1} &= g(M_i^{t+1}) \label{eq:users}\\
|\mathcal{E}_i^{t+1}| &= |\mathcal{E}_i^t| + h(U_i^{t+1}) \label{eq:audit-growth}
\end{align}
 $Q_i^t$ Yajie，$\theta(\cdot)$ CEC（(3)），$C_i$ ，$U_i^{t+1}$ /，$h(\cdot)$ 。
\end{definition}

(\ref{eq:quality})–(\ref{eq:audit-growth})： →  →  →  →  → CEC → 。（cumulative causation），Myrdal (1957)  Arthur (1990) 。

\textbf{Theorem 3。} ——\textit{}。 $j$， $T$ SCXAI。 $t > T$， $j$  $M_j^t$  $i$  $M_i^t$。 $j$ ：$M_j^t < M_i^t$  (a) $j$  $\mathcal{D}_j^t$  $i$， (b) $j$  $f_j$  $i$  $f_i$， (c) ？

Theorem 3（ (The Honest Person Theorem)）：、，。Yajie——YajieCEC，CEC——。，：

\begin{equation}
\underbrace{M_j^t < M_i^t}_{\text{observable}} \quad \text{} \quad \begin{cases}
\mathcal{D}_j^t \text{ } &\text{()}\\
f_j \text{ } &\text{()}\\
\text{} &\text{()}
\end{cases}
\end{equation}

——CEC——，。：，，、。，。

\begin{remark}[]
SCX，：(1) 、（）；(2) （GPU、、）；(3) （CEC $|\mathcal{E}|^*$）；(4) ，（GDPR），。

。(2)（）(3)（27）。(1)（）(2)（）。。

，SCX————。（normative claim），（structural inference）：。，；、（），。——、、、——（audit receivers）：（），CEC。——，，——（dependency theory）"–"（Prebisch, 1950; Frank, 1967）：，。，：–，。

\textbf{——。} ，CEC\textit{}。5.2.2（Stage 6: ）\ref{cor:first-mover-decay}：。CEC，——CEC，""。——-，CEC——（$S_6$）（\ref{cor:S6-adoption-equivalence}）。，：\textbf{""""，}——。（CEC）（），（CEC）。""，。（）（\ref{cor:first-mover-decay}\ref{thm:universal-harmony-absorption}）。
\end{remark}

\subsection{：}

：。——、。\textit{}。

\begin{definition}[]
 $\mathcal{D}_{\text{phys}}$ （physically embedded data）——，、、。$\mathcal{D}_{\text{phys}}$  $\theta_{\text{audit}}$ ：
\begin{equation}
\theta_{\text{audit}} = \theta\left( \mathcal{A}, \mathcal{S}, \mathcal{H}, \mathcal{D}_{\text{phys}} \right)
\end{equation}
 $\mathcal{A}$ ，$\mathcal{S}$ （、、），$\mathcal{H}$ （、、）。$\mathcal{S}$  $\mathcal{H}$ ——，。，SCX\textit{}：$\partial \theta_{\text{audit}} / \partial (\text{distance}) < 0$。
\end{definition}

：，SCX（+++Spring），、。（），（），（）。：SCX，、、。

\begin{remark}[：]
——。（Wrigley, 1988）；（、、）、（Amsden, 1989; Wade, 1990）。，——，，，，。

SCX。SCX——、、、CECSpring——，：。SCX，，。：(a) （），(b) （CEC），(c) ， (d) （）。

SCX， $\Delta w$  $\Delta \theta_{\text{audit}}$。SCX，——。
\end{remark}

\textbf{。} （reshoring）——、（），。\textbf{，；，，、、。，。}SCX，，。（）Yajie、，。——。

：SCX——、、。SCX（，）：，\textit{}。：（CEC），（），（）。

\subsection{-}

（§11.2–11.3）SCX。：，\textit{-}（Manufacturing-Audit Complex, MAC）。。

\begin{definition}[-]
-：
\begin{enumerate}[label=(\roman*)]
    \item \textbf{}： $\mathcal{D}_{\text{phys}}$ ，SCX；
    \item \textbf{}：、Spring、，CEC；
    \item \textbf{-}：，GPU、。
\end{enumerate}
， $\theta_{\text{audit}}$ ，-。
\end{definition}

MAC\textit{}：，，。——。。

\subsubsection*{}

 $A$  $B$， $A$ （GDP，），$B$ （GDP，、、）。 $A$  $B$ SCX。，，：

\begin{equation}
\pi_i(s_i, s_{-i}) = R_i - C_i(s_i) - L_i(s_i, s_{-i})
\end{equation}

 $s_i \in \{0, 1\}$  $i$ SCX（$1 = $ ），$R_i$ ，$C_i(s_i)$ ，$L_i(s_i, s_{-i})$ 。：

\begin{enumerate}[label=(\roman*)]
    \item \textbf{}：$C_A(1) < C_B(1)$。 $A$ ——、Spring。 $B$ ，。
    \item \textbf{}：$\Delta_A > \Delta_B$， $\Delta_i = L_i(0, 1) - L_i(0, 0)$ ，$i$ 。，，；（、），。
\end{enumerate}

，：

\begin{theorem}[MAC]
\label{thm:mac-asymmetric}
 $C_A(1) < C_B(1)$  $\Delta_A > \Delta_B$。 $\Omega$  $(s_A^* = 1, s_B^* = 0)$——，。，：（$B$ ） $B$ ， $A$ 。
\end{theorem}

\textit{Proof sketch.} 。 $C_B(1) > \Delta_B$ ， $A$ ，$B$  $s_B = 0$。 $\Delta_A > C_A(1)$ ， $s_B = 0$，$A$  $s_A = 1$。 $\Omega$ （，，，）。 $C_B(1) > \Delta_B$ ：$B$ ， $(1, 1)$  $(1, 0)$  $B$ 。$(1, 0)$  $A$ （$A$ ）， $B$ ，。\hfill $\square$

\subsubsection*{Theorem 3：}

§11.2MAC。-，Theorem 3（ (The Honest Person Theorem)），。

 $B$，MAC（$s_B^* = 0$）。$B$ ——，——。 $B$ AI，§11.2： $A$ ， (a) $B$ （SCX）， (b) $B$ ， (c) $B$ \textit{}（、、） $A$ 。

MAC。——、。（、、）"ground truth"——""""。：

\begin{remark}[]
SCX\textit{}。 $\mathcal{D}_{\text{phys}}$ ——（SI），""。 $\mathcal{D}_{\text{service}}$ ""、——ground truth。：
\begin{equation}
\theta_{\text{audit}}(\mathcal{D}_{\text{phys}}) > \theta_{\text{audit}}(\mathcal{D}_{\text{service}})
\end{equation}
SCX $\mathcal{A}$ 。，。
\end{remark}

，：（），\textit{}（）。Theorem 3：""""，""（）（）。

\subsubsection*{}

——。""：

\begin{definition}[]
：
\begin{enumerate}[label=(\alph*)]
    \item \textbf{}（de jure standard）：——ISO/IEC 5259（）、ISO/IEC 42001（AI）——、；
    \item \textbf{}（de facto standard）：——CEC、、Spring——""。
\end{enumerate}
、——ISO/IEC 5259。：——CEC、、——（，）。
\end{definition}

：SCX，。，——（、）（Abbate, 1999; DeNardis, 2009）。，CEC。

MAC，：（），。，——。

\subsubsection*{--}

。SCXSpring， $k$  $n$ 。MAC，，————。

\begin{definition}[-]
 $i$ - $\eta_i$ ：
\begin{equation}
\eta_i = \frac{\partial \theta_{\text{audit}, i} / \theta_{\text{audit}, i}}{\partial E_i / E_i}
\end{equation}
 $E_i$  $i$ SCX（、）。$\eta_i$ 。$\eta_i$ ，。
\end{definition}

——（、、GPU）——。：

\begin{enumerate}[label=(\roman*)]
    \item \textbf{}：、。、、，SCX，CEC。
    \item \textbf{}：SCXGPUAI。（、）。——（）（）——。
\end{enumerate}

，--。、（），——，。（），——、IT——、、。

\begin{remark}[]
：SCX——，-——。、，SCXCEC，。MAC（Theorem~\ref{thm:mac-asymmetric}）， $C_A(1)$  $C_B(1)$ ——，，。
\end{remark}

\subsubsection*{：}

——MAC、、——§11.3：$\partial \theta_{\text{audit}} / \partial (\text{distance}) < 0$。：，。，。：

\begin{enumerate}[label=(\roman*)]
    \item \textbf{}：（-、），（Degen et al., 2017）。SCX，。
    \item \textbf{}：————，。。
\end{enumerate}

。，：\textbf{，。，。}，——，，。

\subsubsection*{}

，SCX（structural inferences）——，：

\begin{enumerate}[label=(\roman*)]
    \item ————SCX。：（§11.3）、（）、--（）。
    \item ——，。，CEC。
    \item 。，——。
    \item 。、、CEC，。。
    \item \textbf{——（）。} (i)--(iv)————（5.2.2）\textit{}，。CEC（\ref{cor:first-mover-decay}）：(a) CEC（\ref{cor:self-reinforcing}）；(b) CEC，，CEC；(c) CEC $\Delta(|\mathcal{E}|)$ （\ref{cor:S6-adoption-equivalence}），；(d) """"（\ref{def:S6}），。 $S_5$（） $S_6$（）——， $\pi_{\text{political}}$ （M8）。：\textbf{（）MAC，}——CEC，（、CEC）。，。
\end{enumerate}

，————。。、、，。

\textbf{SCX。} ：SCX。。Theorem~1（）Yajie。，（7nm5nm3nm）（、、、），——、、——、。、、，；——（、、），，。SCX，，。：。、，。SCX，""————，。

\textbf{：。} SCX，。：§11.3，SCX——、、，。——。、，""（），——。，，——、、——。SCX，。

\textbf{——。} SCX，\textit{}（，§11.4），\textit{}。\ref{cor:first-mover-decay}\ref{thm:universal-harmony-absorption}：CEC，——CEC（ $\Delta(\cdot)$ ，CEC $\Delta^{\text{public}}$  $\lambda - \kappa$），CEC，。（$S_6$），"""CERN"——，。，。SCX""，。

\textbf{。} 。，：（、、）、（）、（）。SCX，，：（i）——、；（ii）——、；（iii）——CEC、。：（，），——CEC，。SCX，。，：CEC，、，。

\textbf{：。} ，：——（§11.3）——。，。，——（Degen et al., 2017），——。，：，，SCX。，，————。，，、。。：，。，。


\section*{References}

\bibitem{adalovelace2020} Ada Lovelace Institute. \textit{Examining the Black Box: Tools for Assessing Algorithmic Systems}. London (2020).

\bibitem{armstrong2006} Armstrong, M. Competition in two-sided markets. \textit{The RAND Journal of Economics} \textbf{37}(3), 668--691 (2006).

\bibitem{arthur1989} Arthur, W. B. Competing technologies, increasing returns, and lock-in by historical events. \textit{The Economic Journal} \textbf{99}(394), 116--131 (1989).

\bibitem{bateman2022} Bateman, J. \textit{U.S.-China Technological ``Decoupling'': A Strategy and Policy Framework}. Carnegie Endowment (2022).

\bibitem{batini2009} Batini, C., Cappiello, C., Francalanci, C. \& Maurino, A. Methodologies for data quality assessment and improvement. \textit{ACM Computing Surveys} \textbf{41}(3), 1--52 (2009).

\bibitem{david1985} David, P. A. Clio and the economics of QWERTY. \textit{The American Economic Review} \textbf{75}(2), 332--337 (1985).

\bibitem{diakopoulos2015} Diakopoulos, N. Algorithmic accountability. \textit{Digital Journalism} \textbf{3}(3), 398--415 (2015).

\bibitem{drezner2015} Drezner, D. W. Targeted sanctions in a world of global finance. \textit{International Interactions} \textbf{41}(4), 755--764 (2015).

\bibitem{ec2021} European Commission. \textit{Proposal for a Regulation Laying Down Harmonised Rules on Artificial Intelligence}. COM(2021) 206 final (2021).

\bibitem{farrell2019} Farrell, H. \& Newman, A. L. Weaponized interdependence. \textit{International Security} \textbf{44}(1), 42--79 (2019).

\bibitem{farrell2007} Farrell, J. \& Klemperer, P. Coordination and lock-in. In \textit{Handbook of Industrial Organization}, Vol.~3, pp.~1967--2072. Elsevier (2007).

\bibitem{farrell1985} Farrell, J. \& Saloner, G. Standardization, compatibility, and innovation. \textit{The RAND Journal of Economics} \textbf{16}(1), 70--83 (1985).

\bibitem{floridi2018} Floridi, L., Cowls, J., Beltrametti, M., et al. AI4People---An ethical framework for a good AI society. \textit{Minds and Machines} \textbf{28}(4), 689--707 (2018).

\bibitem{gebru2021} Gebru, T., Morgenstern, J., Vecchione, B., Vaughan, J. W., Wallach, H., Daum\'e III, H. \& Crawford, K. Datasheets for datasets. \textit{Communications of the ACM} \textbf{64}(12), 86--92 (2021).

\bibitem{gregory2021} Gregory, R. W., Henfridsson, O., Kaganer, E. \& Kyriakou, H. The role of artificial intelligence and data network effects. \textit{Academy of Management Review} \textbf{46}(3), 534--551 (2021).

\bibitem{jobin2019} Jobin, A., Ienca, M. \& Vayena, E. The global landscape of AI ethics guidelines. \textit{Nature Machine Intelligence} \textbf{1}(9), 389--399 (2019).

\bibitem{joyner2011} Joyner, D. H. \textit{Interpreting the Nuclear Non-Proliferation Treaty}. Oxford University Press (2011).

\bibitem{kaplan2020} Kaplan, J., McCandlish, S., Henighan, T., Brown, T. B., Chess, B., Child, R., Gray, S., Radford, A., Wu, J. \& Amodei, D. Scaling laws for neural language models. \textit{arXiv preprint arXiv:2001.08361} (2020).

\bibitem{katz1985} Katz, M. L. \& Shapiro, C. Network externalities, competition, and compatibility. \textit{The American Economic Review} \textbf{75}(3), 424--440 (1985).

\bibitem{khan2017} Khan, L. M. Amazon's antitrust paradox. \textit{Yale Law Journal} \textbf{126}(3), 710--805 (2017).

\bibitem{mearsheimer1993} Mearsheimer, J. J. The case for a Ukrainian nuclear deterrent. \textit{Foreign Affairs} \textbf{72}(3), 50--66 (1993).

\bibitem{mokander2021} M\"{o}kander, J., Morley, J., Taddeo, M. \& Floridi, L. Ethics-based auditing of automated decision-making systems. \textit{Science and Engineering Ethics} \textbf{27}(4), 1--30 (2021).

\bibitem{nist2023} NIST. \textit{Artificial Intelligence Risk Management Framework 1.0}. U.S. Department of Commerce (2023).

\bibitem{northcutt2021} Northcutt, C. G., Athalye, A. \& Mueller, J. Pervasive label errors in test sets destabilize machine learning benchmarks. \textit{Proc. NeurIPS} (2021).

\bibitem{parasuraman1997} Parasuraman, R. \& Riley, V. Humans and automation: Use, misuse, disuse, abuse. \textit{Human Factors} \textbf{39}(2), 230--253 (1997).

\bibitem{parker2005} Parker, G. G. \& Van Alstyne, M. W. Two-sided network effects. \textit{Management Science} \textbf{51}(10), 1494--1504 (2005).

\bibitem{pushkarna2022} Pushkarna, M., Zaldivar, A. \& Kjartansson, O. Data cards: Purposeful and transparent dataset documentation. \textit{Proc. FAccT} 1776--1826 (2022).

\bibitem{raji2020} Raji, I. D., Smart, A., White, R. N., Mitchell, M., Gebru, T., Hutchinson, B., Smith-Loud, J., Theron, D. \& Barnes, P. Closing the AI accountability gap. \textit{Proc. FAccT} 33--44 (2020).

\bibitem{rochet2003} Rochet, J.-C. \& Tirole, J. Platform competition in two-sided markets. \textit{Journal of the European Economic Association} \textbf{1}(4), 990--1029 (2003).

\bibitem{sagan1996} Sagan, S. D. Why do states build nuclear weapons? \textit{International Security} \textbf{21}(3), 54--86 (1996).

\bibitem{sambasivan2021} Sambasivan, N., Kapania, S., Highfill, H., Akrong, D., Paritosh, P. \& Aroyo, L. M. ``Everyone wants to do the model work, not the data work'': Data cascades in high-stakes AI. \textit{Proc. CHI} 1--15 (2021).

\bibitem{sandvig2014} Sandvig, C., Hamilton, K., Karahalios, K. \& Langbort, C. Auditing algorithms: Research methods for detecting discrimination on internet platforms. \textit{Data and Discrimination: Converting Critical Concerns into Productive Inquiry} 1--23 (2014).

\bibitem{segal2020} Segal, A. The coming tech Cold War with China. \textit{Foreign Affairs} \textbf{99}(5), 94--105 (2020).

\bibitem{waltz1981} Waltz, K. N. The spread of nuclear weapons: More may be better. \textit{The Adelphi Papers} \textbf{21}(171), 1--32 (1981).

\bibitem{wang1996} Wang, R. Y. \& Strong, D. M. Beyond accuracy: What data quality means to data consumers. \textit{Journal of Management Information Systems} \textbf{12}(4), 5--33 (1996).

\bibitem{wilson2021} Wilson, C., Ghosh, A., Jiang, S., Mislove, A., Baker, L., Szurley, J., Dud\'ik, M., Danks, D. \& Ramachandran, P. Building and auditing fair algorithms. \textit{Proc. FAccT} 666--677 (2021).

\bibitem{zeff2012} Zeff, S. A. The evolution of the IASC into the IASB. \textit{The Accounting Review} \textbf{87}(3), 807--837 (2012).

\bibitem{zhang2014} Zhang, Y. \& Andrew, J. Financialisation and the conceptual framework. \textit{Critical Perspectives on Accounting} \textbf{25}(1), 17--26 (2014).

\end{thebibliography}

\section*{Author's Statement}

The author publishes under the name SCX. The work is released on GitHub for immediate open access. No institutional affiliation is claimed. The theorems stand on their own — they do not depend on the identity or location of their author. Verification is welcome; refutation is welcome; both must follow the standards of the framework: stated assumptions, verifiable evidence, and rigorous reasoning.

\section*{Personal Note}

Game theory is truly fascinating and deeply enjoyable.

\end{document}
